% La transmission des caractères au sein d'une espèce et quelques aspects de la génétique humaine — SVTEEHB Tle C — Cours signé OnBuch+
% Programme officiel MINESEC. Compiler avec : tectonic N04-transmission-caracteres-genetique-humaine.tex
\newcommand{\DOCMATIERE}{SVTEEHB}
\newcommand{\DOCNIVEAU}{Tle C}
\newcommand{\DOCMODULE}{Module I : Le Monde Vivant}
\newcommand{\DOCLECON}{N04}
\newcommand{\DOCTITRE}{La transmission des caractères au sein d'une espèce et quelques aspects de la génétique humaine}
% OnBuch+ — préambule « cours signés » ENRICHI (compilé avec tectonic / XeLaTeX)
% Système visuel « Pop » d'OnBuch+ : fond crème (jamais blanc), bordures encre
% épaisses, ombres « dures » décalées, titres Archivo Black en capitales.
% Version MATHÉMATIQUES : graphes (pgfplots/tikz), boîtes pédagogiques
% (définition, propriété, méthode, attention, le savais-tu, exercice résolu,
% exercice, TP, bilan), page de garde et signature OnBuch+ ancrée en bas.
%
% Macros attendues dans main.tex AVANT \input{preamble} :
%   \DOCMATIERE  \DOCNIVEAU  \DOCMODULE  \DOCLECON (numéro)  \DOCTITRE
\documentclass[11pt]{article}

\usepackage{fontspec}
\usepackage[french]{babel}
\usepackage[a4paper,margin=2cm,top=3.4cm,bottom=2.8cm,headheight=1.4cm,footskip=1.3cm]{geometry}
\usepackage{amsmath,amssymb,amsfonts}
\usepackage{mathtools} % \xmapsto, \coloneqq... souvent utilisés par les modèles
\usepackage{cancel} % simplifications barrées (bilans de forces)
% \abs{z} (module, valeur absolue) et \norm{u} (norme), utilisés par les modèles
\providecommand{\abs}{}\let\abs\relax\DeclarePairedDelimiter{\abs}{\lvert}{\rvert}
\providecommand{\norm}{}\let\norm\relax\DeclarePairedDelimiter{\norm}{\lVert}{\rVert}
\usepackage[table]{xcolor} % [table] : fournit aussi \rowcolors (alternance de lignes)
\usepackage{pagecolor}
\usepackage{tikz}
\usetikzlibrary{arrows.meta,positioning,calc,shapes.geometric,shapes.misc,shapes.arrows,decorations.pathmorphing,decorations.pathreplacing,decorations.markings,patterns,shadows,mindmap,trees,fit,backgrounds,matrix,intersections,angles,quotes}
\usepackage{pgfplots}
\usepackage[version=4]{mhchem} % équations nucléaires \ce{^{235}_{92}U}
\usepackage[european,siunitx]{circuitikz} % schémas électriques
\pgfplotsset{compat=1.18}
\usepgfplotslibrary{fillbetween,groupplots} % aires entre courbes (calcul intégral)
\usepackage{enumitem}
\usepackage{fancyhdr}
% [most] charge toutes les bibliothèques tcolorbox (listings, minted, raster,
% documentation...) et gonfle inutilement la mémoire TeX sur un document long
% (~300 boîtes) : on ne charge que ce qui est réellement utilisé.
\usepackage[skins,breakable]{tcolorbox}
\usepackage{graphicx}
\usepackage{colortbl}
\usepackage{booktabs}
\usepackage{tabularx}
\usepackage{multirow}
\usepackage{diagbox} % case d'en-tête diagonale des tableaux à double entrée
% Colonnes extensibles centrée / gauche / droite (C, L, R), souvent utilisées par les modèles
\newcolumntype{C}{>{\centering\arraybackslash}X}
\newcolumntype{L}{>{\raggedright\arraybackslash}X}
\newcolumntype{R}{>{\raggedleft\arraybackslash}X}
\usepackage{array}
\usepackage{multicol}
\usepackage{siunitx}
% parse-numbers=false : accepte aussi \SI{2,5 \times 10^{-3}}{...}
\sisetup{locale=FR,output-decimal-marker={,},group-minimum-digits=4,parse-numbers=false}
\DeclareSIUnit\ohm{\ensuremath{\symup{\Omega}}} % U+2126 absent de la police
\DeclareSIUnit\division{div} % réglages d'oscilloscope (ms/div, V/div)
\DeclareSIUnit\tr{tr} % tours (vitesse de rotation, ex. \SI{1500}{\tr\per\minute})
\DeclareSIUnit\molar{M} % concentration molaire (ex. \SI{3}{\molar}, \SI{1}{\micro\molar})
\DeclareSIUnit\an{an} % année (le modèle confond parfois avec \year, un compteur TeX) — ex. \SI{5}{\kilo\meter\per\an}
\DeclareSIUnit\FCFA{FCFA} % franc CFA (ex. \SI{500}{\FCFA\per\kilo\gram})
\DeclareSIUnit\degreeBrix{\textdegree Brix} % teneur en sucre d'un sirop/jus (réfractométrie)
\DeclareSIUnit\month{mois} % le modèle utilise parfois \month, un compteur TeX, comme unité de durée
\DeclareSIUnit\microliter{\micro\liter} % le modèle écrit parfois \microliter en un seul token
\DeclareSIUnit\million{millions} % dénombrement (ex. \SI{15}{\million\per\mL} spermatozoïdes)
\usepackage{qrcode}
% \degree : commande fréquemment utilisée par les modèles pour un angle ou une
% température en dehors de \SI{}{} (non fournie par les paquets chargés).
\providecommand{\degree}{^{\circ}}
% \qed (fin de démonstration), utilisé par les modèles sans amsthm
\providecommand{\qed}{\hfill$\square$}
% \barre : double barre des valeurs interdites dans un tableau de variations (tkz-tab)
\providecommand{\barre}{\ensuremath{\|}}
% environnement proof (amsthm non chargé) utilisé par les modèles
\ifdefined\proof\else\newenvironment{proof}[1][Démonstration]{\par\noindent\textit{#1.}\ }{\qed\par}\fi
\usepackage{lastpage}
\usepackage{hyperref}
\hypersetup{hidelinks}

% ============================================================
% Palette « Pop » OnBuch+ — mêmes hex que la classe P (plus_ui.dart)
% ============================================================
\definecolor{poporange}{HTML}{F59321}
\definecolor{popdark}{HTML}{E07A0C}
\definecolor{popcream}{HTML}{FFF6E8}
\definecolor{popink}{HTML}{1C1714}
\definecolor{poppurple}{HTML}{7A5AE0}
\definecolor{popgreen}{HTML}{2E9E6A}
\definecolor{poppink}{HTML}{E0457B}
\definecolor{popblue}{HTML}{2F7DE1}
\definecolor{popgold}{HTML}{C98A12}
\definecolor{popmuted}{HTML}{8A7F76}
\definecolor{popline}{HTML}{EADFD2}
\definecolor{popgray}{HTML}{8A7F76}
\colorlet{popgrey}{popgray}
% Alias tolérants (le modèle nomme parfois les couleurs différemment)
\colorlet{popwhite}{white}
\colorlet{popblack}{popink}
\colorlet{popred}{poppink}
\colorlet{popyellow}{popgold}
% Teintes claires (fonds de tableaux/encadrés) — le modèle nomme parfois
% les couleurs avec un suffixe L (ex. popblueL) sans les avoir définies.
\colorlet{poporangeL}{poporange!20}
\colorlet{popdarkL}{popdark!20}
\colorlet{poppurpleL}{poppurple!20}
\colorlet{popgreenL}{popgreen!20}
\colorlet{poppinkL}{poppink!20}
\colorlet{popblueL}{popblue!20}
\colorlet{popgoldL}{popgold!20}
\colorlet{popredL}{popred!20}
\colorlet{popgrayL}{popgray!20}
% Teintes claires pour fonds de tableaux / zones
\colorlet{popblueL}{popblue!10!white}
\colorlet{poporangeL}{poporange!14!white}
\colorlet{popgreenL}{popgreen!12!white}
\colorlet{poppurpleL}{poppurple!12!white}
\colorlet{poppinkL}{poppink!12!white}
\colorlet{popgoldL}{popgold!14!white}

\pagecolor{popcream}

% ============================================================
% Typographie
% ============================================================
\defaultfontfeatures{Ligatures=TeX}
\setmainfont{PlusJakartaSans-Medium.ttf}[Path=../../../fonts/,BoldFont=PlusJakartaSans-Bold.ttf]
% Maths en sans-serif (Fira Math) avec chiffres et lettres droites de Plus
% Jakarta, pour que les formules se fondent dans le texte.
\usepackage[math-style=ISO,bold-style=ISO]{unicode-math}
\setmathfont{FiraMath-Regular.otf}[Path=../../../fonts/]
\setmathfont{PlusJakartaSans-Medium.ttf}[Path=../../../fonts/,range={up/num,up/{Latin,latin},\mathup}]
\usepackage{icomma} % virgule décimale sans espace en mode math
% Caractères mathématiques Unicode absents des polices de texte (Plus Jakarta,
% Archivo Black) : rendus par leur équivalent mathématique, en texte comme en
% formule — sinon ils s'affichent en carré vide (ex. « Arithmétique dans ℤ »).
\usepackage{newunicodechar}
\newunicodechar{ℝ}{\ensuremath{\mathbb{R}}}
\newunicodechar{ℤ}{\ensuremath{\mathbb{Z}}}
\newunicodechar{ℂ}{\ensuremath{\mathbb{C}}}
\newunicodechar{ℕ}{\ensuremath{\mathbb{N}}}
\newunicodechar{ℚ}{\ensuremath{\mathbb{Q}}}
\newunicodechar{𝔻}{\ensuremath{\mathbb{D}}}
\newunicodechar{α}{\ensuremath{\alpha}}
\newunicodechar{β}{\ensuremath{\beta}}
\newunicodechar{γ}{\ensuremath{\gamma}}
\newunicodechar{δ}{\ensuremath{\delta}}
\newunicodechar{ε}{\ensuremath{\varepsilon}}
\newunicodechar{θ}{\ensuremath{\theta}}
\newunicodechar{λ}{\ensuremath{\lambda}}
\newunicodechar{μ}{\ensuremath{\mu}}
\newunicodechar{σ}{\ensuremath{\sigma}}
\newunicodechar{τ}{\ensuremath{\tau}}
\newunicodechar{φ}{\ensuremath{\varphi}}
\newunicodechar{ψ}{\ensuremath{\psi}}
\newunicodechar{ω}{\ensuremath{\omega}}
\newunicodechar{Σ}{\ensuremath{\Sigma}}
% majuscules produites par \MakeUppercase dans les titres (α-aminés -> Α)
\newunicodechar{Α}{\ensuremath{\alpha}}
\newunicodechar{Β}{\ensuremath{\beta}}
\newunicodechar{Γ}{\ensuremath{\Gamma}}
\newunicodechar{Δ}{\ensuremath{\Delta}}
\newunicodechar{Ε}{\ensuremath{\varepsilon}}
\newunicodechar{Θ}{\ensuremath{\Theta}}
\newunicodechar{Λ}{\ensuremath{\Lambda}}
\newunicodechar{Μ}{\ensuremath{\mu}}
\newunicodechar{Τ}{\ensuremath{\tau}}
\newunicodechar{Φ}{\ensuremath{\Phi}}
\newunicodechar{Ψ}{\ensuremath{\Psi}}
\newunicodechar{Ω}{\ensuremath{\Omega}}
\newunicodechar{↦}{\ensuremath{\mapsto}}
\newunicodechar{∈}{\ensuremath{\in}}
\newunicodechar{∉}{\ensuremath{\notin}}
\newunicodechar{∧}{\ensuremath{\wedge}}
\newunicodechar{⇔}{\ensuremath{\Leftrightarrow}}
\newunicodechar{⟺}{\ensuremath{\Longleftrightarrow}}
\newunicodechar{⇒}{\ensuremath{\Rightarrow}}
\newunicodechar{⟹}{\ensuremath{\Longrightarrow}}
\newunicodechar{∘}{\ensuremath{\circ}}
\newunicodechar{∪}{\ensuremath{\cup}}
\newunicodechar{∩}{\ensuremath{\cap}}
\newunicodechar{≡}{\ensuremath{\equiv}}
\newunicodechar{⊂}{\ensuremath{\subset}}
\newunicodechar{⊥}{\ensuremath{\perp}}
\newunicodechar{∃}{\ensuremath{\exists}}
\newunicodechar{∀}{\ensuremath{\forall}}
\newunicodechar{∛}{\ensuremath{\sqrt[3]{\,}}}
\newunicodechar{∜}{\ensuremath{\sqrt[4]{\,}}}
\newunicodechar{⃗}{\ensuremath{{}^{\to}}}
\newunicodechar{ⁿ}{\ifmmode^{n}\else\textsuperscript{\ensuremath{n}}\fi}
\newunicodechar{ˣ}{\ifmmode^{x}\else\textsuperscript{\ensuremath{x}}\fi}
\newunicodechar{ᵇ}{\ifmmode^{b}\else\textsuperscript{\ensuremath{b}}\fi}
\newunicodechar{ʳ}{\ifmmode^{r}\else\textsuperscript{\ensuremath{r}}\fi}
\newunicodechar{ᵘ}{\ifmmode^{u}\else\textsuperscript{\ensuremath{u}}\fi}
\newunicodechar{ʸ}{\ifmmode^{y}\else\textsuperscript{\ensuremath{y}}\fi}
\newunicodechar{ᵃ}{\ifmmode^{a}\else\textsuperscript{\ensuremath{a}}\fi}
\newunicodechar{ᵉ}{\ifmmode^{e}\else\textsuperscript{\ensuremath{e}}\fi}
\newunicodechar{ᵏ}{\ifmmode^{k}\else\textsuperscript{\ensuremath{k}}\fi}
\newunicodechar{ᵖ}{\ifmmode^{p}\else\textsuperscript{\ensuremath{p}}\fi}
\newunicodechar{ᴺ}{\ifmmode^{N}\else\textsuperscript{\ensuremath{N}}\fi}
\newunicodechar{ᶜ}{\ifmmode^{c}\else\textsuperscript{\ensuremath{c}}\fi}
\newunicodechar{ʰ}{\ifmmode^{h}\else\textsuperscript{\ensuremath{h}}\fi}
\newunicodechar{ᵠ}{\ifmmode^{q}\else\textsuperscript{\ensuremath{q}}\fi}
\newunicodechar{ₙ}{\ifmmode_{n}\else\textsubscript{\ensuremath{n}}\fi}
\newunicodechar{ₐ}{\ifmmode_{a}\else\textsubscript{\ensuremath{a}}\fi}
\newunicodechar{ᵢ}{\ifmmode_{i}\else\textsubscript{\ensuremath{i}}\fi}
\newunicodechar{ᵣ}{\ifmmode_{r}\else\textsubscript{\ensuremath{r}}\fi}
\newunicodechar{ₖ}{\ifmmode_{k}\else\textsubscript{\ensuremath{k}}\fi}
\newunicodechar{ᵦ}{\ifmmode_{\beta}\else\textsubscript{\ensuremath{\beta}}\fi}
\newunicodechar{ₓ}{\ifmmode_{x}\else\textsubscript{\ensuremath{x}}\fi}
\newfontfamily\popdisplay{ArchivoBlack-Regular.ttf}[Path=../../../fonts/]
\newfontfamily\poplogo{SpaceGrotesk-Bold.ttf}[Path=../../../fonts/]
\newfontfamily\popsemi{PlusJakartaSans-SemiBold.ttf}[Path=../../../fonts/]
\newfontfamily\popxbold{PlusJakartaSans-ExtraBold.ttf}[Path=../../../fonts/]
\newfontfamily\popxboldls{PlusJakartaSans-ExtraBold.ttf}[Path=../../../fonts/,LetterSpace=3.0]

\setlist[enumerate]{leftmargin=1.7em,itemsep=5pt,topsep=4pt}
\setlist[itemize]{leftmargin=1.7em,itemsep=4pt,topsep=4pt,label={\color{poporange}\textbullet}}
\setlength{\parindent}{0pt}
\setlength{\parskip}{5pt}
\renewcommand{\arraystretch}{1.25}

% pgfplots : style Pop par défaut
\pgfplotsset{
  every axis/.append style={axis line style={popink,line width=1pt},tick style={popink},
    grid style={popline,line width=0.5pt},label style={font=\small},tick label style={font=\footnotesize}},
  popcurve/.style={poporange,line width=1.8pt,no markers,smooth},
}
% Même style disponible dans un \draw TikZ simple (hors pgfplots)
\tikzset{popcurve/.style={poporange,line width=1.8pt}}
\tikzset{popgrid/.style={popline,very thin}}
\colorlet{popcurve}{poporange} % le nom du style est parfois utilisé comme couleur

% ============================================================
% Pill / squiggle
% ============================================================
\newcommand{\poppill}[3]{%
  \tikz[baseline=-0.55ex]\node[draw=popink,line width=1.2pt,rounded corners=2.6pt,%
    fill=#2,inner xsep=6.5pt,inner ysep=3pt]{%
    \popxboldls\fontsize{7}{7}\selectfont\color{#3}\MakeUppercase{#1}};%
}
\newcommand{\popsquiggle}[1]{%
  \tikz[baseline=-0.4ex]{
    \draw[#1,line width=2.6pt,line cap=round]
      plot [smooth] coordinates {(0,0.09) (0.34,0.01) (0.68,0.21) (1.02,0.01) (1.36,0.10)};
  }%
}
% Mot-clé surligné (marqueur orange sous le texte)
\newcommand{\cle}[1]{{\popxbold\color{popdark}#1}}

% ============================================================
% En-tête / pied
% ============================================================
\pagestyle{fancy}
\fancyhf{}
\renewcommand{\headrulewidth}{0pt}
\renewcommand{\footrulewidth}{0pt}
\lhead{\raisebox{-0.15ex}{\poplogo\fontsize{13}{13}\selectfont\color{popink}OnBuch\color{poporange}+}}
\rhead{\poppill{\DOCMATIERE\ \textbullet\ \DOCNIVEAU\ \textbullet\ Leçon \DOCLECON}{white}{popink}}
\lfoot{\popxbold\fontsize{7.6}{7.6}\selectfont\color{popmuted}\MakeUppercase{Cours signé OnBuch+}\ \textbullet\ {\popsemi\DOCTITRE}}
\rfoot{\tikz[baseline=-0.6ex]\node[rounded corners=8.5pt,fill=popink,minimum height=17pt,minimum width=17pt,inner xsep=5pt,inner sep=0]{\popxbold\fontsize{8}{8}\selectfont\color{popcream}\thepage\,/\,\pageref*{LastPage}};}

% ============================================================
% Titres
% ============================================================
\newcommand{\chaptitre}[2]{%
  \begin{tcolorbox}[enhanced,colback=poporange,colframe=popink,boxrule=2.6pt,arc=11pt,
      drop shadow=popink,left=15pt,right=15pt,top=12pt,bottom=11pt,before skip=3pt,after skip=18pt]
    \poppill{#2}{popink}{popcream}\par\vspace{9pt}
    {\raggedright\hyphenpenalty=10000\exhyphenpenalty=10000\popdisplay\fontsize{22}{24}\selectfont\color{popcream}\MakeUppercase{#1}\par}
  \end{tcolorbox}
}
% Section numérotée : pastille orange avec le numéro + titre Archivo
\newcounter{popsec}
\newcommand{\coursec}[1]{%
  \stepcounter{popsec}\setcounter{popsub}{0}%
  \par\vspace{16pt}\noindent
  \begin{minipage}{\linewidth}
  \tikz[baseline=-0.7ex]\node[draw=popink,line width=1.6pt,fill=poporange,rounded corners=4pt,minimum size=22pt,inner sep=0]{\popdisplay\fontsize{12}{12}\selectfont\color{popcream}\thepopsec};%
  \hspace{8pt}{\popdisplay\fontsize{15}{17}\selectfont\color{popink}\MakeUppercase{#1}}\par
  \vspace{4pt}\noindent{\color{poporange}\rule{36pt}{3.6pt}}
  \end{minipage}\par\nopagebreak\vspace{8pt}
}
\newcounter{popsub}
\newcommand{\courssub}[1]{%
  \stepcounter{popsub}%
  \par\vspace{9pt}\noindent{\popxbold\fontsize{11.5}{13}\selectfont\color{popink}{\color{poporange}\thepopsec.\thepopsub}\hspace{6pt}#1}\par\nopagebreak\vspace{4pt}
}

% ============================================================
% Boîtes pédagogiques — même recette : fond blanc, bordure épaisse colorée,
% ombre dure encre, badge plein. Titre optionnel [#1].
% ============================================================
\tcbset{popbase/.style={breakable,enhanced,colback=white,boxrule=2.3pt,arc=9pt,
  shadow={3pt}{-3pt}{0pt}{popink},left=14pt,right=14pt,top=11pt,bottom=11pt,before skip=11pt,after skip=15pt,
  fontupper=\popsemi\fontsize{10.6}{14.5}\selectfont\color{popink},
  fontlower=\popsemi\fontsize{10.6}{14.5}\selectfont\color{popink}}}
% Coche dessinée (le glyphe ✓ n'existe pas dans les polices du document)
\newcommand{\popcheck}[1]{\tikz[baseline=-0.5ex]\draw[#1,line width=1.6pt,line cap=round,line join=round](-0.13,0.0)--(-0.04,-0.09)--(0.13,0.1);}
\providecommand{\coche}{\popcheck{popgreen}}
\providecommand{\resultat}[1]{\par\smallskip\noindent\fcolorbox{popgreen}{popgreenL}{\parbox{\dimexpr\linewidth-2\fboxsep-2\fboxrule\relax}{\textbf{Résultat.} #1}}\par\smallskip}
\newcommand{\popboxhead}[3]{\poppill{#1}{#2}{white}\ifx\relax#3\relax\else\hspace{6pt}{\popxbold\fontsize{10.6}{12}\selectfont\color{popink}#3}\fi\par\vspace{7pt}}

\newtcolorbox{aretenir}[1][]{popbase,colframe=popgreen,before upper={\popboxhead{À retenir}{popgreen}{#1}}}
\newtcolorbox{exemplebox}[1][]{popbase,colframe=poporange,before upper={\popboxhead{Exemple}{poporange}{#1}}}
\newtcolorbox{definition}[1][]{popbase,colframe=popblue,before upper={\popboxhead{Définition}{popblue}{#1}}}
\newtcolorbox{propriete}[1][]{popbase,colframe=poppurple,before upper={\popboxhead{Propriété}{poppurple}{#1}}}
\newtcolorbox{methode}[1][]{popbase,colframe=popgold,before upper={\popboxhead{Méthode}{popgold}{#1}}}
\newtcolorbox{attention}[1][]{popbase,colframe=poppink,before upper={\popboxhead{Attention}{poppink}{#1}}}
\newtcolorbox{experience}[1][]{popbase,colframe=popblue,colback=popblueL,before upper={\popboxhead{Expérience}{popblue}{#1}}}
% « Le savais-tu ? » : carte violette pleine (contexte, culture, Cameroun)
\newtcolorbox{savaistu}[1][]{popbase,colback=poppurple,colframe=popink,
  fontupper=\popsemi\fontsize{10.4}{14.2}\selectfont\color{white},
  before upper={\poppill{Le savais-tu\,?}{white}{poppurple}\ifx\relax#1\relax\else\hspace{6pt}{\popxbold\color{white}#1}\fi\par\vspace{7pt}}}
% Exercice résolu : énoncé en haut, \tcblower puis la solution sur fond crème
\newcounter{popexo}\newcounter{popexr}
\newtcolorbox{exoresolu}[1][]{popbase,colframe=poporange,colbacklower=popcream,
  segmentation style={popink,line width=1.2pt,dashed},
  before upper={\stepcounter{popexr}\popboxhead{Exercice résolu \thepopexr}{poporange}{#1}},
  before lower={\poppill{Solution}{popink}{popcream}\par\vspace{6pt}}}
% Exercice (non résolu en place) — la correction est dans la partie Corrigés
\newtcolorbox{exercice}[1][]{popbase,colframe=popink,
  before upper={\stepcounter{popexo}\popboxhead{Exercice \thepopexo}{popink}{#1}}}
% Corrigé
\newtcolorbox{corrige}[1][]{popbase,colframe=popgreen,colback=popgreenL,
  before upper={\popboxhead{Corrigé}{popgreen}{#1}}}
% Objectifs / prérequis / bilan
\newtcolorbox{objectifs}{popbase,colframe=popink,colback=poporangeL,
  before upper={\popboxhead{Objectifs de la leçon}{popink}{}}}
\newtcolorbox{prerequis}{popbase,colframe=popink,colback=popblueL,
  before upper={\popboxhead{Prérequis}{popblue}{}}}
\newtcolorbox{bilan}{popbase,colframe=popink,colback=white,boxrule=2.8pt,
  before upper={\popboxhead{Fiche bilan}{poporange}{L'essentiel à connaître pour l'examen}}}
% Figure encadrée avec légende
\newtcolorbox{popfigure}[1][]{enhanced,breakable=false,colback=white,colframe=popline,boxrule=1.4pt,arc=8pt,
  left=10pt,right=10pt,top=10pt,bottom=8pt,before skip=10pt,after skip=12pt,halign=center,
  lower separated=false,halign lower=center,
  fontlower=\popsemi\fontsize{9}{11.5}\selectfont\color{popmuted},
  #1}
\newcommand{\legende}[1]{\par\vspace{6pt}{\popsemi\fontsize{9}{11.5}\selectfont\color{popmuted}#1}}

% ============================================================
% Signature OnBuch+
% ============================================================
\newcommand{\poplogoinline}[1]{{\poplogo\fontsize{#1}{#1}\selectfont\color{popink}OnBuch\color{poporange}+}}
\newcommand{\signatureonbuch}{%
  \begin{tcolorbox}[enhanced,colback=popink,colframe=popink,boxrule=2.2pt,arc=10pt,
      drop shadow=poporange,left=14pt,right=14pt,top=10pt,bottom=10pt,before skip=0pt,after skip=0pt]
    \tikz[baseline=-0.7ex]\node[circle,fill=popgreen,draw=popcream,line width=1.3pt,minimum size=22pt,inner sep=0]{\popcheck{white}};%
    \hspace{9pt}\begin{minipage}[c]{0.62\linewidth}
      {\popxbold\fontsize{12}{13}\selectfont\color{popcream}Signé\ \textbullet\ {\poplogo OnBuch\color{poporange}+}}\par\vspace{2pt}
      {\raggedright\popsemi\fontsize{8}{10.5}\selectfont\color{popline}\MakeUppercase{Équipe pédagogique OnBuch+}\\\MakeUppercase{Conforme au programme officiel MINESEC}\par}
    \end{minipage}\hfill
    {\poplogo\fontsize{22}{22}\selectfont\color{popcream}OnBuch\color{poporange}+}
  \end{tcolorbox}%
}

% ============================================================
% Page de garde
% #1 titre · #2 matière · #3 niveau · #4 module · #5 n° leçon · #6 durée ·
% #7 description / objectifs (texte ou itemize)
% ============================================================
\newcommand{\pagegarde}[7]{%
  \thispagestyle{empty}
  \newgeometry{a4paper,margin=1.7cm,top=1.6cm,bottom=1.4cm}%
  \begin{tikzpicture}[remember picture,overlay]
    \fill[poporange!18] ([xshift=-3.2cm,yshift=-3.6cm]current page.north east) circle (4.6cm);
    \fill[poppurple!14] ([xshift=2.4cm,yshift=7.6cm]current page.south west) circle (3.8cm);
  \end{tikzpicture}%
  {\poplogo\fontsize{30}{30}\selectfont\color{popink}OnBuch\color{poporange}+}\hfill
  \raisebox{0.6ex}{\poppill{Cours signé \textbullet\ Premium}{popink}{popcream}}\par
  \vspace{1pt}\popsquiggle{poppurple}\par
  \vspace{8pt}
  \begin{tcolorbox}[enhanced,colback=poporange,colframe=popink,boxrule=2.8pt,arc=13pt,
      drop shadow=popink,left=18pt,right=18pt,top=12pt,bottom=13pt,after skip=10pt]
    \poppill{#2 \textbullet\ #3}{popink}{popcream}\hspace{5pt}\poppill{#4}{white}{popink}\par\vspace{8pt}
    {\popdisplay\fontsize{14}{14}\selectfont\color{popink}LEÇON #5}\par\vspace{4pt}
    {\raggedright\hyphenpenalty=10000\exhyphenpenalty=10000\popdisplay\fontsize{25}{27}\selectfont\color{popcream}\MakeUppercase{#1}\par}
    \vspace{8pt}
    {\popxbold\fontsize{9.5}{9.5}\selectfont\color{popink}\MakeUppercase{Durée conseillée : #6}}
  \end{tcolorbox}
  \begin{tcolorbox}[enhanced,colback=white,colframe=popink,boxrule=2.2pt,arc=10pt,
      drop shadow=popink,left=16pt,right=16pt,top=10pt,bottom=10pt,after skip=8pt]
    \poppill{Dans cette leçon}{poppurple}{white}\par\vspace{5pt}
    \popsemi\fontsize{9.3}{12.6}\selectfont\color{popink}#7
  \end{tcolorbox}
  \noindent\begin{minipage}[t]{0.487\linewidth}
    \begin{tcolorbox}[enhanced,colback=popgreen,colframe=popink,boxrule=2.2pt,arc=10pt,
        drop shadow=popink,left=11pt,right=11pt,top=10pt,bottom=10pt,height=3.1cm,valign=center]
      \tcbox[colback=white,colframe=popink,boxrule=1.3pt,arc=5pt,boxsep=0pt,nobeforeafter,
        left=3pt,right=3pt,top=3pt,bottom=3pt,valign=center]{\qrcode[height=1.9cm]{https://play.google.com/store/apps/details?id=com.luvvix.onbuch&hl=fr}}%
      \hspace{8pt}\begin{minipage}[c]{0.5\linewidth}
        {\popdisplay\fontsize{11}{12.5}\selectfont\color{white}\MakeUppercase{Télécharge OnBuch}}\par\vspace{4pt}
        {\raggedright\popsemi\fontsize{8.4}{11}\selectfont\color{white}Gratuit \textbullet\ Google Play\\Tuteur IA, annales, concours, résultats\par}
      \end{minipage}
    \end{tcolorbox}
  \end{minipage}\hfill
  \begin{minipage}[t]{0.487\linewidth}
    \begin{tcolorbox}[enhanced,colback=poppurple,colframe=popink,boxrule=2.2pt,arc=10pt,
        drop shadow=popink,left=11pt,right=11pt,top=10pt,bottom=10pt,height=3.1cm,valign=center]
      \tikz[baseline=-0.7ex]\node[circle,fill=white,draw=popink,line width=1.3pt,minimum size=1.6cm,inner sep=0]{\popdisplay\fontsize{24}{24}\selectfont\color{poppurple}+};%
      \hspace{8pt}\begin{minipage}[c]{0.6\linewidth}
        {\popdisplay\fontsize{11}{12.5}\selectfont\color{white}\MakeUppercase{Passe à OnBuch+}}\par\vspace{4pt}
        {\raggedright\popsemi\fontsize{8.4}{11}\selectfont\color{white}Tous les cours signés, Tuteur IA illimité, fiches PDF — onglet OnBuch+ de l'app\par}
      \end{minipage}
    \end{tcolorbox}
  \end{minipage}
  \vspace{8pt}
  \popcontenu
  \vfill
  \signatureonbuch
  \restoregeometry
  \newpage
}

% Rangée « Ce cours contient » (page de garde)
\newcommand{\poptile}[3]{% #1 couleur #2 gros texte #3 légende
  \begin{tcolorbox}[enhanced,colback=#1,colframe=popink,boxrule=2pt,arc=9pt,shadow={2.5pt}{-2.5pt}{0pt}{popink},
      left=6pt,right=6pt,top=9pt,bottom=9pt,halign=center,valign=center,height=1.75cm,width=\linewidth,nobeforeafter]
    {\popdisplay\fontsize{12.5}{13}\selectfont\color{white}\MakeUppercase{#2}}\par\vspace{3pt}
    {\centering\popsemi\fontsize{7.8}{9.5}\selectfont\color{white}#3\par}
  \end{tcolorbox}}
\newcommand{\popcontenu}{%
  \par\noindent{\popxboldls\fontsize{7.5}{7.5}\selectfont\color{popmuted}CE COURS SIGNÉ CONTIENT}\par\vspace{7pt}
  \noindent\begin{minipage}[t]{0.235\linewidth}\poptile{popblue}{Cours}{complet, illustré, pas à pas}\end{minipage}\hfill
  \begin{minipage}[t]{0.235\linewidth}\poptile{poporange}{Méthodes}{et exercices résolus}\end{minipage}\hfill
  \begin{minipage}[t]{0.235\linewidth}\poptile{poppink}{Exercices}{type examen + corrigés}\end{minipage}\hfill
  \begin{minipage}[t]{0.235\linewidth}\poptile{popgreen}{Fiche bilan}{l'essentiel à retenir}\end{minipage}\par}

% Sommaire
\newtcolorbox{sommaire}{enhanced,colback=white,colframe=popink,boxrule=2.2pt,arc=10pt,
  drop shadow=popink,left=18pt,right=18pt,top=13pt,bottom=13pt,before skip=6pt,after skip=20pt,
  before upper={\poppill{Sommaire}{poppurple}{white}\par\vspace{9pt}\popsemi\fontsize{10.6}{14.5}\selectfont\color{popink}}}

% Signature de fin de cours — poussée tout en bas de la dernière page
\newcommand{\findecours}{%
  \par\vfill\nopagebreak
  \begin{center}\popsquiggle{poporange}\end{center}\vspace{4pt}
  \signatureonbuch
}
\AtBeginDocument{\def\llbracket{\mathopen{[\mkern-3.5mu[}}\def\rrbracket{\mathclose{]\mkern-3.5mu]}}} % intervalles d'entiers (glyphe ⟦ absent de la police)
% Glyphes absents de Fira Math : redéfinis à partir de glyphes disponibles
\AtBeginDocument{%
  \def\setminus{\mathbin{\backslash}}\let\smallsetminus\setminus
  \def\vdots{\mathord{\vbox{\baselineskip3.2pt\lineskiplimit0pt\kern4pt\hbox{.}\hbox{.}\hbox{.}}}}%
  \def\ddots{\mathinner{\mkern1mu\raise6.5pt\hbox{.}\mkern2mu\raise3.6pt\hbox{.}\mkern2mu\raise0.7pt\hbox{.}\mkern1mu}}%
  \def\triangle{\mathord{\scalebox{0.9}{$\symup{\Delta}$}}}%
  \def\nabla{\mathord{\raisebox{\depth}{\scalebox{1}[-1]{$\symup{\Delta}$}}}}%
  \def\checkmark{\popcheck{popdark}}%
  \def\star{\mathbin{\ast}}\def\bigstar{\mathord{\ast}}%
  \def\diamond{\mathbin{\scalebox{0.6}{$\Diamond$}}}%
  \def\bigcup{\mathop{\mathchoice{\raisebox{-0.3em}{\scalebox{1.7}{$\cup$}}}{\scalebox{1.25}{$\cup$}}{\cup}{\cup}}\displaylimits}%
  \def\bigcap{\mathop{\mathchoice{\raisebox{-0.3em}{\scalebox{1.7}{$\cap$}}}{\scalebox{1.25}{$\cap$}}{\cap}{\cap}}\displaylimits}%
  \def\lesseqqgtr{\mathrel{\lessgtr}}\def\bigoplus{\mathop{\oplus}\displaylimits}%
}
\providecommand{\ssi}{si et seulement si} % abréviation fréquente
\AtBeginDocument{\providecommand{\wideparen}{\overparen}} % arc de cercle

%% FIGFIX-BEGIN v5 — correctifs globaux des figures (docs/onbuch-generation/tools/figfix)
\usepackage{adjustbox}\usepackage{etoolbox}\tcbuselibrary{raster}
% toute figure TikZ/pgfplots plus large que la ligne est réduite à la largeur disponible
\BeforeBeginEnvironment{tikzpicture}{\begin{adjustbox}{max width=\linewidth}}
\AfterEndEnvironment{tikzpicture}{\end{adjustbox}}
% légendes pgfplots lisibles par-dessus les courbes
\pgfplotsset{every axis/.append style={legend style={fill=white,fill opacity=0.92,text opacity=1,draw=black!15,inner sep=3pt}}}
% étiquettes d'arêtes (syntaxe edge["..."]) avec fond blanc
\tikzset{every edge quotes/.append style={fill=white,inner sep=1.2pt}}
%% FIGFIX-END
\begin{document}

\pagegarde{La transmission des caractères au sein d'une espèce et quelques aspects de la génétique humaine}{SVTEEHB}{Tle C}{Module I : Le Monde Vivant}{N04}{6 h}{Cette leçon étudie les mécanismes de transmission héréditaire des caractères au sein d'une espèce : brassages inter et intrachromosomiques, hérédité autosomique et gonosomique chez l'Homme, mutations et calcul des risques génétiques. Elle fournit les outils d'analyse des croisements et des pedigrees indispensables à l'épreuve de SVT du Baccalauréat C et TI.
\vspace{4pt}
\begin{multicols}{2}\small
\begin{itemize}
\item À la fin de cette leçon, je sais définir les formes alléliques d'un gène et les relations de dominance, codominance et récessivité
\item Analyser et interpréter les résultats d'un croisement de monohybridisme (dominance et codominance) et en déduire les génotypes des parents
\item Analyser et interpréter les résultats d'un dihybridisme à gènes indépendants et mettre en évidence le brassage interchromosomique
\item Exploiter un test-cross pour déterminer le génotype d'un individu de phénotype dominant
\item Mettre en évidence le linkage total et le linkage partiel (crossing-over) et établir une carte factorielle des gènes
\item Construire et exploiter un pedigree pour distinguer une hérédité autosomique d'une hérédité gonosomique
\end{itemize}\end{multicols}}

\begin{sommaire}
\begin{multicols}{2}
\begin{enumerate}
\item Formes alléliques d'un gène et monohybridisme
\item Dihybridisme et brassage interchromosomique
\item Brassage intrachromosomique : linkage et crossing-over
\item Pedigrees et hérédité autosomique humaine
\item Hérédité gonosomique (liée au sexe)
\item Mutations, risques génétiques et prévention
\item Activité d'intégration
\item Méthodes et exercices résolus
\item Exercices
\item Corrigés des exercices
\item Fiche bilan
\end{enumerate}
\end{multicols}
\end{sommaire}

\begin{objectifs}
\begin{itemize}
\item À la fin de cette leçon, je sais définir les formes alléliques d'un gène et les relations de dominance, codominance et récessivité
\item Analyser et interpréter les résultats d'un croisement de monohybridisme (dominance et codominance) et en déduire les génotypes des parents
\item Analyser et interpréter les résultats d'un dihybridisme à gènes indépendants et mettre en évidence le brassage interchromosomique
\item Exploiter un test-cross pour déterminer le génotype d'un individu de phénotype dominant
\item Mettre en évidence le linkage total et le linkage partiel (crossing-over) et établir une carte factorielle des gènes
\item Construire et exploiter un pedigree pour distinguer une hérédité autosomique d'une hérédité gonosomique
\item Citer et caractériser des exemples d'hérédité humaine : albinisme, drépanocytose, achondroplasie, groupes ABO, rhésus, daltonisme, hémophilie
\item Identifier les différents types de mutations géniques et chromosomiques
\item Calculer un risque génétique et proposer des mesures de prévention des anomalies héréditaires
\end{itemize}
\end{objectifs}

\begin{prerequis}
\begin{itemize}
\item Structure du chromosome, ADN et notion de gène (Première)
\item Méiose et formation des gamètes (Première)
\item Caryotype humain : autosomes et gonosomes (Première)
\item Notions de phénotype, génotype, homozygote et hétérozygote (début d'année de Terminale)
\item Probabilités simples et échiquiers de croisement (mathématiques)
\end{itemize}
\end{prerequis}

\newpage

\coursec{Formes alléliques d'un gène et monohybridisme}

\noindent
Dans un village de la région du Centre, un couple dont les deux parents ont un teint normal donne naissance à un enfant albinos : les voisins s'interrogent, certains accusent la mère d'infidélité. À l'hôpital de district, un médecin explique à une famille qu'un enfant drépanocytaire est né de deux parents pourtant en bonne santé. Ailleurs, on constate que l'hémophilie touche presque exclusivement les garçons d'une même lignée. Comment expliquer scientifiquement la réapparition de caractères « cachés », la transmission différentielle selon le sexe, et peut-on prévoir le risque pour les enfants à naître ? C'est tout l'objet de la génétique de la transmission des caractères.

\courssub{Rappels fondamentaux : gène, locus, allèle, génotype et phénotype}

\noindent
Avant d'aborder les lois de la transmission, rappelons le vocabulaire précis qui structure notre raisonnement. Un \cle{gène} est un segment d'ADN portant l'information pour la synthèse d'un produit fonctionnel (protéine ou ARN). Il occupe une position fixe sur un chromosome : c'est son \cle{locus} (pluriel : loci). Chez les eucaryotes diploïdes, les chromosomes vont par paires \cle{homologues} (un d'origine paternelle, un d'origine maternelle). Les deux gènes situés au même locus sur ces chromosomes homologues peuvent être identiques ou différents : ce sont des \cle{allèles}.

\begin{definition}[Allèle]
Une \textbf{forme allélique} (ou allèle) d'un gène est une version particulière de ce gène, occupant le même locus sur des chromosomes homologues, et différant par sa séquence nucléotidique. Les allèles d'un même gène sont notés par la même lettre (souvent la majuscule de l'allèle « normal » ou dominant), avec un exposant ou un indice pour les distinguer (ex. : $A$ et $a$, $I^A$ et $I^B$, $X^h$ et $X^H$).
\end{definition}

\noindent
L'ensemble des allèles présents chez un individu pour un gène donné constitue son \cle{génotype} (ex. : $A/a$ ou $A/A$). L'expression observable de ce génotype, sous l'influence de l'environnement, est le \cle{phénotype} (ex. : graines jaunes, groupe sanguin A, phénotype drépanocytaire). Un individu est \cle{homozygote} pour un gène s'il possède deux allèles identiques ($A/A$ ou $a/a$) ; il est \cle{hétérozygote} s'il possède deux allèles différents ($A/a$). Un homozygote ne produit qu'un seul type de gamète, mais un \textbf{hétérozygote produit deux types de gamètes} (porteurs chacun d'un allèle différent) lors de la méiose.

\courssub{L'approche mendélienne : le monohybridisme et la dominance}

\noindent
Gregor Mendel (1822--1884), moine augustin en Moravie (actuelle République tchèque), a posé les bases de la génétique formelle en croisant des lignées pures de pois (\emph{Pisum sativum}) sur plusieurs générations. Son choix du pois était judicieux : autogame (fécondation naturelle croisée facile à contrôler), cycle court, caractères discrets et bien distincts.

\begin{experience}[Expériences de Mendel sur la couleur des graines du pois (1856--1863)]
\textbf{Matériel :} Lignées pures de pois à graines jaunes (phénotype « jaune ») et lignées pures à graines vertes (phénotype « vert »). \textbf{Protocole :}
\begin{enumerate}
    \item \textbf{Génération P (Parents) :} Croisement réciproque (pollen jaune $\times$ ovaire vert, puis pollen vert $\times$ ovaire jaune) entre lignées pures.
    \item \textbf{Génération F1 :} Auto-fécondation (ou croisement entre individus F1) des plantes issues de P.
    \item \textbf{Génération F2 :} Auto-fécondation des plantes F1. Comptage phénotypique rigoureux sur de grands effectifs.
\end{enumerate}
\textbf{Observations (résultats historiques de Mendel) :}
\begin{itemize}
    \item Génération F1 : \textbf{100\%} des graines sont \textbf{jaunes}. Le caractère « vert » a disparu.
    \item Génération F2 : Sur \num{8023} graines, \num{6022} sont jaunes et \num{2001} sont vertes. Soit un rapport \textbf{3 jaunes : 1 vert} ($\approx 2,99 : 1$).
\end{itemize}
\textbf{Interprétation :} Le caractère « jaune » est \cle{dominant} sur le caractère « vert » qui est \cle{récessif}. Les facteurs (gènes) se séparent lors de la formation des gamètes (loi de la ségrégation) et se recombinent au hasard à la fécondation.
\legende{Expérience fondatrice du monohybridisme : croisement de deux lignées pures différant par un seul caractère (couleur des graines). La réapparition du caractère récessif en F2 dans une proportion 1/4 prouve sa conservation latente en F1.}
\end{experience}

\begin{savaistu}[Mendel, le génie méconnu]
Mendel présenta ses résultats en 1865 et les publia en 1866 dans les \emph{Verhandlungen des naturforschenden Vereines in Brünn}. Son travail, trop mathématique pour les biologistes de l'époque, passa inaperçu. Ce n'est qu'en 1900, soit 34 ans plus tard, que trois botanistes (Hugo de Vries, Carl Correns, Erich von Tschermak) redécouvrirent indépendamment ses lois. Au Cameroun, comme partout, la génétique mendélienne est la clé pour comprendre l'hérédité de l'albinisme, de la drépanocytose ou des groupes sanguins.
\end{savaistu}

\noindent
L'interprétation chromosomique moderne (Sutton, Boveri, Morgan, ~1902--1915) ancre ces lois dans la cytologie : les allèles se séparent lors de l'\textbf{anaphase I de la méiose} (ségrégation des chromosomes homologues), puis les chromatides sœurs se séparent en anaphase II. Chaque gamète reçoit un seul allèle pour ce gène.

\begin{definition}[Loi de l'uniformité des hybrides de première génération (1ère loi de Mendel)]
Lorsqu'on croise deux individus homozygotes de phénotypes différents pour un même caractère (lignées pures), tous les descendants de la première génération (F1) sont phénotypiquement identiques et ressemblent à l'un des parents : celui qui porte l'allèle \textbf{dominant}.
\end{definition}

\begin{definition}[Loi de la ségrégation (2ème loi de Mendel)]
Chez un hétérozygote ($A/a$), les deux allèles se séparent lors de la méiose pour se retrouver dans des gamètes différents. À la fécondation, les allèles se recombinent au hasard. En croisant deux hétérozygotes F1 ($A/a \times A/a$), on observe en F2 un rapport phénotypique \textbf{3 (dominant) : 1 (récessif)} et un rapport génotypique \textbf{1 $A/A$ : 2 $A/a$ : 1 $a/a$}.
\end{definition}

\noindent
Pour visualiser ces recombinaisons aléatoires, on utilise l'\cle{échiquier de croisement} (ou carré de Punnett).

\begin{popfigure}
\begin{tikzpicture}[scale=0.9, every node/.style={font=\small}]
% Grille 3x3 : en-têtes + 2x2 cases
\draw[popdark, thick] (0,0) grid (3,3);
% En-têtes colonnes (gamètes mâle) - ligne y=2 à 3
\node[popblue, font=\bfseries] at (0.5,2.5) {$A$};
\node[popblue, font=\bfseries] at (1.5,2.5) {$a$};
\node[popmuted, anchor=east] at (0,2.5) {$\sigma$};
% En-têtes lignes (gamètes femelle) - colonne x=0 à 1
\node[poporange, font=\bfseries, rotate=90] at (-0.3,2) {$A$};
\node[poporange, font=\bfseries, rotate=90] at (-0.3,1) {$a$};
\node[popmuted, anchor=north] at (-0.3,1.5) {$\varphi$};
% Contenu cases (centres aux coordonnées .5, .5 ; 1.5, .5 ; .5, 1.5 ; 1.5, 1.5)
\node at (0.5,1.5) {$A/A$ \textbf{(Jaune)}};
\node at (1.5,1.5) {$A/a$ \textbf{(Jaune)}};
\node at (0.5,0.5) {$A/a$ \textbf{(Jaune)}};
\node at (1.5,0.5) {$a/a$ \textbf{(Vert)}};
% Légendes proportions
\node[popdark, align=center] at (1.5,-0.6) {Génotypes F2 : 1 $A/A$ : 2 $A/a$ : 1 $a/a$};
\node[popdark, align=center] at (1.5,-1.3) {Phénotypes F2 : \textbf{3 Jaunes} : \textbf{1 Vert}};
\end{tikzpicture}
\legende{Échiquier de croisement F1 $\times$ F1 ($A/a \times A/a$) pour un caractère à dominance complète. Les gamètes sont portés par les axes ; les cases internes donnent les génotypes zygotiques et le phénotype correspondant.}
\end{popfigure}

\begin{methode}[Rédiger un croisement monohybride complet (attendu au Bac)]
Pour tout croisement, suivez ces 5 étapes rigoureuses :
\begin{enumerate}
    \item \textbf{Choix des symboles :} Définir la lettre (ex. $A$ = allèle dominant « jaune », $a$ = allèle récessif « vert »).
    \item \textbf{Génotypes des parents (P) :} Préciser s'ils sont homozygotes ou hétérozygotes (ex. $A/A \times a/a$).
    \item \textbf{Gamètes :} Lister les types de gamètes produits par chaque parent (méiose). Homozygote $\rightarrow$ 1 type ; Hétérozygote $\rightarrow$ 2 types (ségrégation).
    \item \textbf{Échiquier de croisement :} Croiser \textbf{tous} les gamètes mâles (colonnes) avec \textbf{tous} les gamètes femelles (lignes). Remplir les cases avec les génotypes zygotiques.
    \item \textbf{Bilan :} Compter les cases pour donner les proportions \textbf{génotypiques} et \textbf{phénotypiques} (en fractions, pourcentages ou rapports). Conclure.
\end{enumerate}
\end{methode}

\begin{exoresolu}[Croisement de pois : vérification des proportions mendéliennes]
\textbf{Énoncé :} Mendel a obtenu en F2 \num{6022} graines jaunes et \num{2001} graines vertes. Vérifier la conformité au rapport théorique 3/1.
\textbf{Solution détaillée :}
Effectif total $N = 6022 + 2001 = 8023$.
Fréquence observée (jaune) $f_J = 6022 / 8023 \approx 0,7506$.
Fréquence observée (vert) $f_V = 2001 / 8023 \approx 0,2494$.
Fréquences théoriques : $3/4 = 0,75$ et $1/4 = 0,25$.
Les valeurs observées sont extrêmement proches des valeurs théoriques (écart $< 0,1\%$). L'hypothèse mendélienne (dominance complète, ségrégation équitable, fécondation aléatoire) est validée par l'expérience.
\tcblower
\textbf{Remarque :} Au Bac, on ne demande pas de test du $\chi^2$ (khî-deux) en SVT, mais on doit savoir comparer les proportions observées aux proportions théoriques et conclure à la validation ou non du modèle.
\end{exoresolu}

\courssub{Dominance incomplète et codominance : quand les allèles s'expriment différemment}

\noindent
La dominance complète ($A$ masque $a$ chez $A/a$) n'est qu'un cas particulier. Deux autres relations alléliques majeures existent, illustrant que la relation entre génotype et phénotype n'est pas binaire.

\begin{experience}[Dominance incomplète : la Belle-de-nuit (\emph{Mirabilis jalapa})]
\textbf{Observation :} Croisement de lignées pures à fleurs rouges ($R/R$) $\times$ fleurs blanches ($r/r$).
\textbf{F1 :} 100\% fleurs \textbf{roses} (phénotype intermédiaire).
\textbf{F2 (auto-fécondation F1) :} 1/4 Rouge : 1/2 Rose : 1/4 Blanc. Rapport phénotypique \textbf{1 : 2 : 1} (identique au rapport génotypique).
\textbf{Interprétation :} L'allèle $R$ code pour un pigment rouge, l'allèle $r$ ne code pour aucun pigment (ou un pigment non fonctionnel). Chez $R/r$, la dose de pigment est moitié moindre $\rightarrow$ phénotype intermédiaire (rose). \textbf{Aucun allèle n'est dominant.}
\legende{Dominance incomplète : l'hétérozygote présente un phénotype intermédiaire entre les deux homozygotes. Le rapport phénotypique F2 devient 1:2:1, superposable au rapport génotypique.}
\end{experience}

\begin{experience}[Codominance : le système MN du sang humain]
\textbf{Contexte :} Découvert par Landsteiner et Levine (1927). Deux allèles $L^M$ et $L^N$ (notation $M$ et $N$ souvent utilisée) déterminent des antigènes de surface des globules rouges (glycophorines A et B).
\textbf{Phénotypes observés :}
\begin{itemize}
    \item $L^M/L^M$ $\rightarrow$ Groupe M (antigène M seul).
    \item $L^N/L^N$ $\rightarrow$ Groupe N (antigène N seul).
    \item $L^M/L^N$ $\rightarrow$ Groupe \textbf{MN} (les \textbf{deux} antigènes M \textbf{et} N sont présents simultanément à la surface des globules rouges).
\end{itemize}
\textbf{Interprétation :} Les deux allèles s'expriment \textbf{pleinement et indépendamment} chez l'hétérozygote. Il n'y a pas de phénotype intermédiaire, mais une \textbf{co-expression}.
\legende{Codominance : les deux allèles s'expriment simultanément et distinctement chez l'hétérozygote. Le rapport phénotypique F2 est aussi 1:2:1, mais pour une raison différente de la dominance incomplète.}
\end{experience}

\begin{attention}[Piège majeur : Codominance $\neq$ Dominance incomplète]
\begin{itemize}
    \item \textbf{Dominance incomplète :} Phénotype \textbf{intermédiaire} (mélange quantitatif). Ex. : Belle-de-nuit (Rouge + Blanc = Rose), drépanocytose (HbA/HbS = trait drépanocytaire asymptomatique \emph{mais} biologiquement intermédiaire : résistance au paludisme + hémolyse modérée en hypoxie).
    \item \textbf{Codominance :} \textbf{Co-expression} qualitative des deux produits alléliques. Ex. : Groupes sanguins MN, Groupes sanguins ABO ($I^A/I^B$ = groupe AB : antigènes A \textbf{et} B), Protéines sériques (haptoglobine).
\end{itemize}
Au Bac, si on vous demande de « distinguer », citez l'exemple du groupe AB (codominance) vs Belle-de-nuit (dominance incomplète) et précisez : « En codominance, les deux produits alléliques sont détectables ; en dominance incomplète, le phénotype est un mélange/atténuation. »
\end{attention}

\begin{popfigure}
\begin{tikzpicture}
\begin{axis}[
    ybar,
    bar width=12pt,
    width=\linewidth,
    height=8cm,
    ymin=0, ymax=1.1,
    xtick={1,2,3},
    xticklabels={Dominance complète, Dominance incomplète, Codominance},
    ylabel={Fréquence phénotypique en F2},
    ymajorgrids=true,
    grid style={popline},
    legend style={at={(0.5,-0.15)}, anchor=north, legend columns=3, font=\small},
    nodes near coords,
    nodes near coords align={vertical},
    every node near coord/.append style={font=\tiny, popdark}
]
\addplot[popblue!70, fill=popblue!50] coordinates {(1,0.75) (2,0.25) (3,0.25)};
\addplot[poporange!70, fill=poporange!50] coordinates {(1,0.25) (2,0.50) (3,0.50)};
\addplot[popgreen!70, fill=popgreen!50] coordinates {(1,0) (2,0.25) (3,0.25)};
\legend{Phénotype « Dominant / Hom. 1 » (ex: Jaune, Rouge, M), Phénotype « Hétérozygote / Intermédiaire » (ex: Rose, MN, AB), Phénotype « Récessif / Hom. 2 » (ex: Vert, Blanc, N)}
\end{axis}
\end{tikzpicture}
\legende{Comparaison des proportions phénotypiques en F2 (auto-fécondation d'hétérozygotes) selon le type de relation allélique. En dominance complète : 3/4 - 1/4. En dominance incomplète et codominance : 1/4 - 2/4 - 1/4 (le rapport phénotypique = rapport génotypique).}
\end{popfigure}

\courssub{Allèles multiples : le système ABO, modèle de référence}

\noindent
Jusqu'ici, nous avons considéré un gène à \textbf{deux} allèles dans la population. Or, une population peut abriter \textbf{plus de deux allèles} pour un même gène : ce sont des \cle{allèles multiples}. Un individu diploïde ne peut en posséder que \textbf{deux} (un sur chaque chromosome homologue), mais la diversité allélique dans la population est supérieure à 2.

\begin{definition}[Allèles multiples]
On parle d'allèles multiples lorsqu'un gène existe sous plus de deux formes alléliques différentes au sein d'une population. La série allélique est notée par une même lettre de base avec des exposants (ex. : $I^A, I^B, i$ pour ABO).
\end{definition}

\noindent
Le système \textbf{ABO} (découvert par Karl Landsteiner, 1900, prix Nobel 1930) est l'exemple type. Trois allèles : $I^A$, $I^B$, $i$ (ou $I^O$).
\begin{itemize}
    \item $I^A$ et $I^B$ sont \textbf{codominants} l'un par rapport à l'autre (groupe AB = antigènes A et B).
    \item $I^A$ et $I^B$ sont \textbf{dominants} sur $i$ (récessif).
    \item $i$ ne produit pas d'antigène fonctionnel (antigène H non modifié).
\end{itemize}

\begin{center}
\begin{tabularx}{\linewidth}{|l|X|X|X|}
\hline
\rowcolor{popblueL}
\textbf{Génotype} & \textbf{Phénotype (Groupe)} & \textbf{Antigènes sur GR} & \textbf{Anticorps dans le sérum (anti-)} \\
\hline
$I^A/I^A$ ou $I^A/i$ & A & A & anti-B \\
$I^B/I^B$ ou $I^B/i$ & B & B & anti-A \\
$I^A/I^B$ & AB & A \textbf{et} B & \textbf{aucun} (donneur universel de plasma) \\
$i/i$ & O & \textbf{aucun} (H seul) & anti-A \textbf{et} anti-B (donneur universel de GR) \\
\hline
\end{tabularx}
\end{center}

\begin{exoresolu}[Prévoir les groupes sanguins possibles d'un enfant]
\textbf{Énoncé :} Un homme de groupe A (phénotype) et une femme de groupe B ont un enfant de groupe O. Déterminer les génotypes des parents et la probabilité d'avoir un enfant de groupe AB.
\textbf{Solution détaillée :}
\begin{enumerate}
    \item Enfant groupe O $\rightarrow$ génotype obligatoire $i/i$. Il a reçu un allèle $i$ de chaque parent.
    \item Père groupe A + allèle $i$ transmis $\rightarrow$ Génotype père = $I^A/i$.
    \item Mère groupe B + allèle $i$ transmis $\rightarrow$ Génotype mère = $I^B/i$.
    \item Croisement : $I^A/i \times I^B/i$.
    \item Gamètes père : $I^A$, $i$. Gamètes mère : $I^B$, $i$.
    \item Échiquier :
    \begin{center}
    \begin{tabular}{|c|c|c|}\hline
    & $I^B$ & $i$ \\ \hline
    $I^A$ & $I^A/I^B$ (AB) & $I^A/i$ (A) \\ \hline
    $i$ & $I^B/i$ (B) & $i/i$ (O) \\ \hline
    \end{tabular}
    \end{center}
    \item Proportions : 1/4 A, 1/4 B, 1/4 AB, 1/4 O.
    \item Probabilité groupe AB = \textbf{1/4 (25\%)}.
\end{enumerate}
\tcblower
\textbf{Application clinique :} Ce raisonnement est la base des \emph{preuves de paternité} par exclusion (ex. : père O, mère O $\rightarrow$ enfant ne peut être que O ; père AB, mère O $\rightarrow$ enfant A ou B, jamais O ni AB). Au Cameroun, les centres de transfusion sanguine (CNTS) utilisent ces règles pour la compatibilité donneur/receveur.
\end{exoresolu}

\courssub{Allèles létaux : une nuance importante pour le Bac}

\noindent
Certains allèles, à l'état homozygote, entraînent la mort de l'individu avant la naissance (léthalité embryonnaire) ou peu après. Cela modifie les proportions mendéliennes attendues en F2.

\begin{exemplebox}[Allèle létal récessif : la souris jaune ($A^y$)]
L'allèle $A^y$ (jaune) est dominant sur $a$ (gris/agouti) pour la couleur du pelage.
\begin{itemize}
    \item $A^y/a$ : Souris \textbf{jaune} (viable).
    \item $a/a$ : Souris \textbf{grise} (viable).
    \item $A^y/A^y$ : \textbf{Létal embryonnaire} (mort in utero, résorption).
\end{itemize}
\textbf{Croisement :} Souris jaune $\times$ Souris jaune ($A^y/a \times A^y/a$).
Échiquier théorique : 1 $A^y/A^y$ (mort) : 2 $A^y/a$ (jaune) : 1 $a/a$ (grise).
\textbf{Proportions observées chez les survivants (nés vivants) :} \textbf{2 Jaunes : 1 Grise} (soit 2/3 jaune, 1/3 grise). Le rapport 3/1 devient 2/1.
\legende{Exemple d'allèle létal récessif conditionnel (homozygote). La mortalité des homozygotes dominants fausse les proportions phénotypiques attendues.}
\end{exemplebox}

\begin{attention}[Allèle létal : ne pas confondre génotype zygotique et phénotype observé]
Au Bac, si un énoncé donne des proportions F2 de type 2/3 - 1/3 (ou 2:1), \textbf{pensez immédiatement à un allèle létal récessif (homozygote mortel)}. Il faut :
\begin{enumerate}
    \item Identifier l'allèle létal (souvent l'allèle dominant pour le caractère visible).
    \item Écrire l'échiquier complet (avec la case létale).
    \item Calculer les proportions \textbf{parmi les individus viables} (néonés ou adultes observés).
    \item Préciser : « Le rapport génotypique zygotique est 1:2:1, mais le rapport phénotypique observé est 2:1 car le génotype $A^y/A^y$ est létal. »
\end{enumerate}
\end{attention}

\courssub{Le test-cross (test de descendance) : vérifier le génotype d'un phénotype dominant}

\noindent
Un individu exprimant le phénotype dominant peut être homozygote ($A/A$) ou hétérozygote ($A/a$). Comment le savoir sans séquençage ? En le croisant avec un individu \textbf{homozygote récessif} ($a/a$). C'est le \cle{test-cross} (ou croisement test).

\begin{propriete}[Intérêt du test-cross]
\begin{itemize}
    \item Si l'individu testé est \textbf{homozygote dominant} ($A/A$) : 100\% de la descendance exprime le phénotype dominant. Gamètes : que des $A$. Croisement $A/A \times a/a \rightarrow$ 100\% $A/a$.
    \item Si l'individu testé est \textbf{hétérozygote} ($A/a$) : La descendance ségrège \textbf{1 dominant : 1 récessif}. Gamètes : $A$ et $a$. Croisement $A/a \times a/a \rightarrow$ 1/2 $A/a$ (dom) : 1/2 $a/a$ (rec).
\end{itemize}
C'est une application directe de la ségrégation méiotique. Ce principe sera généralisé en dihybridisme (section 2) pour détecter le linkage.
\end{propriete}

\begin{methode}[Réaliser et interpréter un test-cross]
\begin{enumerate}
    \item Identifier le phénotype dominant à tester (ex. : graines jaunes, groupe sanguin A, phénotype normal pour drépanocytose).
    \item Croiser cet individu avec un homozygote récessif pur (ex. : graines vertes $a/a$, groupe O $i/i$, drépanocytaire majeur $S/S$).
    \item Observer la descendance (effectif suffisant, $n \ge 20$ idéalement).
    \item \textbf{Conclusion :}
    \begin{itemize}
        \item Que dominants $\rightarrow$ Parent testé = Homozygote dominant.
        \item 1 Dominant : 1 Récessif $\rightarrow$ Parent testé = Hétérozygote.
    \end{itemize}
\end{enumerate}
\end{methode}

\begin{exoresolu}[Test-cross et conseil génétique : drépanocytose au Cameroun]
\textbf{Contexte :} La drépanocytose (anémie falciforme) est une maladie autosomique récessive fréquente au Cameroun (prévalence du trait $AS \approx 20-25\%$ en zone forestière). Allèle $A$ (HbA normal) dominant sur $S$ (HbS muté). $A/A$ = normal, $A/S$ = trait drépanocytaire (asymptomatique, résistant au paludisme), $S/S$ = drépanocytaire majeur (malade).
\textbf{Énoncé :} Un homme phénotypiquement normal (non drépanocytaire) épouse une femme drépanocytaire majeure ($S/S$). Ils ont 4 enfants : 2 drépanocytaires majeurs ($S/S$), 2 porteurs du trait ($A/S$). Quel est le génotype du père ? Quel était le risque pour chaque grossesse ?
\textbf{Solution détaillée :}
\begin{enumerate}
    \item Mère $S/S$ $\rightarrow$ Gamètes : 100\% $S$.
    \item Enfants malades ($S/S$) : ont reçu $S$ du père $\rightarrow$ Le père possède l'allèle $S$.
    \item Père phénotypiquement normal + allèle $S$ $\rightarrow$ Génotype père = \textbf{$A/S$ (Hétérozygote, porteur sain)}.
    \item C'est un \textbf{test-cross} déguisé : Père ($A/?$) $\times$ Mère ($S/S$). Descendance 1 $A/S$ : 1 $S/S$ $\rightarrow$ Père hétérozygote.
    \item Risque par grossesse : \textbf{1/2 (50\%)} enfant malade ($S/S$), \textbf{1/2} porteur sain ($A/S$).
\end{enumerate}
\tcblower
\textbf{Message de prévention :} Au Cameroun, le dépistage prénuptial (électrophorèse de l'hémoglobine) est recommandé. Un couple $A/S \times A/S$ a 25\% de risque d'enfant $S/S$ par grossesse. Le conseil génétique informe, ne décide pas à la place du couple.
\end{exoresolu}

\begin{aretenir}[Bilan : Monohybridisme et relations alléliques]
\begin{itemize}
    \item \textbf{Gène / Locus / Allèle :} Unité de l'hérédité, position, version.
    \item \textbf{Dominance complète :} $A/a \rightarrow$ Phénotype $A$. F2 : 3/4 Dom : 1/4 Rec. (Rapport phénotypique $\neq$ génotypique).
    \item \textbf{Dominance incomplète :} $A/a \rightarrow$ Phénotype intermédiaire. F2 : 1/4 : 2/4 : 1/4. (Rapport phénotypique = génotypique).
    \item \textbf{Codominance :} $A/a \rightarrow$ Co-expression des deux allèles. F2 : 1/4 : 2/4 : 1/4. (Rapport phénotypique = génotypique).
    \item \textbf{Allèles multiples :} $>2$ allèles dans la population (ex. ABO : $I^A, I^B, i$). Relations : codominance $I^A/I^B$, dominance sur $i$.
    \item \textbf{Allèle létal (homozygote) :} Mortalité précoce. F2 observé : 2/3 : 1/3 (au lieu de 3/4 : 1/4).
    \item \textbf{Test-cross :} Individu phénotype dominant $\times$ homozygote récessif. Permet de discriminer $A/A$ (100\% Dom) de $A/a$ (1 Dom : 1 Rec).
    \item \textbf{Base cellulaire :} Ségrégation des allèles en Anaphase I de la méiose (séparation chromosomes homologues).
\end{itemize}
\end{aretenir}

\noindent
Nous avons maintenant les outils pour analyser la transmission d'un seul caractère. Mais les caractères sont portés par des chromosomes qui en portent des milliers. Comment se comportent \textbf{deux gènes} (ou plus) situés sur \textbf{chromosomes différents} (indépendants) ou sur le \textbf{même chromosome} (liés) ? C'est l'objet de la section suivante : le dihybridisme et le brassage inter- puis intrachromosomique.

\coursec{Dihybridisme et brassage interchromosomique}

Après avoir étudié la transmission d'un seul caractère (monohybridisme) et découvert que les allèles d'un même gène se séparent lors de la méiose (loi de la pureté des gamètes), nous nous posons maintenant une question naturelle : \textbf{que se passe-t-il lorsque l'on suit simultanément deux caractères différents ?} Les allèles de ces deux gènes restent-ils toujours associés comme dans le phénotype parental, ou bien se mélangent-ils librement ? C'est la question que Mendel a résolue par ses expériences de dihybridisme, révélant le \textbf{brassage interchromosomique}, source majeure de diversité génétique.

\courssub{L'expérience fondatrice de Mendel : le dihybridisme chez le pois}

\emph{Observation.} Mendel a croisé deux lignées pures de pois (\emph{Pisum sativum}) différant par deux caractères : la couleur du tégument des graines (jaune dominant $J$ / vert récessif $j$) et la texture des graines (lisse dominant $A$ / ridée récessif $a$).
\begin{itemize}
    \item Parental (P) : Graines \textbf{jaunes lisses} ($JJAA$) $\times$ Graines \textbf{vertes ridées} ($jjaa$).
    \item F1 : 100\% graines \textbf{jaunes lisses} ($JjAa$) — les allèles dominants masquent les récessifs.
\end{itemize}

\emph{Expérience décisive : l'autofécondation de la F1 (F1 $\times$ F1).} Mendel a semé les graines F1 et a récolté la génération F2. Il a compté 556 graines et obtenu les effectifs suivants :
\begin{center}
\begin{tabularx}{\linewidth}{|l|X|X|X|X|}\hline
Phénotype & Jaune lisse & Jaune ridée & Verte lisse & Verte ridée \\ \hline
Effectif observé & 315 & 101 & 108 & 32 \\ \hline
Proportion observée & 0,567 & 0,182 & 0,194 & 0,058 \\ \hline
Proportion théorique & 9/16 = 0,5625 & 3/16 = 0,1875 & 3/16 = 0,1875 & 1/16 = 0,0625 \\ \hline
\end{tabularx}
\end{center}

\begin{experience}[Le dihybridisme de Mendel (Pois, 1865)]
\textbf{Matériel :} Lignées pures de pois : jaune lisse ($JJAA$) et verte ridée ($jjaa$).
\textbf{Protocole :}
\begin{enumerate}
    \item Croisement P : $JJAA \times jjaa$ $\rightarrow$ F1 toute $JjAa$ (jaune lisse).
    \item Autofécondation de la F1 : $JjAa \times JjAa$ $\rightarrow$ F2.
    \item Comptage phénotypique rigoureux sur grande descendance ($N = 556$).
\end{enumerate}
\textbf{Observations :} Réapparition des deux phénotypes parentaux (jaune lisse, verte ridée) \textbf{et} apparition de deux phénotypes \textbf{recombinés} (jaune ridée, verte lisse) dans des proportions très proches de $9 : 3 : 3 : 1$.
\textbf{Interprétation :} Les deux couples d'allèles ($J/j$ et $A/a$) se comportent \textbf{indépendamment} l'un de l'autre lors de la formation des gamètes. Le double hétérozygote $JjAa$ produit \textbf{4 types de gamètes équiprobables} : $JA$, $Ja$, $jA$, $ja$ (chacun à 25\%).
\textbf{Conclusion :} \textbf{Loi de l'indépendance des caractères (2\textsuperscript{e} loi de Mendel)} : lorsque deux gènes sont portés par des chromosomes différents (ou sont très éloignés sur le même chromosome), la ségrégation des allèles de l'un est indépendante de la ségrégation des allèles de l'autre.
\tcblower
\textbf{Échiquier de croisement (carré de Punnett) 4 $\times$ 4 :}
\end{experience}

\begin{popfigure}
\begin{tikzpicture}[scale=0.65, every node/.style={font=\small}]
% Styles
\tikzset{
    cell/.style={minimum width=1.6cm, minimum height=1.6cm, draw=popline, thick, anchor=south west},
    gamlabel/.style={font=\small\bfseries, text=popdark},
    genolabel/.style={font=\scriptsize, text=popdark},
    legbox/.style={draw=popline, thick, minimum width=0.6cm, minimum height=0.35cm, rounded corners=2pt}
}
% Gametes
\def\glist{{"JA","Ja","jA","ja"}}
% Colors per phenotype
% 0: Yellow Smooth (J_ A_) -> popgoldL
% 1: Yellow Wrinkled (J_ aa) -> poporangeL
% 2: Green Smooth (jj A_) -> popgreenL
% 3: Green Wrinkled (jj aa) -> poppinkL

% Draw grid and cells
\foreach \r in {0,1,2,3} {
    \foreach \c in {0,1,2,3} {
        % Determine phenotype index 0,1,2,3
        % Gamete row \r, Gamete col \c
        % Extract alleles
        % Row gamete
        \pgfmathparse{\glist[\r]} \let\rgam\pgfmathresult
        % Col gamete
        \pgfmathparse{\glist[\c]} \let\cgam\pgfmathresult
        % Alleles: J/j from char 0, A/a from char 1
        \edef\rJ{\rgam:0:1} \edef\rA{\rgam:1:1} % Not easy in tikz math, use ifthenelse logic via macros
    }
}
% Simpler: Hardcode the 16 genotypes/phenotypes for reliability in TikZ
% Row 0 (Female JA)
\foreach \c/\cgam/\pheno/\geno in {0/JA/0/JA, 1/Ja/0/Ja, 2/jA/0/jA, 3/ja/0/ja} {
    \pgfmathparse{\c*1.6} \let\x\pgfmathresult
    \pgfmathparse{0*1.6} \let\y\pgfmathresult
    \ifnum\pheno=0 \fill[popgoldL] (\x,\y) rectangle (\x+1.6,\y+1.6); \fi
    \ifnum\pheno=1 \fill[poporangeL] (\x,\y) rectangle (\x+1.6,\y+1.6); \fi
    \ifnum\pheno=2 \fill[popgreenL] (\x,\y) rectangle (\x+1.6,\y+1.6); \fi
    \ifnum\pheno=3 \fill[poppinkL] (\x,\y) rectangle (\x+1.6,\y+1.6); \fi
    \node[genolabel] at (\x+0.8,\y+0.8) {$JjAa$}; % Placeholder, will fix below
}
% Actually, let's just manually place the 16 nodes with correct genotypes for perfect control.
% Coordinates: bottom-left (0,0) to top-right (6.4, 6.4). Row 0 bottom (Female JA), Row 3 top (Female ja).
% Col 0 left (Male JA), Col 3 right (Male ja).

% Row 0 (Female JA) - y = 0 to 1.6
\fill[popgoldL] (0,0) rectangle (1.6,1.6); \node[genolabel] at (0.8,0.8) {$JJAA$}; % JA x JA
\fill[popgoldL] (1.6,0) rectangle (3.2,1.6); \node[genolabel] at (2.4,0.8) {$JjAA$}; % JA x Ja
\fill[popgoldL] (3.2,0) rectangle (4.8,1.6); \node[genolabel] at (4.0,0.8) {$JJAa$}; % JA x jA
\fill[popgoldL] (4.8,0) rectangle (6.4,1.6); \node[genolabel] at (5.6,0.8) {$JjAa$}; % JA x ja

% Row 1 (Female Ja) - y = 1.6 to 3.2
\fill[popgoldL] (0,1.6) rectangle (1.6,3.2); \node[genolabel] at (0.8,2.4) {$JjAA$}; % Ja x JA
\fill[popgoldL] (1.6,1.6) rectangle (3.2,3.2); \node[genolabel] at (2.4,2.4) {$JJAA$}; % Ja x Ja -> JJAA? No Ja x Ja -> JJ Aa? Wait. Ja (J,a) x Ja (J,a) -> JJ aa. Yellow Wrinkled.
% CORRECTION:
% Gametes: JA (J,A), Ja (J,a), jA (j,A), ja (j,a)
% Row 1 Female Ja (J,a)
% Col 0 Male JA (J,A) -> Jj Aa (Yellow Smooth) -> popgoldL
% Col 1 Male Ja (J,a) -> JJ aa (Yellow Wrinkled) -> poporangeL
% Col 2 Male jA (j,A) -> Jj Aa (Yellow Smooth) -> popgoldL
% Col 3 Male ja (j,a) -> Jj aa (Yellow Wrinkled) -> poporangeL

\fill[popgoldL] (0,1.6) rectangle (1.6,3.2); \node[genolabel] at (0.8,2.4) {$JjAa$};
\fill[poporangeL] (1.6,1.6) rectangle (3.2,3.2); \node[genolabel] at (2.4,2.4) {$JJaa$};
\fill[popgoldL] (3.2,1.6) rectangle (4.8,3.2); \node[genolabel] at (4.0,2.4) {$JjAa$};
\fill[poporangeL] (4.8,1.6) rectangle (6.4,3.2); \node[genolabel] at (5.6,2.4) {$Jjaa$};

% Row 2 (Female jA) - y = 3.2 to 4.8
% Female jA (j,A)
% Col 0 Male JA (J,A) -> Jj AA (Yellow Smooth) -> popgoldL
% Col 1 Male Ja (J,a) -> Jj Aa (Yellow Smooth) -> popgoldL
% Col 2 Male jA (j,A) -> jj AA (Green Smooth) -> popgreenL
% Col 3 Male ja (j,a) -> jj Aa (Green Smooth) -> popgreenL
\fill[popgoldL] (0,3.2) rectangle (1.6,4.8); \node[genolabel] at (0.8,4.0) {$JjAA$};
\fill[popgoldL] (1.6,3.2) rectangle (3.2,4.8); \node[genolabel] at (2.4,4.0) {$JjAa$};
\fill[popgreenL] (3.2,3.2) rectangle (4.8,4.8); \node[genolabel] at (4.0,4.0) {$jjAA$};
\fill[popgreenL] (4.8,3.2) rectangle (6.4,4.8); \node[genolabel] at (5.6,4.0) {$jjAa$};

% Row 3 (Female ja) - y = 4.8 to 6.4
% Female ja (j,a)
% Col 0 Male JA (J,A) -> Jj Aa (Yellow Smooth) -> popgoldL
% Col 1 Male Ja (J,a) -> Jj aa (Yellow Wrinkled) -> poporangeL
% Col 2 Male jA (j,A) -> jj Aa (Green Smooth) -> popgreenL
% Col 3 Male ja (j,a) -> jj aa (Green Wrinkled) -> poppinkL
\fill[popgoldL] (0,4.8) rectangle (1.6,6.4); \node[genolabel] at (0.8,5.6) {$JjAa$};
\fill[poporangeL] (1.6,4.8) rectangle (3.2,6.4); \node[genolabel] at (2.4,5.6) {$Jjaa$};
\fill[popgreenL] (3.2,4.8) rectangle (4.8,6.4); \node[genolabel] at (4.0,5.6) {$jjAa$};
\fill[poppinkL] (4.8,4.8) rectangle (6.4,6.4); \node[genolabel] at (5.6,5.6) {$jjaa$};

% Grid lines
\draw[thick, popline] (0,0) grid[step=1.6] (6.4,6.4);
\draw[very thick, popdark] (0,0) rectangle (6.4,6.4);

% Gamete Labels (Male top, Female left)
\foreach \c/\cgam in {0/JA, 1/Ja, 2/jA, 3/ja} {
    \pgfmathparse{\c*1.6+0.8} \let\x\pgfmathresult
    \node[gamlabel, anchor=south, rotate=45] at (\x, 6.6) {\cgam};
}
\foreach \r/\rgam in {0/JA, 1/Ja, 2/jA, 3/ja} {
    \pgfmathparse{\r*1.6+0.8} \let\y\pgfmathresult
    \node[gamlabel, anchor=east, rotate=-45] at (-0.2, \y) {\rgam};
}
\node[font=\small\bfseries] at (3.2, 7.2) {Gamètes mâles (colonnes)};
\node[font=\small\bfseries, rotate=90] at (-1.0, 3.2) {Gamètes femelles (lignes)};

% Legend
\begin{scope}[yshift=-1.8cm]
\node[legbox, fill=popgoldL] (l1) at (0,0) {};
\node[right=2pt] at (l1.east) {Jaune lisse ($J\_A\_$) -- 9/16};
\node[legbox, fill=poporangeL] (l2) at (5,0) {};
\node[right=2pt] at (l2.east) {Jaune ridée ($J\_aa$) -- 3/16};
\node[legbox, fill=popgreenL] (l3) at (0,-0.6) {};
\node[right=2pt] at (l3.east) {Verte lisse ($jjA\_$) -- 3/16};
\node[legbox, fill=poppinkL] (l4) at (5,-0.6) {};
\node[right=2pt] at (l4.east) {Verte ridée ($jjaa$) -- 1/16};
\end{scope}
\end{tikzpicture}
\legende{Échiquier de croisement (carré de Punnett) 4 $\times$ 4 du dihybridisme $JjAa \times JjAa$. Les 16 cases équiprobables se regroupent en 4 phénotypes selon le ratio 9 : 3 : 3 : 1. Couleurs : \textcolor{popgoldL}{\rule{4pt}{4pt}} Jaune lisse ($J\_A\_$), \textcolor{poporangeL}{\rule{4pt}{4pt}} Jaune ridée ($J\_aa$), \textcolor{popgreenL}{\rule{4pt}{4pt}} Verte lisse ($jjA\_$), \textcolor{poppinkL}{\rule{4pt}{4pt}} Verte ridée ($jjaa$). Les génotypes des 16 cases sont indiqués.}
\end{popfigure}

\emph{Interprétation génétique.} Le ratio $9:3:3:1$ s'obtient par le produit des deux monohybridismes indépendants :
\[
(Jj \times Jj) \times (Aa \times Aa) = (\tfrac{3}{4}J\_ + \tfrac{1}{4}jj) \times (\tfrac{3}{4}A\_ + \tfrac{1}{4}aa)
\]
\[
= \tfrac{9}{16}J\_A\_ \;+\; \tfrac{3}{16}J\_aa \;+\; \tfrac{3}{16}jjA\_ \;+\; \tfrac{1}{16}jjaa
\]

\begin{propriete}[Loi de l'indépendance des caractères (2\textsuperscript{e} loi de Mendel) — Admise au Bac]
Lorsque deux gènes sont situés sur \textbf{chromosomes différents} (non homologues) \textbf{ou très éloignés sur le même chromosome}, les allèles de ces gènes se répartissent \textbf{indépendamment} les uns des autres lors de la méiose.
\begin{itemize}
    \item Un double hétérozygote $AaBb$ produit \textbf{4 types de gamètes équiprobables} : $AB$, $Ab$, $aB$, $ab$ (fréquence 1/4 chacun).
    \item L'autofécondation ($AaBb \times AaBb$) donne en F2 un ratio phénotypique \textbf{9 : 3 : 3 : 1} (si dominance complète pour les deux gènes).
    \item Le test-cross ($AaBb \times aabb$) donne un ratio \textbf{1 : 1 : 1 : 1}.
\end{itemize}
Cette propriété est la traduction génétique du \textbf{brassage interchromosomique}.
\end{propriete}

\courssub{Le brassage interchromosomique : base cytologique de l'indépendance des gènes}

Comment l'indépendance des gènes s'explique-t-elle au niveau cellulaire ? Rappelons que les gènes sont portés par les chromosomes. Si le gène $J/j$ est sur le chromosome 1 et le gène $A/a$ sur le chromosome 2, alors la paire de chromosomes 1 et la paire de chromosomes 2 se comportent indépendamment lors de la méiose.

\begin{popfigure}
\begin{tikzpicture}[scale=0.75, every node/.style={font=\small}]
% Styles
\tikzset{
    chrom/.style={draw=popdark, thick, minimum width=0.5cm, minimum height=1.3cm, rounded corners=2pt},
    chromM/.style={chrom, fill=popblueL}, % Maternal (Blue tint)
    chromP/.style={chrom, fill=poporangeL}, % Paternal (Orange tint)
    allele/.style={font=\small\bfseries, text=popdark},
    centromere/.style={fill=popdark, minimum width=0.5cm, minimum height=0.15cm},
    arrow/.style={->, >=Stealth, thick, popdark},
    labelbox/.style={draw=popdark, rounded corners, fill=popcream, inner sep=3pt, font=\small\bfseries, text=popdark},
    plate/.style={dashed, popmuted, thick}
}

% --- Configuration 1 : Orientation 1 ---
\node[labelbox] at (0, 7.0) {\textbf{Orientation 1 (aléatoire)}};
% Metaphase Plate
\draw[plate] (-3.5, 5.0) -- (3.5, 5.0) node[right=2mm] {Plaque équatoriale};

% Bivalent 1 (Chr 1 pair) - Left side
% Maternal (Top) J, Paternal (Bottom) j
\node[chromM] (b1m) at (-1.5, 5.8) {}; \node[allele, popblue] at (-1.5, 5.8) {$J$};
\node[chromP] (b1p) at (-1.5, 4.2) {}; \node[allele, poporange] at (-1.5, 4.2) {$j$};
\draw[popblue, thick] (-1.5, 5.45) -- (-1.5, 4.55); % Centromere connection

% Bivalent 2 (Chr 2 pair) - Right side
% Maternal (Top) A, Paternal (Bottom) a
\node[chromM] (b2m) at (1.5, 5.8) {}; \node[allele, popblue] at (1.5, 5.8) {$A$};
\node[chromP] (b2p) at (1.5, 4.2) {}; \node[allele, poporange] at (1.5, 4.2) {$a$};
\draw[poporange, thick] (1.5, 5.45) -- (1.5, 4.55);

% Anaphase I Arrows (Separation of homologs)
\draw[arrow] (-1.5, 5.8) -- (-1.5, 6.8); % J up
\draw[arrow] (1.5, 5.8) -- (1.5, 6.8);   % A up
\draw[arrow] (-1.5, 4.2) -- (-1.5, 3.2); % j down
\draw[arrow] (1.5, 4.2) -- (1.5, 3.2);   % a down

% Telophase I / Gametes Result 1 (Top pole gets J and A; Bottom gets j and a)
\node[labelbox] at (0, 1.8) {\textbf{Pôle 1 : Gamètes $JA$ \quad Pôle 2 : Gamètes $ja$}};
% Pole 1 (Top)
\node[chromM] at (-1.5, 1.8) {}; \node[allele, popblue] at (-1.5, 1.8) {$J$};
\node[chromM] at (0.5, 1.8) {}; \node[allele, popblue] at (0.5, 1.8) {$A$};
% Pole 2 (Bottom)
\node[chromP] at (-1.5, 0.2) {}; \node[allele, poporange] at (-1.5, 0.2) {$j$};
\node[chromP] at (0.5, 0.2) {}; \node[allele, poporange] at (0.5, 0.2) {$a$};

% --- Configuration 2 : Orientation 2 (shifted right) ---
\begin{scope}[xshift=10cm]
\node[labelbox] at (0, 7.0) {\textbf{Orientation 2 (aléatoire)}};
\draw[plate] (-3.5, 5.0) -- (3.5, 5.0) node[right=2mm] {Plaque équatoriale};

% Bivalent 1 (Chr 1) - Maternal on Bottom, Paternal on Top
\node[chromP] (b1p) at (-1.5, 5.8) {}; \node[allele, poporange] at (-1.5, 5.8) {$j$};
\node[chromM] (b1m) at (-1.5, 4.2) {}; \node[allele, popblue] at (-1.5, 4.2) {$J$};
\draw[popblue, thick] (-1.5, 5.45) -- (-1.5, 4.55);

% Bivalent 2 (Chr 2) - Maternal on Top, Paternal on Bottom (Same as Orient 1 for Chr2)
\node[chromM] (b2m) at (1.5, 5.8) {}; \node[allele, popblue] at (1.5, 5.8) {$A$};
\node[chromP] (b2p) at (1.5, 4.2) {}; \node[allele, poporange] at (1.5, 4.2) {$a$};
\draw[poporange, thick] (1.5, 5.45) -- (1.5, 4.55);

% Anaphase I Arrows
\draw[arrow] (-1.5, 5.8) -- (-1.5, 6.8); % j up
\draw[arrow] (1.5, 5.8) -- (1.5, 6.8);   % A up
\draw[arrow] (-1.5, 4.2) -- (-1.5, 3.2); % J down
\draw[arrow] (1.5, 4.2) -- (1.5, 3.2);   % a down

% Gametes Result 2
\node[labelbox] at (0, 1.8) {\textbf{Pôle 1 : Gamètes $jA$ \quad Pôle 2 : Gamètes $Ja$}};
% Pole 1 (Top) gets j and A
\node[chromP] at (-1.5, 1.8) {}; \node[allele, poporange] at (-1.5, 1.8) {$j$};
\node[chromM] at (0.5, 1.8) {}; \node[allele, popblue] at (0.5, 1.8) {$A$};
% Pole 2 (Bottom) gets J and a
\node[chromM] at (-1.5, 0.2) {}; \node[allele, popblue] at (-1.5, 0.2) {$J$};
\node[chromP] at (0.5, 0.2) {}; \node[allele, poporange] at (0.5, 0.2) {$a$};
\end{scope}
\end{tikzpicture}
\legende{Mécanisme du brassage interchromosomique (Anaphase I). Deux paires de chromosomes homologues (Bleu = Chromosome maternel, Orange = Chromosome paternel). Chr 1 porte $J/j$, Chr 2 porte $A/a$. \textbf{Gauche :} Orientation 1 $\rightarrow$ Gamètes parentaux ($JA$, $ja$). \textbf{Droite :} Orientation 2 $\rightarrow$ Gamètes recombinés ($Ja$, $jA$). Les deux orientations sont équiprobables (50\% chacune), d'où 4 types de gamètes à 25\% chacun.}
\end{popfigure}

\emph{Mécanisme cytologique.} En \textbf{métaphase I}, les bivalents (paires de chromosomes homologues dupliqués) s'alignent sur la plaque équatoriale. L'orientation de chaque bivalent (quel pôle recevra le chromosome maternel et quel pôle recevra le paternel) est \textbf{aléatoire et indépendante} de l'orientation des autres bivalents. Pour $n$ paires de chromosomes, il y a $2^n$ orientations possibles. Chez l'Homme ($n=23$), cela fait $2^{23}$ combinaisons chromosomiques différentes rien que par ce mécanisme.

\courssub{Le test-cross (croisement test) : vérifier l'indépendance et déterminer un génotype}

Le test-cross consiste à croiser un individu de phénotype dominant (génotype inconnu : $J\_A\_$) avec un individu \textbf{double récessif homozygote} ($jjaa$). Ce dernier ne produit qu'un seul type de gamète : $ja$. Le phénotype de la descendance reflète \textbf{directement} le génotype des gamètes produits par l'individu testé.

\begin{exemplebox}[Test-cross d'un double hétérozygote F1 ($JjAa \times jjaa$)]
\textbf{Gamètes F1 :} $JA$, $Ja$, $jA$, $ja$ (chacun 1/4).
\textbf{Gamètes double récessif :} $ja$ (100\%).
\textbf{Descendance attendue :}
\begin{center}
\begin{tabular}{|c|c|c|c|}\hline
Gamète F1 & Génotype descendance & Phénotype & Fréquence \\ \hline
$JA$ & $JjAa$ & Jaune lisse & 1/4 \\ \hline
$Ja$ & $Jjaa$ & Jaune ridée & 1/4 \\ \hline
$jA$ & $jjAa$ & Verte lisse & 1/4 \\ \hline
$ja$ & $jjaa$ & Verte ridée & 1/4 \\ \hline
\end{tabular}
\end{center}
\textbf{Ratio phénotypique : 1 : 1 : 1 : 1.} C'est la signature du brassage interchromosomique (gènes indépendants).
\end{exemplebox}

\begin{exemplebox}[Test-cross pour déterminer le génotype d'un individu « Jaune lisse »]
Un agriculteur camerounais a une plante de pois à graines jaunes lisses (phénotype dominant). Il ignore si elle est homozygote ($JJAA$), hétérozygote pour un gène ($JjAA$ ou $JJAA$), ou double hétérozygote ($JjAa$). Il la croise avec une plante verte ridée ($jjaa$).
\begin{itemize}
    \item \textbf{1 seule classe phénotypique} (100\% Jaune lisse) $\rightarrow$ Génotype **$JJAA$** (homozygote pour les deux).
    \item \textbf{2 classes phénotypiques 1:1} (ex: Jaune lisse / Jaune ridée) $\rightarrow$ Hétérozygote pour \textbf{un seul} gène (ex: $JjAA$).
    \item \textbf{4 classes phénotypiques 1:1:1:1} $\rightarrow$ **Double hétérozygote $JjAa$** (gènes indépendants).
\end{itemize}
\end{exemplebox}

\begin{methode}[Reconnaître un dihybridisme à gènes indépendants]
\emph{Données :} Effectifs phénotypiques d'une F2 (autofécondation F1) ou d'un test-cross (F1 $\times$ double récessif).
\begin{enumerate}
    \item \textbf{Calculer les fréquences observées} (effectif / total).
    \item \textbf{Comparer aux ratios théoriques} :
    \begin{itemize}
        \item F2 : $9/16$, $3/16$, $3/16$, $1/16$ (somme = 1).
        \item Test-cross : $1/4$, $1/4$, $1/4$, $1/4$.
    \end{itemize}
    \item \textbf{Test du $\chi^2$ (si effectifs fournis) :} Vérifier si les écarts sont dus au hasard ($p > 0,05$).
    \item \textbf{Conclure :} Si les proportions correspondent $\rightarrow$ \textbf{Gènes indépendants (brassage interchromosomique)}. Si écart significatif (excès de parentaux, déficit de recombinés) $\rightarrow$ \textbf{Gènes liés (brassage intrachromosomique, voir Section 3)}.
\end{enumerate}
\end{methode}

\courssub{Généralisation : nombre de combinaisons gamétiques possibles}

Pour un individu hétérozygote pour $n$ couples d'allèles \textbf{indépendants} (portés par $n$ paires de chromosomes différentes ou très éloignés), le nombre de types de gamètes différents produits est :
\[
N_{\text{gamètes}} = 2^n
\]
Chaque gamète reçoit un allèle par couple, choisi au hasard parmi les deux (maternel ou paternel).

\begin{savaistu}[La puissance du brassage interchromosomique chez l'Homme]
Chez l'Homme, $n = 23$ paires de chromosomes. Le brassage interchromosomique seul permet de produire :
\[
2^{23} = 8\,388\,608 \approx \mathbf{8,4 \times 10^6} \text{ combinaisons chromosomiques différentes par méiose.}
\]
C'est-à-dire qu'un seul individu peut théoriquement produire plus de \textbf{8 millions de gamètes génétiquement distincts} rien qu'en mélangeant ses chromosomes maternels et paternels ! Ajoutez à cela le crossing-over (brassage intrachromosomique) et la fécondation aléatoire, et l'unicité génétique de chaque être humain devient une évidence mathématique.
\end{savaistu}

\begin{attention}[Pièges fréquents et erreurs à éviter]
\begin{itemize}
    \item \textbf{Confondre brassage inter- et intra-chromosomique :} L'indépendance (9:3:3:1) ne vaut que pour des gènes sur \textbf{chromosomes différents} (ou très loin sur le même chromosome). Sur le même chromosome et proches, c'est le linkage (voir Section 3).
    \item \textbf{Oublier la condition « effectifs suffisants » :} Le ratio 9:3:3:1 est une \textbf{espérance mathématique}. Sur de petits effectifs (ex: 16 graines), le hasard des rencontres gamétiques crée de forts écarts. Au Bac, on vous demandera souvent de faire un \textbf{test du $\chi^2$} pour valider l'hypothèse d'indépendance.
    \item \textbf{Mauvaise lecture de l'échiquier :} Ne comptez pas les \emph{génotypes} (16 cases) mais regroupez-les en \emph{phénotypes} (4 classes). Vérifiez la dominance : $J\_$ = Jaune, $jj$ = Vert ; $A\_$ = Lisse, $aa$ = Ridée.
    \item \textbf{Test-cross :} Ne pas confondre le génotype du parent testé (inconnu) et celui du parent testeur (toujours double récessif homozygote $aabb$). Le parent testeur ne sert qu'à « révéler » les gamètes de l'autre.
    \item \textbf{Notation des gamètes :} Écrire les allèles dans l'ordre des gènes (ex: $JA$, pas $AJ$) pour éviter la confusion lors du croisement.
\end{itemize}
\end{attention}

\courssub{Exercice résolu : Analyse d'un dihybridisme chez le maïs}

\begin{exoresolu}[Dihybridisme et test-cross chez le maïs]
\textbf{Énoncé.} Chez le maïs, la couleur du grain est déterminée par un gène ($C/c$) : coloré ($C$, dominant) / incolore ($c$, récessif). La forme du grain est déterminée par un second gène ($S/s$) : lisse ($S$, dominant) / ridé ($s$, récessif). On croise une lignée pure à grains colorés lisses avec une lignée pure à grains incolores ridés. On obtient une F1 uniforme. On réalise ensuite un test-cross : F1 $\times$ lignée incolore ridée. On observe 400 grains dans la descendance :
\begin{itemize}
    \item 98 grains colorés lisses
    \item 102 grains colorés ridés
    \item 101 grains incolores lisses
    \item 99 grains incolores ridés
\end{itemize}
\textbf{Questions :}
\begin{enumerate}
    \item Déterminer les génotypes des parents P, de la F1 et du parent testeur.
    \item Préciser les types de gamètes et leurs fréquences attendues pour la F1 si les gènes sont indépendants.
    \item Vérifier par un test du $\chi^2$ (seuil 5\%) si les résultats expérimentaux sont compatibles avec l'hypothèse d'indépendance des gènes.
    \item Conclure sur la localisation des gènes $C$ et $S$.
\end{enumerate}

\textbf{Solution détaillée.}
\begin{enumerate}
    \item \textbf{Génotypes :}
    \begin{itemize}
        \item P1 (coloré lisse pur) : $CCSS$. Gamètes : $CS$.
        \item P2 (incolore ridé pur) : $ccss$. Gamètes : $cs$.
        \item F1 (tous colorés lisses) : $CcSs$. Gamètes attendus si indépendance : $CS$, $Cs$, $cS$, $cs$ (1/4 chacun).
        \item Parent testeur (incolore ridé) : $ccss$. Gamètes : $cs$ uniquement.
    \end{itemize}
    \item \textbf{Descendance théorique du test-cross (indépendance) :}
    \begin{center}
    \begin{tabular}{|c|c|c|c|}\hline
    Gamète F1 & Génotype descendance & Phénotype & Effectif théorique ($N=400$) \\ \hline
    $CS$ (1/4) & $CcSs$ & Coloré lisse & 100 \\ \hline
    $Cs$ (1/4) & $Ccss$ & Coloré ridé & 100 \\ \hline
    $cS$ (1/4) & $ccSs$ & Incolore lisse & 100 \\ \hline
    $cs$ (1/4) & $ccss$ & Incolore ridé & 100 \\ \hline
    \end{tabular}
    \end{center}
    Ratio attendu : \textbf{1 : 1 : 1 : 1}.
    \item \textbf{Test du $\chi^2$ (indépendance) :}
    \begin{align*}
    \chi^2_{\text{calc}} &= \sum \frac{(O_i - E_i)^2}{E_i} \\
    &= \frac{(98-100)^2}{100} + \frac{(102-100)^2}{100} + \frac{(101-100)^2}{100} + \frac{(99-100)^2}{100} \\
    &= \frac{4}{100} + \frac{4}{100} + \frac{1}{100} + \frac{1}{100} = \frac{10}{100} = \mathbf{0,10}
    \end{align*}
    \textbf{Degrés de liberté :} $dl = \text{nombre de classes} - 1 = 4 - 1 = 3$.
    \textbf{Valeur critique ($\alpha = 5\%$, $dl=3$) :} $\chi^2_{\text{lim}} = 7,815$ (valeur à connaître ou fournie).
    \textbf{Comparaison :} $\chi^2_{\text{calc}} = 0,10 < 7,815 = \chi^2_{\text{lim}}$.
    \item \textbf{Conclusion :} On \textbf{ne rejette pas} l'hypothèse nulle ($H_0$ : les gènes sont indépendants). Les écarts observés sont dus aux fluctuations d'échantillonnage (hasard). Les gènes $C$ et $S$ sont \textbf{non liés} (portés par des chromosomes différents ou très éloignés). Le brassage interchromosomique s'applique.
\end{enumerate}
\end{exoresolu}

\begin{aretenir}[Synthèse : Dihybridisme et Brassage Interchromosomique]
\begin{itemize}
    \item \textbf{Dihybridisme :} Étude simultanée de 2 caractères (4 allèles, 2 gènes).
    \item \textbf{Condition :} Gènes sur \textbf{chromosomes différents} (non homologues) \textbf{ou très éloignés sur le même chromosome}.
    \item \textbf{Mécanisme :} \textbf{Brassage interchromosomique} = orientation aléatoire et indépendante des bivalents en Anaphase I de la méiose.
    \item \textbf{Conséquence gamétique :} Double hétérozygote $\rightarrow$ \textbf{4 gamètes équiprobables} (25\% chacun : 2 parentaux + 2 recombinés).
    \item \textbf{Ratios caractéristiques :}
    \begin{itemize}
        \item F2 (autofécondation F1) : \textbf{9 : 3 : 3 : 1} (phénotypes).
        \item Test-cross (F1 $\times$ double récessif) : \textbf{1 : 1 : 1 : 1} (phénotypes = gamètes F1).
    \end{itemize}
    \item \textbf{Généralisation :} $n$ gènes indépendants $\rightarrow 2^n$ types de gamètes.
    \item \textbf{Outil de validation :} Test du $\chi^2$ sur effectifs observés vs théoriques.
\end{itemize}
\end{aretenir}

\coursec{Brassage intrachromosomique : linkage et crossing-over}

\courssub{Transition depuis le brassage interchromosomique}
Dans la section précédente, nous avons vu que lorsque deux gènes sont portés par des chromosomes \emph{différents} (ou très éloignés sur le même chromosome), ils se séparent indépendamment lors de la méiose : c’est le \textbf{brassage interchromosomique}. Les proportions phénotypiques observées dans un test‑cross dihybride sont alors de $1:1:1:1$ (quatre classes égales).  
Or, de nombreux croisements (notamment ceux de \textbf{Thomas Hunt Morgan} sur la drosophile) donnent des proportions \emph{très différentes} : deux classes sont majoritaires, deux sont minoritaires. Cela révèle que les deux gènes sont \textbf{liés} sur le \emph{même} chromosome. Nous étudions ici ce \textbf{brassage intrachromosomique}.

\courssub{Expériences de Morgan sur la drosophile}
\begin{experience}[Croisements de Morgan (1910‑1915)]
Morgan croise des drosophiles \emph{Drosophila melanogaster} portant deux caractères : la couleur du corps (gris $G$ dominant sur noir $g$) et la forme des ailes (normales $N$ dominantes sur vestigiales $n$).  
Il réalise un \textbf{test‑cross} : un double hétérozygote $Gg\,Nn$ (phénotype gris‑ailes normales), de phase connue \emph{cis} ($\frac{G\,N}{g\,n}$), est croisé avec un double récessif homozygote $gg\,nn$ (noir‑ailes vestigiales, $\frac{g\,n}{g\,n}$).  
\emph{Observation} : parmi la descendance, on compte environ $42\,\%$ gris‑ailes normales, $42\,\%$ noir‑ailes vestigiales, $8\,\%$ gris‑ailes vestigiales et $8\,\%$ noir‑ailes normales.  
\emph{Interprétation} : les deux gènes sont portés par le même chromosome. Les combinaisons parentales (gris‑normales et noir‑vestigiales) sont majoritaires ; les combinaisons recombinées (gris‑vestigiales et noir‑normales) sont minoritaires. Cela prouve l’existence d’un \textbf{linkage partiel} avec \textbf{crossing‑over}.
\end{experience}

\begin{savaistu}[Morgan et le prix Nobel]
Thomas Hunt Morgan reçoit le \textbf{prix Nobel de physiologie ou médecine en 1933} pour ses travaux sur le rôle des chromosomes dans l’hérédité, établissant la \emph{théorie chromosomique de l’hérédité} à partir de la drosophile.
\end{savaistu}

\courssub{Linkage total}
Lorsque deux gènes sont \emph{très proches} sur le même chromosome, aucun crossing‑over ne se produit entre eux (ou il est extrêmement rare).  
\begin{definition}[Linkage total]
Deux gènes sont en \textbf{linkage total} s’ils sont portés par le même chromosome et qu’aucun échange de segments (crossing‑over) n’a lieu entre eux au cours de la méiose.  
\emph{Conséquence} : un test‑cross ne produit que \textbf{deux types de descendants} (les deux combinaisons parentales) dans des proportions $50\,\%$ – $50\,\%$.
\end{definition}

\begin{exemplebox}[Exemple de linkage total]
Chez la drosophile, les gènes \emph{sepia} (yeux bruns, allèle $se$) et \emph{ebony} (corps noir, allèle $e$) sont si proches qu’aucun recombinant n’est observé. Un test‑cross $\frac{se^+\,e^+}{se\,e} \times \frac{se\,e}{se\,e}$ donne uniquement $se^+\,e^+$ (phénotype sauvage) et $se\,e$ (double mutant) à parts égales.
\end{exemplebox}

\courssub{Linkage partiel et crossing-over}
Lorsque la distance entre deux gènes sur un chromosome est suffisante, des \textbf{crossing‑overs} peuvent se produire en prophase I de la méiose.

\begin{definition}[Crossing‑over]
Le \textbf{crossing‑over} est l’échange \emph{réciproque} de fragments d’ADN entre deux \textbf{chromatides non‑sœurs} de chromosomes \textbf{homologues} au cours de la \textbf{prophase I} de la méiose (stade de \emph{pachytène}). Il se manifeste cytologiquement par la formation de \textbf{chiasmas}.
\end{definition}

\begin{popfigure}
\begin{tikzpicture}[scale=0.85, every node/.style={font=\small}]
% Chromosome homologue 1 (bleu) - deux chromatides sœurs
\draw[popblue, thick] (0,2.2) -- (7,2.2);
\draw[popblue, thick] (0,1.8) -- (7,1.8);
% Chromosome homologue 2 (orange) - deux chromatides sœurs
\draw[poporange, thick] (0,0.2) -- (7,0.2);
\draw[poporange, thick] (0,-0.2) -- (7,-0.2);
% Centromères
\draw[popdark, thick] (3.5,2.2) -- (3.5,-0.2);
% Allèles avant crossing-over
\node[popblue] at (1.5,2.5) {$A$};
\node[popblue] at (5.5,2.5) {$B$};
\node[popblue] at (1.5,1.5) {$A$};
\node[popblue] at (5.5,1.5) {$B$};
\node[poporange] at (1.5,-0.5) {$a$};
\node[poporange] at (5.5,-0.5) {$b$};
\node[poporange] at (1.5,-0.9) {$a$};
\node[poporange] at (5.5,-0.9) {$b$};
% Chiasma entre chromatides non-sœurs (haut du 1, bas du 2)
\draw[popgreen, thick, decorate, decoration={snake, amplitude=1.5pt, segment length=4pt}] (3.5,1.8) -- (3.5,0.2);
\node[popgreen, right] at (3.7,1.0) {Chiasma};
% Légendes
\node[align=left, popdark] at (-1.5,2.0) {Chromatides\\sœurs (identiques)};
\node[align=left, popdark] at (-1.5,0.0) {Chromatides\\sœurs (identiques)};
\node[align=center, popblue] at (3.5,2.7) {Chromosome homologue 1 (paternel)};
\node[align=center, poporange] at (3.5,-1.3) {Chromosome homologue 2 (maternel)};
\end{tikzpicture}
\legende{Crossing-over en prophase I (stade pachytène). Échange réciproque entre deux chromatides non-sœurs (une de chaque homologue). Les allèles parentaux $A\,B$ et $a\,b$ donnent naissance à des chromatides recombinées $A\,b$ et $a\,B$.}
\end{popfigure}

\begin{popfigure}
\begin{tikzpicture}[scale=0.85, every node/.style={font=\small}]
% Après crossing-over : 4 chromatides distinctes
% Chromatide 1 (parentale 1)
\draw[popblue, thick] (0,2.5) -- (7,2.5);
\node[popblue] at (1.5,2.8) {$A$};
\node[popblue] at (5.5,2.8) {$B$};
\node[align=left] at (-1.5,2.5) {Parentale 1\\$A\,B$};
% Chromatide 2 (recombinée 1)
\draw[popblue, thick] (0,1.5) -- (3.5,1.5);
\draw[poporange, thick] (3.5,1.5) -- (7,1.5);
\node[popblue] at (1.5,1.8) {$A$};
\node[poporange] at (5.5,1.2) {$b$};
\node[align=left] at (-1.5,1.5) {Recombinée 1\\$A\,b$};
% Chromatide 3 (recombinée 2)
\draw[poporange, thick] (0,0.5) -- (3.5,0.5);
\draw[popblue, thick] (3.5,0.5) -- (7,0.5);
\node[poporange] at (1.5,0.8) {$a$};
\node[popblue] at (5.5,0.8) {$B$};
\node[align=left] at (-1.5,0.5) {Recombinée 2\\$a\,B$};
% Chromatide 4 (parentale 2)
\draw[poporange, thick] (0,-0.5) -- (7,-0.5);
\node[poporange] at (1.5,-0.2) {$a$};
\node[poporange] at (5.5,-0.2) {$b$};
\node[align=left] at (-1.5,-0.5) {Parentale 2\\$a\,b$};
% Centromères
\draw[popdark, thick] (3.5,2.5) -- (3.5,-0.5);
\end{tikzpicture}
\legende{Résultat du crossing-over : sur les quatre chromatides, deux sont parentales ($A\,B$ et $a\,b$) et deux sont recombinées ($A\,b$ et $a\,B$). Lors de l'anaphase I, les chromosomes homologues se séparent, distribuant ces chromatides dans les gamètes.}
\end{popfigure}

\begin{propriete}[Linkage partiel]
Si le crossing‑over a lieu entre deux gènes liés, le test‑cross produit \textbf{quatre classes de descendants} :
\begin{itemize}
\item deux \textbf{types parentaux} (combinaisons d’allèles identiques aux chromosomes parentaux) — majoritaires ;
\item deux \textbf{types recombinés} (nouvelles combinaisons issues du crossing‑over) — minoritaires.
\end{itemize}
Les fréquences des types parentaux sont égales entre elles, de même pour les types recombinés.
\end{propriete}

\begin{popfigure}
\begin{tikzpicture}
\begin{axis}[
    ybar,
    bar width=12pt,
    width=\linewidth,
    height=7cm,
    ylabel={Fréquence (\% )},
    xtick={1,2,3,4},
    xticklabels={Parentaux 1, Parentaux 2, Recombinés 1, Recombinés 2},
    ymin=0, ymax=50,
    ymajorgrids=true,
    grid style=dashed,
    legend style={at={(0.5,-0.15)},anchor=north,legend columns=-1},
    nodes near coords,
    nodes near coords align={vertical},
]
\addplot[popblue!70, fill=popblue!70] coordinates {(1,42) (2,42) (3,8) (4,8)};
\legend{Test‑cross Morgan (exemple)}
\end{axis}
\end{tikzpicture}
\legende{Diagramme en barres des quatre classes phénotypiques observées dans le test‑cross de Morgan (linkage partiel). Les deux classes parentales (42\,\% chacune) sont majoritaires ; les deux classes recombinées (8\,\% chacune) sont minoritaires.}
\end{popfigure}

\courssub{Mécanisme du crossing-over}
\begin{methode}[Étapes du crossing-over en prophase I]
\begin{enumerate}
\item \textbf{Appariement (synapse)} : les chromosomes homologues se rapprochent et s’alignent gène par gène, formant un \emph{bivalent} (tétra­de de quatre chromatides).
\item \textbf{Formation du complexe synaptonémal} : une structure protéique maintient les homologues ensemble.
\item \textbf{Cassures double brin} : des enzymes (Spo11) créent des cassures dans l’ADN des chromatides non‑sœurs.
\item \textbf{Réparation par échange} : les brins cassés envahissent la chromatide non‑sœur homologue, entraînant un échange réciproque de segments.
\item \textbf{Chiasmas} : les points d’échange deviennent visibles au microscope comme des \emph{chiasmas} (au moins un par bras chromosomique).
\item \textbf{Séparation} : en anaphase I, les chromosomes homologues (chacun composé de deux chromatides sœurs, dont une éventuellement recombinée) migrent vers des pôles opposés.
\end{enumerate}
\end{methode}

\courssub{Taux de recombinaison et carte factorielle}
\begin{definition}[Taux de recombinaison]
Pour un couple de gènes liés, le \textbf{taux de recombinaison} ($T_R$) est le pourcentage d’individus recombinés parmi la descendance totale d’un test‑cross :
\[
T_R = \frac{N_{\text{recombinés}}}{N_{\text{total}}} \times 100 \quad (\text{en \%}).
\]
\end{definition}

\begin{propriete}[Distance génétique et centimorgan]
La \textbf{distance génétique} entre deux gènes s’exprime en \textbf{centimorgans (cM)}. Par définition, pour de faibles distances :
\[
1\,\% \text{ de recombinaison } = 1 \text{ cM}.
\]
Ainsi, $T_R = 16,2\,\% \;\Rightarrow\; d = 16,2 \text{ cM}$.
\end{propriete}

\begin{attention}[Limite de la relation $d = T_R$]
Cette relation linéaire n'est valable que pour des distances inférieures à environ $10-15$ cM. Au-delà, les doubles crossing-overs (non détectés phénotypiquement) font sous-estimer la distance réelle. Au Bac, on utilise l'additivité des distances pour ordonner les gènes, en supposant l'absence d'interférence significative sur les intervalles étudiés.
\end{attention}

\begin{methode}[Construire une carte factorielle à partir d’un test-cross]
\begin{enumerate}
\item Effectuer un test‑cross entre un double hétérozygote (phase \emph{cis} ou \emph{trans} connue) et un double récessif.
\item Compter les effectifs des quatre classes phénotypiques.
\item Identifier les deux classes \textbf{parentales} (les plus nombreuses) et les deux classes \textbf{recombinées} (les moins nombreuses).
\item Calculer $T_R = \frac{n_{\text{rec1}}+n_{\text{rec2}}}{N_{\text{total}}}\times 100$.
\item La distance en cM entre les deux gènes est $d = T_R$.
\item Pour trois gènes ou plus : répéter les croisements deux à deux. Le gène du milieu est celui pour lequel la \textbf{somme de ses distances aux deux autres est la plus petite} (cette somme correspond alors à la distance entre les deux gènes extrêmes).
\end{enumerate}
\end{methode}

\begin{exoresolu}[Calcul du taux de recombinaison et carte factorielle à trois gènes]
\textbf{Énoncé.} Un test‑cross chez la drosophile (gènes \emph{vg} (ailes vestigiales) et \emph{se} (yeux sépia)) donne la descendance suivante :  
$vg^+ se^+$ : 420 ; $vg\,se$ : 418 ; $vg^+ se$ : 78 ; $vg\,se^+$ : 84.  
Calculer le taux de recombinaison, la distance en cM et proposer l’ordre des gènes si un troisième gène \emph{e} (ébony) donne $vg$–$e = 12$ cM et $se$–$e = 28$ cM.

\tcblower

\textbf{Solution.}
\begin{enumerate}
\item Total $N = 420+418+78+84 = 1000$.
\item Classes parentales (majoritaires) : $vg^+ se^+$ (420) et $vg\,se$ (418).  
  Classes recombinées : $vg^+ se$ (78) et $vg\,se^+$ (84).
\item $N_{\text{rec}} = 78+84 = 162$.
\item $T_R = \frac{162}{1000}\times 100 = 16,2\,\%$.
\item Distance $vg$–$se = 16,2 \text{ cM}$.
\item Recherche de l'ordre pour les trois gènes ($vg$, $se$, $e$) :  
  Distances connues : $d(vg, se) = 16,2$ cM ; $d(vg, e) = 12$ cM ; $d(se, e) = 28$ cM.  
  On teste les trois ordres possibles en vérifiant l'additivité (distance extrêmes = somme des deux intervalles) :
  \begin{itemize}
  \item Ordre \textbf{$vg$ – $se$ – $e$} : $d(vg,e)$ prévu $= 16,2 + 28 = 44,2$ cM $\neq 12$ cM (rejeté).
  \item Ordre \textbf{$se$ – $vg$ – $e$} : $d(se,e)$ prévu $= 16,2 + 12 = 28,2$ cM $\approx 28$ cM (\textbf{validé}).
  \item Ordre \textbf{$vg$ – $e$ – $se$} : $d(vg,se)$ prévu $= 12 + 28 = 40$ cM $\neq 16,2$ cM (rejeté).
  \end{itemize}
  Le gène du milieu est \textbf{$vg$} (somme des distances aux autres $= 16,2 + 12 = 28,2$ cM, la plus petite des trois sommes).  
  Ordre final : \textbf{$se$ — $vg$ — $e$} avec distances $se$–$vg = 16,2$ cM, $vg$–$e = 12$ cM, $se$–$e = 28,2$ cM.
\end{enumerate}
\end{exoresolu}

\begin{popfigure}
\begin{tikzpicture}[scale=0.9, every node/.style={font=\small}]
% Ligne chromosomique
\draw[popdark, thick] (0,0) -- (10,0);
% Gènes dans l'ordre se - vg - e
\node[poporange, circle, draw, minimum width=8pt, inner sep=1pt] (se) at (0,0) {};
\node[popblue, circle, draw, minimum width=8pt, inner sep=1pt] (vg) at (5.8,0) {}; % 16.2 * 0.36 approx
\node[popgreen, circle, draw, minimum width=8pt, inner sep=1pt] (e) at (10,0) {};  % 28.2 * 0.36 approx
% Étiquettes
\node[above=4pt] at (se) {$se$};
\node[above=4pt] at (vg) {$vg$};
\node[above=4pt] at (e) {$e$};
% Distances
\draw[<->, popmuted] (0.2,-0.6) -- (5.6,-0.6) node[midway, below=2pt] {16,2 cM};
\draw[<->, popmuted] (6,-0.6) -- (9.8,-0.6) node[midway, below=2pt] {12 cM};
\draw[<->, popmuted, dashed] (0.2,-1.2) -- (9.8,-1.2) node[midway, below=2pt] {28,2 cM};
\end{tikzpicture}
\legende{Carte factorielle linéaire à trois gènes ($se$, $vg$, $e$) avec distances en centimorgans. Le gène $vg$ est central car la somme de ses distances aux deux autres ($16,2 + 12 = 28,2$ cM) est égale à la distance entre les gènes extrêmes ($se$–$e$).}
\end{popfigure}

\courssub{Comparaison des deux brassages}
\begin{center}
\begin{tabularx}{\linewidth}{|l|X|X|}
\hline
\rowcolor{popblueL}
\textbf{Critère} & \textbf{Brassage interchromosomique} & \textbf{Brassage intrachromosomique} \\ \hline
Localisation des gènes & Chromosomes \emph{différents} (ou très éloignés) & Même chromosome (gènes \emph{liés}) \\ \hline
Ségrégation & Indépendante (loi de Mendel 2) & Dépendante (linkage) \\ \hline
Test‑cross dihybride & 4 classes \textbf{égales} ($1:1:1:1$) & 2 classes parentales majoritaires + 2 classes recombinées minoritaires \\ \hline
Taux de recombinaison & $50\,\%$ (maximum) & $< 50\,\%$ (strictement) \\ \hline
Base cytologique & Assortiment indépendant des bivalents en métaphase I & Crossing‑over entre chromatides non‑sœurs en prophase I \\ \hline
Unité de distance & — & centimorgan (cM) : $1\,\%$ recombinaison = 1 cM \\ \hline
\end{tabularx}
\end{center}

\begin{attention}[Pièges fréquents]
\begin{itemize}
\item \textbf{Confondre $50\,\%$ de recombinaison et linkage.} Un taux de $50\,\%$ signifie que les gènes se comportent \emph{comme s’ils étaient sur des chromosomes différents} (indépendance). Il n’y a \emph{pas} de linkage détectable.
\item \textbf{Oublier que le taux de recombinaison est toujours $< 50\,\%$ pour des gènes liés.} Si vous trouvez $> 50\,\%$, c’est une erreur de comptage ou d’identification des classes parentales.
\item \textbf{Ne pas vérifier la phase du double hétérozygote.} Le test‑cross ne donne des proportions interprétables que si l’on connaît quels allèles sont sur le même chromosome (phase \emph{cis} : $\frac{AB}{ab}$ ou \emph{trans} : $\frac{Ab}{aB}$).
\item \textbf{Inverser la règle du gène du milieu.} Le gène du milieu est celui dont la \textbf{somme des distances aux deux autres est la plus petite} (elle équivaut à la distance entre les gènes extrêmes).
\item \textbf{Additionner les distances sans vérifier l’ordre des gènes.} Les distances ne sont additives que si le gène du milieu est correctement identifié.
\end{itemize}
\end{attention}

\begin{aretenir}[Points clés du brassage intrachromosomique]
\begin{itemize}
\item Deux gènes sur le même chromosome sont \textbf{liés} (linkage).
\item \textbf{Linkage total} : pas de crossing-over $\rightarrow$ seulement 2 classes (50/50).
\item \textbf{Linkage partiel} : crossing-over possible $\rightarrow$ 4 classes inégales (parentaux $>$ recombinés).
\item Le \textbf{crossing-over} est un échange réciproque entre chromatides non‑sœurs en prophase I (chiasmas).
\item \textbf{Taux de recombinaison} $T_R = \frac{\text{recombinés}}{\text{total}}\times 100$ ; distance $d = T_R$ cM (pour $d < 15$ cM).
\item Une \textbf{carte factorielle} ordonne les gènes et donne leurs distances en cM ; le gène central a la plus petite somme de distances aux extrêmes.
\item $T_R = 50\,\%$ $\Leftrightarrow$ gènes indépendants (brassage interchromosomique).
\end{itemize}
\end{aretenir}

\coursec{Pedigrees et hérédité autosomique humaine}

Dans la section précédente, nous avons vu que le brassage intrachromosomique (linkage et crossing-over) permet d'expliquer la transmission conjointe de gènes portés par un même chromosome, et même d'établir des cartes factorielles. Ces analyses reposaient sur des croisements expérimentaux contrôlés, réalisables chez la drosophile ou le pois. Mais chez l'Homme, on ne peut évidemment pas organiser de croisements : le nombre de descendants par couple est faible, la durée d'une génération est longue, et l'éthique l'interdit. Comment alors étudier l'hérédité des caractères humains, en particulier celle des maladies génétiques qui frappent certaines familles ? La réponse est à la fois simple et puissante : on \emph{observe} les familles telles qu'elles sont, on recense les individus atteints et les individus sains sur plusieurs générations, et on raisonne sur cette documentation. C'est toute la démarche du \cle{pedigree}, outil fondamental de la génétique humaine et du conseil génétique pratiqué dans les hôpitaux, y compris au Cameroun.

\courssub{La notion de pedigree : lire l'histoire d'une famille}

\begin{definition}[Pedigree]
Un \cle{pedigree} (ou \cle{arbre généalogique}) est la représentation schématique et conventionnelle de la transmission d'un caractère (souvent une maladie héréditaire) au sein d'une famille, sur plusieurs générations. Il permet de visualiser quels individus sont atteints, quels individus sont sains, et comment le caractère se transmet des parents aux descendants.
\end{definition}

Comme tout langage scientifique, le pedigree obéit à des conventions graphiques strictes, qu'il faut connaître parfaitement pour le jour du Baccalauréat :

\begin{itemize}
\item un \textbf{carré} représente un homme, un \textbf{rond} (cercle) représente une femme ;
\item un symbole \textbf{plein} (noir) représente un individu \textbf{atteint} (qui présente le caractère étudié) ; un symbole \textbf{vide} (blanc) représente un individu \textbf{sain} ;
\item un \textbf{trait horizontal} reliant un homme et une femme représente une \textbf{union} (couple) ;
\item un \textbf{trait vertical} descendant de ce trait horizontal mène aux \textbf{enfants}, reliés entre eux par un trait horizontal (fratrie) ; les enfants sont classés de gauche à droite, de l'aîné au cadet ;
\item les \textbf{générations} sont numérotées en chiffres romains (I, II, III\ldots) de haut en bas ; au sein d'une génération, chaque individu est repéré par un numéro arabe (I-1, I-2, II-1, II-2\ldots) ;
\item un losange indique parfois un individu de sexe non précisé ; un symbole barré d'une diagonale désigne un individu décédé.
\end{itemize}

\begin{methode}[Analyser un pedigree : déterminer le mode de transmission d'une maladie]
Face à un arbre généalogique, procède toujours dans cet ordre :
\begin{enumerate}
\item \textbf{La maladie est-elle dominante ou récessive ?} Cherche un couple de parents \emph{sains} ayant au moins un enfant \emph{atteint}. Si un tel couple existe, les parents possèdent l'allèle de la maladie sans l'exprimer : l'allèle morbide est donc \textbf{récessif} (il est masqué par l'allèle sain). Au contraire, si tout enfant atteint a toujours au moins un parent atteint, et que la maladie apparaît à \emph{toutes les générations}, l'allèle est vraisemblablement \textbf{dominant}.
\item \textbf{Le gène est-il autosomique ou gonosomique (porté par un chromosome sexuel) ?} Observe la répartition entre les sexes. Si la maladie frappe indifféremment les garçons et les filles, et si un père peut transmettre la maladie à son fils, le gène est \textbf{autosomique}. Si seuls les garçons (ou presque) sont atteints, et qu'aucun père ne transmet la maladie à son fils, on soupçonne un gène porté par le chromosome X (voir section suivante).
\item \textbf{Vérifier la cohérence} en attribuant des génotypes à chaque individu : tous les génotypes déduits doivent être compatibles avec les phénotypes observés. C'est cette vérification logique qui valide définitivement l'hypothèse.
\end{enumerate}
\end{methode}

\begin{attention}[Un caractère qui saute des générations n'est pas forcément récessif]
Le « saut de génération » est un indice fort de récessivité, mais ce n'est pas une preuve absolue : une maladie dominante peut sembler sauter une génération si le parent porteur est peu ou pas atteint (pénétrance incomplète), ou si la descendance observée est trop petite pour que le hasard ait permis la transmission. Inversement, une maladie récessive peut apparaître à chaque génération si des porteurs sains épousent des individus atteints. \textbf{Ne conclus jamais sur un seul indice} : attribue les génotypes et vérifie que tout l'arbre est cohérent. C'est la logique des génotypes, cas par cas, qui fait foi.
\end{attention}

\courssub{L'hérédité autosomique récessive : l'exemple de l'albinisme}

\textbf{Observation.} Dans certaines familles camerounaises, deux parents à la peau normalement pigmentée donnent naissance à un enfant \cle{albinos}, dont la peau, les cheveux et les iris sont dépigmentés. Comment deux parents sains peuvent-ils avoir un enfant atteint ?

\textbf{Hypothèse.} L'allèle responsable de l'albinisme est récessif : les parents sont \emph{porteurs sains} (hétérozygotes).

\textbf{Exploitation du document.} Considérons le pedigree suivant, relevé sur trois générations.

\begin{popfigure}
\begin{center}
\begin{tikzpicture}[
  male/.style={draw=popblue, thick, minimum size=0.55cm, inner sep=0pt},
  female/.style={draw=poppink, thick, circle, minimum size=0.55cm, inner sep=0pt},
  affected/.style={fill=popdark},
  lbl/.style={font=\small, text=popmuted},
  gen/.style={font=\small\bfseries, text=poporange},
  union/.style={thick, popline},
  desc/.style={thick, popline}]
% Génération I
\node[male] (I1) at (0,0) {};
\node[female, affected] (I2) at (1.4,0) {};
\node[lbl, below=2pt of I1] {I-1};
\node[lbl, below=2pt of I2] {I-2};
\draw[union] (I1) -- (I2);
% Génération II
\coordinate (midI) at (0.7,0);
\draw[desc] (midI) -- (0.7,-1.1);
\draw[desc] (-1.4,-1.1) -- (2.8,-1.1);
\node[male] (II1) at (-1.4,-1.9) {};
\node[female] (II2) at (0,-1.9) {};
\node[male] (II3) at (1.4,-1.9) {};
\node[female, affected] (II4) at (2.8,-1.9) {};
\draw[desc] (-1.4,-1.1) -- (II1.north);
\draw[desc] (0,-1.1) -- (II2.north);
\draw[desc] (1.4,-1.1) -- (II3.north);
\draw[desc] (2.8,-1.1) -- (II4.north);
\node[lbl, below=2pt of II1] {II-1};
\node[lbl, below=2pt of II2] {II-2};
\node[lbl, below=2pt of II3] {II-3};
\node[lbl, below=2pt of II4] {II-4};
% Conjoint de II2
\node[male] (IIc) at (-2.4,-1.9) {};
\draw[union] (IIc) -- (II2);
\node[lbl, below=2pt of IIc] {II-c};
% Génération III (enfants de II2 x IIc)
\coordinate (midII) at (-1.2,-1.9);
\draw[desc] (midII) -- (-1.2,-2.9);
\draw[desc] (-2.1,-2.9) -- (-0.3,-2.9);
\node[female] (III1) at (-2.1,-3.7) {};
\node[male, affected] (III2) at (-1.2,-3.7) {};
\node[female] (III3) at (-0.3,-3.7) {};
\draw[desc] (-2.1,-2.9) -- (III1.north);
\draw[desc] (-1.2,-2.9) -- (III2.north);
\draw[desc] (-0.3,-2.9) -- (III3.north);
\node[lbl, below=2pt of III1] {III-1};
\node[lbl, below=2pt of III2] {III-2};
\node[lbl, below=2pt of III3] {III-3};
% Numéros de générations
\node[gen] at (-3.6,0) {I};
\node[gen] at (-3.6,-1.9) {II};
\node[gen] at (-3.6,-3.7) {III};
\end{tikzpicture}
\end{center}
\legende{Pedigree d'une famille atteinte d'albinisme sur trois générations. Carré = homme ; rond = femme ; symbole plein = individu atteint (albinos) ; symbole vide = individu sain. Les générations sont numérotées I, II, III ; les individus I-1, I-2, etc.}
\end{popfigure}

\textbf{Interprétation.} Notons $A$ l'allèle dominant commandant une pigmentation normale et $a$ l'allèle récessif responsable de l'albinisme. L'individu I-2 est atteinte : son génotype est nécessairement $a/a$. Ses enfants II-1, II-2 et II-3 sont sains mais ont obligatoirement reçu l'allèle $a$ de leur mère : ils sont donc tous hétérozygotes $A/a$, \emph{porteurs sains}. Le couple II-2 ($A/a$) $\times$ conjoint (II-c) a un fils atteint III-2 ($a/a$) : le conjoint de II-2 est donc lui aussi porteur sain $A/a$.

\textbf{Conclusion.} L'albinisme se transmet selon un mode \cle{autosomique récessif} : les deux sexes sont atteints, des parents sains peuvent avoir des enfants atteints, et la maladie peut « sauter » des générations.

\begin{exemplebox}[Croisement de deux porteurs sains : les proportions à connaître]
Pour un couple de porteurs sains $A/a \times A/a$, l'échiquier de croisement donne :
\begin{center}
\begin{tabular}{|c|c|c|}
\hline
\rowcolor{popblueL} Gamètes & $A$ & $a$ \\
\hline
$A$ & $A/A$ sain & $A/a$ porteur sain \\
\hline
$a$ & $A/a$ porteur sain & $a/a$ \textbf{albinos} \\
\hline
\end{tabular}
\end{center}
Soit, à chaque naissance : \textbf{\SI{25}{\%}} d'enfants albinos ($a/a$), \textbf{\SI{50}{\%}} de porteurs sains ($A/a$), \textbf{\SI{25}{\%}} d'homozygotes sains ($A/A$). Attention : ces proportions sont des probabilités par naissance, non des effectifs garantis dans une fratrie réelle de petite taille.
\end{exemplebox}

\courssub{L'hérédité autosomique dominante : l'exemple de l'achondroplasie}

L'\cle{achondroplasie} est une maladie du développement osseux : les cartilages de conjugaison des os longs se développent mal, ce qui donne un nanisme disproportionné (membres courts, tronc de taille quasi normale, tête volumineuse). L'observation des pedigrees montre des faits très différents de ceux de l'albinisme :

\begin{itemize}
\item tout enfant atteint a \textbf{au moins un parent atteint} ;
\item la maladie apparaît à \textbf{toutes les générations} (pas de saut) ;
\item les \textbf{deux sexes} sont atteints, et un père atteint peut transmettre la maladie à son fils (ce qui exclut un gène porté par X).
\end{itemize}

\begin{propriete}[Critères d'une transmission autosomique dominante]
Un caractère est transmis selon le mode \cle{autosomique dominant} lorsque : (1) chaque individu atteint a au moins un parent atteint ; (2) le caractère apparaît à toutes les générations ; (3) les deux sexes sont atteints de façon comparable ; (4) la transmission père-fils est possible. Un individu atteint est le plus souvent \textbf{hétérozygote} $A/a$ (l'allèle morbide dominant étant noté ici $A$) ; le génotype homozygote $A/A$ est souvent létal ou exceptionnel, car deux individus atteints se rencontrent rarement.
\end{propriete}

Ainsi, un achondroplaste hétérozygote $A/a$ uni à une personne saine $a/a$ transmet l'allèle morbide à \textbf{\SI{50}{\%}} de ses enfants en moyenne : c'est le risque génétique que le médecin conseiller annoncera au couple.

\begin{attention}[Ne pas confondre les deux modes de transmission]
Deux réflexes à acquérir : (1) « parents sains, enfant atteint » $\Rightarrow$ allèle \textbf{récessif} ; (2) « toutes les générations touchées, au moins un parent atteint » $\Rightarrow$ allèle \textbf{dominant}. Et dans les deux cas, si garçons et filles sont atteints sans différence, le gène est \textbf{autosomique}. Ne jamais écrire « la maladie est dominante parce qu'elle est fréquente » : dominance et fréquence sont deux notions indépendantes (l'achondroplasie est dominante mais rare).
\end{attention}

\courssub{La drépanocytose : une hérédité autosomique à deux visages}

La \cle{drépanocytose} (ou anémie falciforme) est une maladie du globule rouge très fréquente en Afrique subsaharienne, donc au Cameroun. Elle est due à une mutation du gène codant la chaîne $\beta$ de l'hémoglobine : l'allèle normal $HbA$ devient l'allèle muté $HbS$ (substitution d'un acide glutamique par une valine en position 6 de la chaîne $\beta$).

\textbf{Observation.} Au microscope, les hématies d'un malade se déforment en « faucilles » (drépanocytes) lorsque la tension en dioxygène baisse ; ces cellules fragiles se lysent (anémie) et obstruent les capillaires (crises douloureuses). En revanche, les individus hétérozygotes sont cliniquement sains, mais l'électrophorèse de leur hémoglobine révèle \emph{deux} bandes : HbA \emph{et} HbS.

\textbf{Interprétation.} Il faut distinguer deux niveaux d'analyse du phénotype :

\begin{itemize}
\item \textbf{Au niveau moléculaire} (électrophorèse de l'hémoglobine), les allèles $HbA$ et $HbS$ sont \textbf{codominants} : l'hétérozygote $HbA/HbS$ synthétise les deux hémoglobines et présente un phénotype intermédiaire distinct des deux homozygotes.
\item \textbf{Au niveau de la maladie} (phénotype clinique), l'allèle $HbS$ est \textbf{récessif} : seul l'homozygote $HbS/HbS$ est malade ; l'hétérozygote $HbA/HbS$ est un porteur sain.
\end{itemize}

\begin{savaistu}[Drépanocytose et paludisme : l'avantage de l'hétérozygote]
Au Cameroun, environ \textbf{20 à 25 \%} de la population est hétérozygote $HbA/HbS$ (on dit « AS »). Pourquoi un allèle potentiellement mortel se maintient-il à une fréquence si élevée ? Parce que l'hétérozygote AS possède un \textbf{avantage sélectif en zone d'endémie palustre} : le plasmodium (parasite du paludisme) se développe mal dans ses hématies, qui falciforment légèrement et sont éliminées par la rate. L'hétérozygote résiste donc mieux au paludisme que l'homozygote $HbA/HbA$, sans souffrir de la drépanocytose comme l'homozygote $HbS/HbS$. C'est un cas célèbre d'\emph{équilibre polymorphe} par sélection naturelle. D'où l'importance capitale du \textbf{dépistage} (test d'électrophorèse, test de falciformation) avant le mariage : un couple AS $\times$ AS a, à chaque grossesse, \SI{25}{\%} de risque d'avoir un enfant drépanocytaire SS.
\end{savaistu}

\begin{exoresolu}[Prévision génétique : couple de porteurs sains du caractère drépanocytaire]
Un couple, tous deux hétérozygotes $HbA/HbS$, consulte avant d'avoir un enfant. Quel est le risque que l'enfant soit drépanocytaire ? Quelle est la probabilité qu'il soit porteur sain ?
\tcblower
Chaque parent produit deux types de gamètes en proportions égales : $1/2$ $HbA$ et $1/2$ $HbS$. L'échiquier de croisement donne :
\begin{align*}
P(HbS/HbS) &= \tfrac{1}{2}\times\tfrac{1}{2} = \tfrac{1}{4} = \SI{25}{\%} \quad \text{(enfant drépanocytaire)}\\
P(HbA/HbS) &= 2\times\tfrac{1}{2}\times\tfrac{1}{2} = \tfrac{1}{2} = \SI{50}{\%} \quad \text{(porteur sain)}\\
P(HbA/HbA) &= \tfrac{1}{4} = \SI{25}{\%} \quad \text{(non porteur)}
\end{align*}
Le risque que l'enfant soit drépanocytaire est donc de \textbf{\SI{25}{\%} à chaque grossesse}, et la probabilité qu'il soit porteur sain est de \textbf{\SI{50}{\%}}. Ce calcul de \cle{risque génétique} fonde le conseil génétique : il permet au couple de prendre une décision éclairée (dépistage prénatal possible).
\end{exoresolu}

\courssub{Les groupes sanguins ABO : allèles multiples et codominance}

\textbf{Observation.} Lors des transfusions sanguines pratiquées dans les hôpitaux (par exemple après une hémorragie du post-partum ou un accident de la route), on vérifie scrupuleusement le groupe sanguin du donneur et du receveur : une erreur peut être mortelle. Pourquoi ?

\textbf{Document expérimental.} Si l'on mélange sur lame le sang d'un individu de groupe A avec un sérum anti-A, les hématies s'agglutinent ; avec un sérum anti-B, rien ne se passe. Le sang d'un individu AB s'agglutine avec les deux sérums ; celui d'un individu O avec aucun. On en déduit que les hématies portent à leur surface des molécules appelées \cle{agglutinogènes} (antigènes A et/ou B), tandis que le plasma contient des anticorps naturels, les \cle{agglutinines} (anti-A et/ou anti-B), dirigés contre les agglutinogènes que l'individu ne possède pas.

\textbf{Interprétation génétique.} Le système ABO est déterminé par \emph{un seul gène} autosomique (chromosome 9) existant sous \textbf{trois formes alléliques} : c'est un cas d'\cle{allélisme multiple}.

\begin{itemize}
\item l'allèle $I^{A}$ commande la synthèse de l'agglutinogène A ;
\item l'allèle $I^{B}$ commande la synthèse de l'agglutinogène B ;
\item l'allèle $i$ (ou $O$) ne commande aucun agglutinogène.
\end{itemize}

$I^{A}$ et $I^{B}$ sont \textbf{codominants} entre eux (un hétérozygote $I^{A}/I^{B}$ exprime les deux agglutinogènes : groupe AB), et tous deux sont \textbf{dominants} sur $i$, qui est récessif.

\begin{popfigure}
\begin{center}
\begin{tikzpicture}
\node[font=\small\bfseries, text=popdark] at (5.4,4.6) {Système ABO : génotypes, phénotypes et compatibilités};
\foreach \x/\g/\ph/\agg/\agl/\c in {
  0/{$I^{A}/I^{A}$ ou $I^{A}/i$}/{A}/{A}/{anti-B}/popblue,
  1/{$I^{B}/I^{B}$ ou $I^{B}/i$}/{B}/{B}/{anti-A}/popgreen,
  2/{$I^{A}/I^{B}$}/{AB}/{A et B}/{aucune}/poporange,
  3/{$i/i$}/{O}/{aucun}/{anti-A + anti-B}/poppurple}{
  \pgfmathsetmacro{\yy}{3.4 - \x*1.05}
  \node[draw=\c, thick, rounded corners=2pt, fill=\c!12, minimum width=3.4cm, minimum height=0.8cm, font=\small] at (1.9,\yy) {\g};
  \node[draw=\c, thick, rounded corners=2pt, fill=\c!25, minimum width=1.2cm, minimum height=0.8cm, font=\small\bfseries] at (4.4,\yy) {\ph};
  \node[draw=\c, thick, rounded corners=2pt, fill=white, minimum width=2.2cm, minimum height=0.8cm, font=\small] at (6.4,\yy) {\agg};
  \node[draw=\c, thick, rounded corners=2pt, fill=white, minimum width=2.6cm, minimum height=0.8cm, font=\small] at (9.1,\yy) {\agl};
}
\node[font=\small\bfseries, text=popmuted] at (1.9,4.05) {Génotypes};
\node[font=\small\bfseries, text=popmuted] at (4.4,4.05) {Groupe};
\node[font=\small\bfseries, text=popmuted] at (6.4,4.05) {Agglutinogènes};
\node[font=\small\bfseries, text=popmuted] at (9.1,4.05) {Agglutinines};
\end{tikzpicture}
\end{center}
\legende{Correspondances entre génotypes, groupes sanguins (phénotypes), agglutinogènes portés par les hématies et agglutinines du plasma dans le système ABO. Les allèles $I^{A}$ et $I^{B}$ sont codominants et dominants sur l'allèle récessif $i$.}
\end{popfigure}

\begin{propriete}[Règles de compatibilité transfusionnelle]
On peut transfuser un sang dont les hématies ne portent \textbf{aucun agglutinogène} contre lequel le receveur possède une agglutinine. Ainsi : l'individu O, sans agglutinogène, est \textbf{donneur universel} ; l'individu AB, sans agglutinine, est \textbf{receveur universel}. Un individu A peut recevoir du A et du O ; un individu B du B et du O.
\end{propriete}

\begin{exemplebox}[Détermination des génotypes possibles]
Un homme de groupe A et une femme de groupe B peuvent-ils avoir un enfant de groupe O ? Oui, si et seulement si les deux parents sont hétérozygotes : $I^{A}/i \times I^{B}/i$ donne $1/4$ $I^{A}/I^{B}$ (AB), $1/4$ $I^{A}/i$ (A), $1/4$ $I^{B}/i$ (B), $1/4$ $i/i$ (O). En revanche, un couple $I^{A}/I^{A} \times I^{B}/I^{B}$ ne peut avoir que des enfants AB. Ce type de raisonnement est classique au Baccalauréat (et en médecine légale pour les contestations de paternité, avec prudence).
\end{exemplebox}

\courssub{Le facteur rhésus : dominance autosomique et incompatibilité fœto-maternelle}

Le \cle{facteur rhésus} est un autre antigène membranaire des hématies, découvert chez le singe \emph{Macacus rhesus}. Deux phénotypes existent :

\begin{itemize}
\item \textbf{Rh+} : les hématies portent l'agglutinogène rhésus (environ \SI{85}{\%} des individus) ;
\item \textbf{Rh-} : absence de l'agglutinogène.
\end{itemize}

La transmission est \cle{autosomique dominante} : l'allèle $Rh^{+}$ est dominant, l'allèle $Rh^{-}$ récessif. Un individu Rh+ est de génotype $Rh^{+}/Rh^{+}$ ou $Rh^{+}/Rh^{-}$ ; un individu Rh- est obligatoirement $Rh^{-}/Rh^{-}$.

\begin{aretenir}[L'incompatibilité fœto-maternelle rhésus (rappel de Première)]
Contrairement au système ABO, il n'existe pas d'agglutinines anti-rhésus naturelles : elles n'apparaissent qu'après \textbf{immunisation}. Si une mère Rh- porte un fœtus Rh+ (allèle transmis par le père), des hématies fœtales peuvent passer dans le sang maternel, surtout \emph{à l'accouchement} : la mère produit alors des anticorps anti-Rh. Lors d'une \textbf{grossesse suivante} d'un nouveau fœtus Rh+, ces anticorps (des IgG) traversent le placenta et détruisent les hématies du fœtus : c'est la \textbf{maladie hémolytique du nouveau-né} (anémie, ictère grave). La prévention repose sur l'injection d'immunoglobulines anti-D à la mère Rh- juste après chaque accouchement d'un enfant Rh+, ce qui détruit les hématies fœtales avant qu'elles n'immunisent la mère.
\end{aretenir}

\begin{exercice}[Pour s'entraîner]
Dans une famille, le père est de groupe AB Rh+, la mère de groupe O Rh-. (1) Quels sont les groupes sanguins ABO possibles pour leurs enfants ? (2) La mère peut-elle être confrontée à un risque d'incompatibilité fœto-maternelle ? Justifier. (3) Le couple a un enfant de groupe B : quels génotypes complets (ABO) peut-on attribuer au père et à l'enfant ?
\end{exercice}

\begin{corrige}[Exercice 1]
(1) Le père AB est $I^{A}/I^{B}$ ; la mère O est $i/i$. Les gamètes paternels sont $I^{A}$ ou $I^{B}$, les gamètes maternels sont tous $i$. Les enfants sont donc $I^{A}/i$ (groupe A) ou $I^{B}/i$ (groupe B), à parts égales : \textbf{\SI{50}{\%} A, \SI{50}{\%} B}. Aucun enfant AB ni O n'est possible. (2) Oui : la mère étant Rh- ($Rh^{-}/Rh^{-}$), si le père transmet l'allèle dominant $Rh^{+}$, le fœtus sera Rh+ ; lors d'une seconde grossesse Rh+, les anticorps maternels anti-Rh pourront provoquer une maladie hémolytique du nouveau-né. (3) Le père est nécessairement $I^{A}/I^{B}$ ; l'enfant de groupe B, ayant reçu $i$ de sa mère et $I^{B}$ de son père, est $I^{B}/i$.
\end{corrige}

En résumé, l'étude des pedigrees permet, sans aucun croisement expérimental, de déterminer le mode de transmission des caractères humains : récessif ou dominant, autosomique ou — comme nous le verrons dans la prochaine section — lié au sexe, avec des exemples célèbres comme le daltonisme et l'hémophilie.

\coursec{Hérédité gonosomique (liée au sexe)}

\courssub{Particularités des chromosomes sexuels et définition d'un gène lié au sexe}

Rappelons-nous : dans la section précédente, nous avons étudié l'hérédité autosomique, où les gènes concernés sont portés par des chromosomes non sexuels (autosomes), présents en deux exemplaires chez l'homme comme chez la femme. La transmission ne dépend donc pas du sexe de l'individu. Or, certaines maladies humaines — le daltonisme, l'hémophilie — touchent très majoritairement les garçons. Cette observation clinique oriente immédiatement le généticien vers une hypothèse : \emph{les gènes responsables sont portés par les chromosomes sexuels, les gonosomes}.

\begin{experience}[Observation clinique : prévalence sexuée du daltonisme et de l'hémophilie]
\textbf{Question de départ.} Pourquoi le daltonisme et l'hémophilie touchent-ils environ 8\,\% des hommes contre seulement 0,5\,\% des femmes au Cameroun comme ailleurs ?
\textbf{Document.} Statistiques épidémiologiques (OMS, 2023) : prévalence du daltonisme rouge-vert : hommes $\approx$ 8\,\%, femmes $\approx$ 0,5\,\% ; hémophilie A : incidence $\approx$ 1/5\,000 naissances mâles, quasi nulle chez les femmes.
\textbf{Hypothèse.} Les allèles mutés responsables sont portés par le chromosome X. Comme l'homme ne possède qu'un seul X (caryotype 46,XY), un seul allèle muté suffit à exprimer le phénotype. La femme (46,XX) a deux X ; elle ne sera atteinte que si elle reçoit l'allèle muté de \emph{son père} \emph{et} de \emph{sa mère} — événement rare.
\textbf{Interprétation.} Cette différence de prévalence entre sexes est la signature expérimentale d'une \textbf{hérédité liée au chromosome X} (ou gonosomique récessive liée à X).
\textbf{Conclusion.} Les gènes du daltonisme et de l'hémophilie sont localisés sur la \textbf{région différentielle du chromosome X}, absente du Y.
\end{experience}

\begin{definition}[Gène lié au sexe (ou gène gonosomique)]
Un \cle{gène lié au sexe} est un gène porté par l'un des chromosomes sexuels (gonosomes). Dans le programme de Terminale, on étudie presque exclusivement les gènes portés par la \textbf{région différentielle du chromosome X} (portion non homologue au Y). Le chromosome Y possède aussi sa région différentielle (gène \emph{SRY} déterminant la masculinité), mais elle ne porte pas de gènes de caractères « classiques » comme le daltonisme ou l'hémophilie.
\end{definition}

\begin{propriete}[Transmission différentielle des gonosomes par le père]
\textbf{Admise au Bac.} Lors de la méiose chez l'homme (46,XY), les deux gonosomes se séparent :
\begin{itemize}
\item La moitié des spermatozoïdes reçoit le chromosome \textbf{X} ;
\item L'autre moitié reçoit le chromosome \textbf{Y}.
\end{itemize}
\emph{Conséquence directe :} un père transmet \textbf{son X à toutes ses filles} (qui reçoivent son X + le X maternel) et \textbf{son Y à tous ses fils} (qui reçoivent son Y + le X maternel). Il n'y a \textbf{jamais} de transmission père $\to$ fils du chromosome X.
\end{propriete}

\begin{attention}[Piège décisif en analyse de pedigree]
Si un pedigree montre un père atteint transmettant le caractère à son fils, \textbf{l'hérédité ne peut pas être liée au X récessive}. Ce critère (absence de transmission père-fils) est l'indice le plus fiable pour reconnaître une hérédité liée à X sur un arbre généalogique.
\end{attention}

\courssub{Le daltonisme : modèle d'hérédité récessive liée à X}

Le daltonisme rouge-vert (protanopie/deutéranopie) est dû à une mutation d'un gène codant pour une opsin (pigment photorécepteur) sur le chromosome X. L'allèle normal (vision trichromatique) est noté $X^{D}$ (dominant), l'allèle muté (daltonisme) $X^{d}$ (récessif).

\begin{exemplebox}[Croisement type : femme conductrice $\times$ homme normal]
Soit une femme \textbf{conductrice} (hétérozygote, phénotype normal) de génotype $X^{D}X^{d}$ et un homme \textbf{normal} de génotype $X^{D}Y$. Construisons l'échiquier de croisement (car les gamètes mâles sont de deux types selon le gonosome).
\end{exemplebox}

\begin{popfigure}
\begin{center}
\begin{tabularx}{\linewidth}{|>{\centering\arraybackslash}X|>{\centering\arraybackslash}X|>{\centering\arraybackslash}X|}
\hline
\rowcolor{popblueL}
\textbf{Gamètes maternels} $\downarrow$ / \textbf{Gamètes paternels} $\rightarrow$ & $X^{D}$ (0,5) & $Y$ (0,5) \\
\hline
$X^{D}$ (0,5) & $X^{D}X^{D}$ \newline Fille normale (0,25) & $X^{D}Y$ \newline Garçon normal (0,25) \\
\hline
$X^{d}$ (0,5) & $X^{D}X^{d}$ \newline Fille conductrice (0,25) & $X^{d}Y$ \newline \textbf{Garçon daltonien} (0,25) \\
\hline
\end{tabularx}
\end{center}
\legende{Échiquier de croisement $X^{D}X^{d} \times X^{D}Y$ (daltonisme). Les fréquences gamétiques sont indiquées entre parenthèses.}
\end{popfigure}

\begin{aretenir}[Bilan du croisement femme conductrice $\times$ homme normal]
\begin{itemize}
\item \textbf{Filles :} 50\,\% normales ($X^{D}X^{D}$), 50\,\% conductrices ($X^{D}X^{d}$) — \textbf{0\,\% atteintes}.
\item \textbf{Garçons :} 50\,\% normaux ($X^{D}Y$), 50\,\% \textbf{daltoniens} ($X^{d}Y$) — \textbf{50\,\% atteints}.
\item Globalement : 25\,\% de la descendance atteinte (tous des garçons).
\end{itemize}
\end{aretenir}

\begin{exemplebox}[Pourquoi les garçons sont-ils plus souvent atteints ? — L'hémizygotie]
Chez l'homme (XY), le chromosome X est présent en \textbf{un seul exemplaire} pour les gènes de sa région différentielle. On dit que l'homme est \cle{hémizygote} pour ces gènes. Il n'y a pas d'allèle sur le Y pour masquer un allèle récessif muté sur X. Donc :
\begin{itemize}
\item Un seul allèle muté $X^{d}$ suffit à exprimer le phénotype mutant chez le garçon ($X^{d}Y$).
\item Chez la fille (XX), il faut \textbf{deux} allèles mutés ($X^{d}X^{d}$) pour être atteinte — ce qui impose que le père soit atteint \emph{et} que la mère soit au moins conductrice.
\end{itemize}
C'est la base mathématique de la différence de prévalence (8\,\% vs 0,5\,\%).
\end{exemplebox}

\begin{savaistu}[Prévalence du daltonisme au Cameroun et dans le monde]
Environ 8\,\% des hommes et 0,5\,\% des femmes sont daltoniens (rouge-vert). Ce rapport $\approx$ 16:1 s'explique parfaitement par la liaison à X récessive : si la fréquence de l'allèle $d$ dans la population est $q \approx 0,08$, alors fréquence des hommes atteints $= q = 8\,\%$ ; fréquence des femmes atteintes $= q^{2} = 0,0064 \approx 0,6\,\%$ (proche de 0,5\,\% observé). La différence n'est pas due à une « sensibilité » biologique, mais à la \textbf{structure chromosomique}.
\end{savaistu}

\courssub{L'hémophilie : un exemple historique et médical majeur}

L'hémophilie A (déficit en facteur VIII) et B (déficit en facteur IX) sont des maladies hémorragiques graves, transmises comme caractères \textbf{récessifs liés au X}. Le gène \emph{F8} (hémophilie A) est en Xq28 ; le gène \emph{F9} (hémophilie B) en Xq27.1.

\begin{experience}[L'hémophilie dans la descendance de la reine Victoria — étude de document historique]
\textbf{Contexte.} La reine Victoria du Royaume-Uni (1819–1901) était conductrice de l'hémophilie B (mutation \emph{de novo} dans sa lignée germinale ou héritée de sa mère). Elle a transmis l'allèle muté à plusieurs de ses neuf enfants.
\textbf{Document simplifié (pedigree partiel).}
\begin{itemize}
\item Victoria (conductrice) $\times$ Albert (normal) $\to$ parmi leurs enfants : Alice (conductrice), Léopold (atteint), Béatrice (conductrice).
\item Alice (conductrice) $\times$ Louis IV de Hesse $\to$ Alexis (atteint, tsarévitch de Russie), Irène (conductrice), Alix (conductrice, future tsarine Alexandra).
\item Alix (conductrice) $\times$ Nicolas II de Russie $\to$ Alexis (atteint, unique fils, hémophile) — sa maladie a influencé l'histoire politique (Raspoutine, révolution).
\item Béatrice (conductrice) $\times$ Henri de Battenberg $\to$ Léopold (atteint), Maurice (atteint), Victoria-Eugénie (conductrice, reine d'Espagne) $\to$ deux fils atteints.
\end{itemize}
\textbf{Question.} Représentez ce pedigree simplifié et vérifiez les critères de l'hérédité liée à X récessive.
\end{experience}

\begin{popfigure}
\begin{center}
\begin{tikzpicture}[
  person/.style={draw, thick, minimum width=1.2cm, minimum height=0.8cm, rounded corners=3pt},
  male/.style={person, rectangle, fill=popblueL},
  female/.style={person, circle, fill=poppinkL},
  affected/.style={fill=poporange!80},
  carrier/.style={pattern=north east lines, pattern color=poporange},
  generation/.style={font=\bfseries\small, text=popdark},
  link/.style={thick, popdark}
]
% Generation 1
\node[generation] (G1) at (0,4) {G1};
\node[female, carrier] (V) at (2,4) {Victoria};
\node[male] (A) at (5,4) {Albert};
\draw[link] (V) -- (A);
% Generation 2
\node[generation] (G2) at (0,2.5) {G2};
\node[female, carrier] (Alice) at (0.5,2.5) {Alice};
\node[male, affected] (Leopold) at (2.5,2.5) {Léopold};
\node[female, carrier] (Beatrice) at (4.5,2.5) {Béatrice};
\draw[link] (V) -- (Alice);
\draw[link] (V) -- (Leopold);
\draw[link] (V) -- (Beatrice);
% Generation 3 - Alice line
\node[generation] (G3a) at (0,1) {G3};
\node[male, affected] (AlexisR) at (-1,1) {Alexis (Russie)};
\node[female, carrier] (Irene) at (0.5,1) {Irène};
\node[female, carrier] (Alix) at (2,1) {Alix (Tsarine)};
\draw[link] (Alice) -- (AlexisR);
\draw[link] (Alice) -- (Irene);
\draw[link] (Alice) -- (Alix);
% Generation 3 - Beatrice line
\node[male, affected] (LeopoldB) at (4,1) {Léopold};
\node[male, affected] (Maurice) at (5,1) {Maurice};
\node[female, carrier] (VicEug) at (6,1) {Vic.-Eug.};
\draw[link] (Beatrice) -- (LeopoldB);
\draw[link] (Beatrice) -- (Maurice);
\draw[link] (Beatrice) -- (VicEug);
% Generation 4 - Alix line
\node[generation] (G4) at (0,-0.5) {G4};
\node[male, affected] (Alexis2) at (2,-0.5) {Alexis (héritier)};
\draw[link] (Alix) -- (Alexis2);
% Generation 4 - VicEug line
\node[male, affected] (Alfonso) at (5.5,-0.5) {Alphonse};
\node[male, affected] (Gonzalo) at (6.5,-0.5) {Gonzalo};
\draw[link] (VicEug) -- (Alfonso);
\draw[link] (VicEug) -- (Gonzalo);
% Legend
\node[anchor=west] at (8,3.5) {\textbf{Légende :}};
\node[male, anchor=west] at (8,3) {}; \node[anchor=west] at (9.3,3) {Homme normal};
\node[male, affected, anchor=west] at (8,2.5) {}; \node[anchor=west] at (9.3,2.5) {Homme atteint};
\node[female, anchor=west] at (8,2) {}; \node[anchor=west] at (9.3,2) {Femme normale};
\node[female, carrier, anchor=west] at (8,1.5) {}; \node[anchor=west] at (9.3,1.5) {Femme conductrice};
\end{tikzpicture}
\end{center}
\legende{Pedigree simplifié de l'hémophilie B dans la descendance de la reine Victoria. Carrés = hommes, cercles = femmes. Rempli orange = atteint. Hachuré = conductrice (phénotype normal, un allèle muté).}
\end{popfigure}

\begin{aretenir}[Critères de reconnaissance d'une hérédité liée à X récessive sur un pedigree]
\begin{enumerate}
\item \textbf{Prédominance masculine :} beaucoup plus d'hommes atteints que de femmes.
\item \textbf{Transmission par les femmes conductrices :} les hommes atteints ont une mère conductrice (ou atteinte) ; les femmes conductrices ont un père atteint.
\item \textbf{Absence de transmission père $\to$ fils :} un père atteint transmet son Y à ses fils, jamais son X muté. Ses fils sont donc toujours sains (pour ce caractère).
\item \textbf{Toutes les filles d'un père atteint sont conductrices :} elles reçoivent obligatoirement son X muté.
\item \textbf{Saut de génération apparent :} le caractère peut passer d'un grand-père maternel à un petit-fils via une fille conductrice phénotypiquement normale.
\end{enumerate}
\end{aretenir}

\courssub{Calcul de risques génétiques dans les croisements gonosomiques}

La prévision génétique (conseil génétique) consiste à calculer la probabilité qu'un enfant soit atteint ou conducteur, connaissant les génotypes des parents. La méthode est rigoureuse : \textbf{1. Identifier les génotypes parentaux. 2. Déterminer les gamètes possibles et leurs fréquences. 3. Croiser les gamètes (échiquier ou produit de probabilités). 4. Lire les phénotypes par sexe.}

\begin{methode}[Calcul du risque génétique pour un couple : femme conductrice $\times$ homme normal]
\textbf{Étape 1 : Génotypes.} Mère $X^{D}X^{d}$ (conductrice), Père $X^{D}Y$ (normal).
\textbf{Étape 2 : Gamètes.}
Mère : $X^{D}$ (1/2) et $X^{d}$ (1/2) — \emph{pas de Y}.
Père : $X^{D}$ (1/2) et $Y$ (1/2).
\textbf{Étape 3 : Croisement (produit des probabilités).}
\begin{align*}
P(X^{D}X^{D}) &= \tfrac{1}{2} \times \tfrac{1}{2} = \tfrac{1}{4} \quad \text{(fille normale)} \\
P(X^{D}X^{d}) &= \tfrac{1}{2} \times \tfrac{1}{2} = \tfrac{1}{4} \quad \text{(fille conductrice)} \\
P(X^{D}Y) &= \tfrac{1}{2} \times \tfrac{1}{2} = \tfrac{1}{4} \quad \text{(garçon normal)} \\
P(X^{d}Y) &= \tfrac{1}{2} \times \tfrac{1}{2} = \tfrac{1}{4} \quad \text{(garçon atteint)}
\end{align*}
\textbf{Étape 4 : Risques conditionnés par le sexe.}
\begin{itemize}
\item Si l'enfant est une \textbf{fille} : risque d'être atteinte = 0 ; risque d'être conductrice = 1/2.
\item Si l'enfant est un \textbf{garçon} : risque d'être atteint = 1/2 ; risque d'être normal = 1/2.
\item \textbf{Risque global} (sans connaître le sexe) : 1/4 d'enfant atteint (toujours un garçon).
\end{itemize}
\end{methode}

\begin{exoresolu}[Conseil génétique : couple dont le fils aîné est hémophile]
\textbf{Énoncé.} Un couple consulte. Leur premier fils, âgé de 3 ans, est atteint d'hémophilie A (récessive liée à X). La mère est phénotypiquement normale. Le père est normal. Ils souhaitent avoir un deuxième enfant. Calculer : (a) le génotype de la mère ; (b) la probabilité que le deuxième enfant soit une fille atteinte ; (c) la probabilité qu'un deuxième garçon soit atteint.
\tcblower
\textbf{Solution détaillée.}
\textbf{(a) Génotype de la mère.} Le fils atteint a pour génotype $X^{h}Y$ (allèle $h$ = hémophilie). Il a reçu son Y de son père, donc son X muté ($X^{h}$) \textbf{vient obligatoirement de sa mère}. La mère est phénotypiquement normale $\to$ elle est \textbf{conductrice} : $X^{H}X^{h}$.
\textbf{(b) Fille atteinte ?} Croisement $X^{H}X^{h} \times X^{H}Y$.
Gamètes mère : $X^{H}$ (1/2), $X^{h}$ (1/2). Gamètes père : $X^{H}$ (1/2), $Y$ (1/2).
Filles possibles : $X^{H}X^{H}$ (normale) et $X^{H}X^{h}$ (conductrice). \textbf{Aucune fille $X^{h}X^{h}$ n'est possible car le père donne $X^{H}$.} Probabilité = \textbf{0}.
\textbf{(c) Garçon atteint ?} Garçons possibles : $X^{H}Y$ (normal) et $X^{h}Y$ (atteint). Probabilité = \textbf{1/2} (50\,\%).
\textbf{Conclusion.} Pour chaque grossesse : 25\,\% de chance d'avoir un garçon atteint, 25\,\% d'une fille conductrice, 25\,\% d'un garçon normal, 25\,\% d'une fille normale. Risque global d'enfant atteint = 25\,\%.
\end{exoresolu}

\begin{attention}[Erreurs fréquentes dans les calculs de risques gonosomiques]
\begin{enumerate}
\item \textbf{Oublier le conditionnement par le sexe :} écrire « 50\,\% de risque d'enfant atteint » au lieu de « 50\,\% de risque \emph{si c'est un garçon}, 0\,\% \emph{si c'est une fille}, 25\,\% global ».
\item \textbf{Confondre conductrice et atteinte :} une femme $X^{H}X^{h}$ est \emph{conductrice} (phénotype normal), pas atteinte. Elle ne saigne pas anormalement (pour l'hémophilie) ou voit normalement (pour le daltonisme) — sauf inactivation X biaisée, hors programme.
\item \textbf{Attribuer un allèle au chromosome Y :} le Y ne porte pas l'allèle du daltonisme ni de l'hémophilie. Noter $X^{h}Y$, jamais $X^{h}Y^{h}$.
\item \textbf{Inverser la transmission :} croire qu'un père atteint peut avoir un fils atteint. \textbf{Jamais} pour un caractère lié à X.
\end{enumerate}
\end{attention}

\courssub{Le cas rare de la femme atteinte : homozygotie imposée}

Une femme ne peut être atteinte (phénotype mutant) que si elle est \textbf{homozygote mutée} ($X^{d}X^{d}$). Cela impose deux conditions \textbf{simultanées} :
\begin{enumerate}
\item Son \textbf{père est atteint} (il lui transmet son unique X muté).
\item Sa \textbf{mère est au moins conductrice} (elle lui transmet un X muté avec probabilité 1/2 si conductrice, 1 si atteinte).
\end{enumerate}
C'est un événement rare (fréquence $\approx q^{2}$ dans la population), ce qui explique pourquoi les femmes atteintes sont exceptionnelles. Au Cameroun, pour la drépanocytose (autosomique) on voit des femmes atteintes fréquemment ; pour l'hémophilie, c'est quasi exclusivement masculin.

\begin{exemplebox}[Croisement femme atteinte $\times$ homme normal]
Génotypes : Mère $X^{d}X^{d}$, Père $X^{D}Y$.
Gamètes mère : uniquement $X^{d}$. Gamètes père : $X^{D}$ (1/2), $Y$ (1/2).
Descendance :
\begin{itemize}
\item Filles : toutes $X^{D}X^{d}$ $\to$ \textbf{toutes conductrices} (phénotype normal).
\item Garçons : tous $X^{d}Y$ $\to$ \textbf{tous atteints}.
\end{itemize}
C'est l'unique situation où \textbf{100\,\% des garçons sont atteints} et \textbf{0\,\% des filles sont atteintes} (mais 100\,\% conductrices).
\end{exemplebox}

\courssub{Synthèse comparative : autosomique vs gonosomique}

\begin{center}
\begin{tabularx}{\linewidth}{|l|X|X|}
\hline
\rowcolor{popblueL}
\textbf{Critère} & \textbf{Hérédité autosomique} & \textbf{Hérédité gonosomique liée à X (récessive)} \\
\hline
Chromosome porteur & Autosome (paires 1–22) & Chromosome X (région différentielle) \\
\hline
Fréquence selon le sexe & Identique chez ♂ et ♀ & ♂ $\gg$ ♀ (hémizygotie) \\
\hline
Transmission père $\to$ fils & Possible & \textbf{Impossible} \\
\hline
Transmission père $\to$ fille & 50\,\% (si hétérozygote) & \textbf{100\,\% conductrices} (si père atteint) \\
\hline
Femme atteinte & Homozygote mutée & Homozygote mutée (père atteint + mère conductrice) \\
\hline
Conducteur sain & Hétérozygote (♂ et ♀) & \textbf{Uniquement la femme} (hétérozygote $X^{D}X^{d}$) \\
\hline
Exemples (programme) & Albinisme, drépanocytose, achondroplasie, ABO, Rhésus & Daltonisme, hémophilie A et B \\
\hline
\end{tabularx}
\end{center}

\begin{exercice}[Application : analyse d'un pedigree familial camerounais]
\textbf{Énoncé.} Dans une famille de Yaoundé, le grand-père maternel est daltonien. Sa fille (la mère de l'enfant consulté) a une vision normale. Elle épouse un homme à vision normale. Ils ont un fils de 6 ans qui confond le rouge et le vert à l'école.
\begin{enumerate}
\item Représentez le pedigree sur 3 générations (symboles standards).
\item Déduisez le génotype de la mère.
\item Calculez la probabilité qu'un deuxième enfant soit une fille daltonienne.
\item Calculez la probabilité qu'un deuxième garçon soit daltonien.
\end{enumerate}
\end{exercice}

\begin{corrige}[Exercice 1]
\textbf{1. Pedigree.}
G1 : Grand-père maternel (carré rempli orange) $\times$ Grand-mère (cercle normal, génotype inconnu mais au moins $X^{D}X^{?}$).
G2 : Mère (cercle hachuré = conductrice) $\times$ Père (carré normal).
G3 : Fils aîné (carré rempli = daltonien) ; deuxième enfant (sexe inconnu).
\textbf{2. Génotype de la mère.} Le grand-père est $X^{d}Y$. Il transmet son $X^{d}$ à \textbf{toutes} ses filles. La mère a une vision normale $\to$ elle est \textbf{conductrice} $X^{D}X^{d}$.
\textbf{3. Fille daltonienne ?} Croisement $X^{D}X^{d} \times X^{D}Y$. Filles possibles : $X^{D}X^{D}$ et $X^{D}X^{d}$. Aucune $X^{d}X^{d}$ possible (père donne $X^{D}$). Probabilité = \textbf{0}.
\textbf{4. Garçon daltonien ?} Garçons possibles : $X^{D}Y$ (normal) et $X^{d}Y$ (daltonien). Probabilité = \textbf{1/2}.
\end{corrige}

\begin{savaistu}[Dépistage néonatal et prise en charge au Cameroun]
Au Cameroun, le dépistage néonatal de la drépanocytose (autosomique) est en cours de généralisation dans les grandes villes (Yaoundé, Douala, Garoua). Pour l'hémophilie, le diagnostic repose sur le dosage des facteurs VIII/IX (Centre National de Transfusion Sanguine, Hôpital Général de Yaoundé). La prise en charge (facteurs recombinants, prophylaxie) reste coûteuse. Le conseil génétique prénuptial, mené par des biologistes médicaux et généticiens cliniciens, permet d'informer les couples conducteurs sur les risques (25\,\% d'enfant atteint par grossesse pour l'hémophilie si mère conductrice) et d'envisager un diagnostic prénatal (biopsie choriales, amniocentèse) dans les structures spécialisées.
\end{savaistu}

\coursec{Mutations, risques génétiques et prévention}

Après avoir exploré les mécanismes de la transmission des caractères selon les lois de Mendel, tant pour les gènes portés par les autosomes que pour ceux localisés sur les gonosomes, nous abordons maintenant la source ultime de toute variabilité génétique : la \cle{mutation}. Sans mutation, pas d'allèles nouveaux, pas de diversité sur laquelle la sélection naturelle peut agir, et donc pas d'évolution. Pourtant, dans nos sociétés, le mot « mutation » évoque souvent la maladie. Il est temps de démêler le biologique du pathologique.

\courssub{Les mutations : définition et classification générale}

\begin{definition}[Mutation]
Une \textbf{mutation} est une modification \textbf{accidentelle, stable et héréditaire} de la séquence nucléotidique de l'ADN (mutation génique) ou de l'organisation des chromosomes (mutation chromosomique). Elle constitue la source première de la \cle{variabilité génétique} au sein des populations.
\end{definition}

Les mutations peuvent affecter :
\begin{itemize}
    \item un seul gène (ou quelques nucléotides) : \textbf{mutations géniques} (ou ponctuelles) ;
    \item la structure d'un chromosome : \textbf{mutations chromosomiques structurales} (délétion, duplication, inversion, translocation) ;
    \item le nombre de chromosomes : \textbf{mutations chromosomiques numériques} (aneuploïdie, polyploïdie).
\end{itemize}

\begin{aretenir}[Classification des mutations]
\begin{center}
\begin{tabularx}{\linewidth}{|l|X|X|}\hline
\textbf{Type} & \textbf{Niveau} & \textbf{Exemples} \\ \hline
Génique & Séquence d'ADN (1 à quelques nucléotides) & Drépanocytose (substitution), mucoviscidose (délétion $\Delta$F508) \\ \hline
Chromosomique structurelle & Organisation du chromosome & Syndrome du cri du chat (délétion 5p), leucémie myéloïde chronique (translocation 9;22) \\ \hline
Chromosomique numérique & Nombre de chromosomes & Trisomie 21 (aneuploïdie), triploïdie (polyploïdie) \\ \hline
\end{tabularx}
\end{center}
\end{aretenir}

\courssub{Mutations géniques : quand la lettre change}

Une mutation génique altère la séquence codante d'un gène. Trois mécanismes élémentaires existent sur une courte portion d'ADN :

\begin{popfigure}
\begin{tikzpicture}[scale=0.9, every node/.style={font=\small}]
% Couleurs
\colorlet{baseA}{poppink!80}
\colorlet{baseT}{popblue!80}
\colorlet{baseC}{popgreen!80}
\colorlet{baseG}{poporange!80}
\colorlet{backbone}{popdark!60}

% Légende bases
\node[anchor=west] at (-6, 5.5) {\textbf{Légende :}};
\foreach \b/\c/\y in {A/baseA/5.0, T/baseT/4.6, C/baseC/4.2, G/baseG/3.8} {
    \fill[\c] (-5.5, \y) rectangle (-5.0, \y+0.3);
    \node[anchor=west] at (-4.8, \y+0.15) {\b};
}

% Séquence originale
\node[anchor=east] at (-1, 5.2) {\textbf{Séquence originale :}};
\draw[backbone, thick] (0, 5.3) -- (10, 5.3);
\draw[backbone, thick] (0, 4.7) -- (10, 4.7);
\foreach \i/\b/\cb in {0/A/T, 1/T/A, 2/C/G, 3/G/C, 4/A/T, 5/T/A, 6/C/G, 7/G/C} {
    \pgfmathsetmacro{\x}{0.5 + \i*1.2}
    \fill[base\b] (\x, 4.85) rectangle (\x+0.4, 5.15);
    \node[white, font=\tiny\bfseries] at (\x+0.2, 5.0) {\b};
    % Brin complémentaire
    \fill[base\cb] (\x, 4.7) rectangle (\x+0.4, 4.95);
    \node[white, font=\tiny\bfseries] at (\x+0.2, 4.825) {\cb};
}

% Substitution
\node[anchor=east] at (-1, 3.2) {\textbf{Substitution (A$\to$G) :}};
\draw[backbone, thick] (0, 3.3) -- (10, 3.3);
\draw[backbone, thick] (0, 2.7) -- (10, 2.7);
\foreach \i/\b/\cb in {0/A/T, 1/T/A, 2/C/G, 3/G/C, 4/G/C, 5/T/A, 6/C/G, 7/G/C} { % 4ème base changée A->G
    \pgfmathsetmacro{\x}{0.5 + \i*1.2}
    \fill[base\b] (\x, 2.85) rectangle (\x+0.4, 3.15);
    \node[white, font=\tiny\bfseries] at (\x+0.2, 3.0) {\b};
    \fill[base\cb] (\x, 2.7) rectangle (\x+0.4, 2.95);
    \node[white, font=\tiny\bfseries] at (\x+0.2, 2.825) {\cb};
}
% Marqueur
\draw[poppink, thick, decorate, decoration={snake, amplitude=1pt, segment length=4pt}] (5.3, 3.3) -- (5.3, 3.6);
\node[poppink, anchor=south, font=\tiny\bfseries] at (5.3, 3.6) {Muté};

% Délétion
\node[anchor=east] at (-1, 1.2) {\textbf{Délétion (perte de C) :}};
\draw[backbone, thick] (0, 1.3) -- (8.8, 1.3);
\draw[backbone, thick] (0, 0.7) -- (8.8, 0.7);
\foreach \i/\b/\cb in {0/A/T, 1/T/A, 2/G/C, 3/A/T, 4/T/A, 5/C/G, 6/G/C} { % base 2 (C) supprimée
    \pgfmathsetmacro{\x}{0.5 + \i*1.2}
    \fill[base\b] (\x, 0.85) rectangle (\x+0.4, 1.15);
    \node[white, font=\tiny\bfseries] at (\x+0.2, 1.0) {\b};
    \fill[base\cb] (\x, 0.7) rectangle (\x+0.4, 0.95);
    \node[white, font=\tiny\bfseries] at (\x+0.2, 0.825) {\cb};
}
% Marqueur
\draw[poppink, thick, decorate, decoration={snake, amplitude=1pt, segment length=4pt}] (2.9, 1.3) -- (2.9, 1.6);
\node[poppink, anchor=south, font=\tiny\bfseries] at (2.9, 1.6) {Supprimé};

% Insertion
\node[anchor=east] at (-1, -0.8) {\textbf{Insertion (ajout de T) :}};
\draw[backbone, thick] (0, -0.7) -- (11.2, -0.7);
\draw[backbone, thick] (0, -1.3) -- (11.2, -1.3);
\foreach \i/\b/\cb in {0/A/T, 1/T/A, 2/C/G, 3/T/A, 4/G/C, 5/A/T, 6/T/A, 7/C/G, 8/G/C} { % T inséré après 2ème
    \pgfmathsetmacro{\x}{0.5 + \i*1.2}
    \fill[base\b] (\x, -1.15) rectangle (\x+0.4, -0.85);
    \node[white, font=\tiny\bfseries] at (\x+0.2, -1.0) {\b};
    \fill[base\cb] (\x, -1.3) rectangle (\x+0.4, -1.05);
    \node[white, font=\tiny\bfseries] at (\x+0.2, -1.175) {\cb};
}
% Marqueur
\draw[poppink, thick, decorate, decoration={snake, amplitude=1pt, segment length=4pt}] (4.1, -0.7) -- (4.1, -0.4);
\node[poppink, anchor=south, font=\tiny\bfseries] at (4.1, -0.4) {Inséré};
\end{tikzpicture}
\legende{Les trois types élémentaires de mutations géniques sur une courte séquence d'ADN double brin. En haut : séquence originale. Au centre : substitution d'une base (A$\to$G). À gauche : délétion d'une base (C). À droite : insertion d'une base (T). Les mutations par délétion ou insertion d'un nombre de bases non multiple de 3 entraînent une \textbf{mutation de décalage de cadre} (frameshift) qui modifie toute la séquence protéique en aval.}
\end{popfigure}

\begin{exemplebox}[La drépanocytose : une substitution ponctuelle classique]
La \textbf{drépanocytose} (anémie falciforme), très fréquente au Cameroun (prévalence de l'allèle $S$ : 10--25 \% selon les régions), résulte d'une \textbf{substitution} unique dans le gène $\beta$-globine (chromosome 11) :
\[
\text{Codon normal (GAG)} \xrightarrow{\text{substitution A$\to$T}} \text{Codon muté (GTG)}
\]
Cela remplace l'acide \textbf{glutamique} (chargé négativement, hydrophile) par de la \textbf{valine} (neutre, hydrophobe) en position 6 de la chaîne $\beta$. L'hémoglobine mutée (HbS) polymérise en conditions d'hypoxie, déformant les globules rouges en \emph{faucilles} (drepanocytes) $\to$ anémie hémolytique, crises vaso-occlusives, sensibilité aux infections.
\begin{itemize}
    \item Génotype $[A//A]$ : normal (HbAA).
    \item Génotype $[A//S]$ : hétérozygote \textbf{AS} (porteur sain, résistance au paludisme \emph{Plasmodium falciparum}).
    \item Génotype $[S//S]$ : drépanocytaire majeur (HbSS), pathologie sévère.
\end{itemize}
C'est un exemple de \textbf{codominance} au niveau moléculaire (les deux allèles s'expriment) et de \textbf{pléiotropie} (un gène, plusieurs effets phénotypiques).
\end{exemplebox}

\begin{attention}[Piège : « mutation = maladie » ?]
\textbf{Non.} La grande majorité des mutations sont \textbf{neutres} (régions non codantes, redondance du code génétique, substitution synonyme). Seule une minorité est délétère (maladies génétiques) ou, plus rarement, favorable (ex. : mutation $CCR5\Delta32$ conférant une résistance au VIH-1, mutation de la lactase permettant la digestion du lait à l'âge adulte). Sans mutations neutres ou favorables, l'évolution serait impossible.
\end{attention}

\courssub{Mutations chromosomiques : le gros lot}

\begin{propriete}[Mutations chromosomiques structurales]
Elles résultent de \textbf{cassures chromosomiques} suivies d'une réparation aberrante :
\begin{itemize}
    \item \textbf{Délétion} : perte d'un fragment (ex. : syndrome du cri du chat, délétion 5p15).
    \item \textbf{Duplication} : doublement d'un fragment (ex. : maladie de Charcot-Marie-Tooth type 1A, duplication 17p12).
    \item \textbf{Inversion} : fragment retourné de 180\textdegree{} (péri- ou paracentrique). Souvent sans effet phénotypique chez le porteur, mais risque de gamètes déséquilibrés à la méiose.
    \item \textbf{Translocation} : échange de fragments entre chromosomes non homologues. \emph{Réciproque équilibrée} : pas de perte de matériel, mais risque de trisomies partielles à la descendance. \emph{Robertsonienne} : fusion de deux chromosomes acrocentriques au niveau du centromère (ex. : translocation 14;21 $\to$ trisomie 21 par translocation, héréditaire).
\end{itemize}
\end{propriete}

\begin{propriete}[Mutations chromosomiques numériques]
\begin{itemize}
    \item \textbf{Aneuploïdie} : nombre de chromosomes non multiple du nombre haploïde ($n$). Principalement par \textbf{non-disjonction} (échec de séparation des chromatides sœurs ou des chromosomes homologues) au cours de la méiose I ou II, ou de la mitose précoce.
    \begin{itemize}
        \item Trisomie : $2n+1$ (ex. : trisomie 21, 18, 13 ; trisomies des chromosomes sexuels : XXX, XXY, XYY).
        \item Monosomie : $2n-1$ (viable seulement pour X : syndrome de Turner XO).
    \end{itemize}
    \item \textbf{Polyploïdie} : nombre de chromosomes multiple de $n$ ($3n$ triploïdie, $4n$ tétraploïdie). Fréquente chez les plantes (spéciation, agronomie : blé tendre hexaploïde, bananier triploïde stérile), létale chez l'homme (avortement spontané précoce).
\end{itemize}
\end{propriete}

\courssub{La trisomie 21 : modèle d'aneuploïdie humaine}

La trisomie 21 (syndrome de Down) est l'aneuploïdie viable la plus fréquente ($\approx$ 1/700 naissances vivantes). Elle résulte dans 95 \% des cas d'une \textbf{non-disjonction méiotique maternelle} (majoritairement en Méiose I).

\begin{popfigure}
\begin{tikzpicture}[scale=0.7, every node/.style={font=\tiny}]
% Caryotype simplifié : 23 paires, focus sur chr 21
\foreach \i/\x/\y in {
    1/0.5/5.5, 2/2.0/5.5, 3/3.5/5.5, 4/5.0/5.5, 5/6.5/5.5, 6/8.0/5.5, 7/9.5/5.5,
    8/0.5/4.0, 9/2.0/4.0, 10/3.5/4.0, 11/5.0/4.0, 12/6.5/4.0, 13/8.0/4.0, 14/9.5/4.0,
    15/0.5/2.5, 16/2.0/2.5, 17/3.5/2.5, 18/5.0/2.5, 19/6.5/2.5, 20/8.0/2.5, 21/9.5/2.5,
    22/0.5/1.0, 23/2.0/1.0
} {
    \pgfmathsetmacro{\w}{1.2}
    \pgfmathsetmacro{\h}{1.2}
    \pgfmathsetmacro{\cen}{\x+\w/2}
    \pgfmathsetmacro{\ceny}{\y+\h/2}
    \ifnum\i=21
        % Trois chromosomes 21
        \fill[poppink!30] (\x, \y) rectangle (\x+\w, \y+\h);
        \draw[poppink, thick] (\x, \y) rectangle (\x+\w, \y+\h);
        \node[poppink, font=\bfseries] at (\cen, \ceny) {21};
        \fill[poppink!30] (\x+1.4, \y) rectangle (\x+1.4+\w, \y+\h);
        \draw[poppink, thick] (\x+1.4, \y) rectangle (\x+1.4+\w, \y+\h);
        \node[poppink, font=\bfseries] at (\cen+1.4, \ceny) {21};
        \fill[poppink!30] (\x+2.8, \y) rectangle (\x+2.8+\w, \y+\h);
        \draw[poppink, thick] (\x+2.8, \y) rectangle (\x+2.8+\w, \y+\h);
        \node[poppink, font=\bfseries] at (\cen+2.8, \ceny) {21};
        \node[poppink, anchor=north, font=\bfseries] at (\cen+1.4, \y-0.3) {Trisomie 21};
    \else
        \ifnum\i=23
            % Chromosomes sexuels XY
            \fill[popblue!20] (\x, \y) rectangle (\x+\w, \y+\h);
            \draw[popblue] (\x, \y) rectangle (\x+\w, \y+\h);
            \node[popblue, font=\bfseries] at (\cen, \ceny) {X};
            \fill[popblue!20] (\x+1.4, \y) rectangle (\x+1.4+\w, \y+\h);
            \draw[popblue] (\x+1.4, \y) rectangle (\x+1.4+\w, \y+\h);
            \node[popblue, font=\bfseries] at (\cen+1.4, \ceny) {Y};
        \else
            \fill[popmuted!20] (\x, \y) rectangle (\x+\w, \y+\h);
            \draw[popmuted] (\x, \y) rectangle (\x+\w, \y+\h);
            \node[popdark, font=\bfseries] at (\cen, \ceny) {\i};
        \fi
    \fi
}

% Flèche indicatrice
\draw[poppink, thick, ->] (12, 3.0) -- (10.5, 3.0);
\node[poppink, anchor=west, font=\small\bfseries] at (12, 3.0) {3 chromosomes 21 (au lieu de 2)};
\end{tikzpicture}
\legende{Caryotype schématisé d'un individu trisomique 21 (47,XX,+21 ou 47,XY,+21). Les 22 paires d'autosomes et la paire de gonosomes sont représentées. La trisomie du chromosome 21 (trois copies au lieu de deux) est mise en évidence en rose. Le chromosome 21 est le plus petit autosome humain (environ 48 Mb, $\approx$ 1,5 \% de l'ADN total).}
\end{popfigure}

\begin{experience}[Mécanisme de la non-disjonction méiotique (modèle admis)]
\textbf{Observation :} Un gamète porte deux chromosomes 21 au lieu d'un.
\textbf{Hypothèse :} Échec de la séparation des chromosomes homologues (Méiose I) ou des chromatides sœurs (Méiose II).
\textbf{Interprétation (schéma) :}
\begin{itemize}
    \item \textbf{Méiose I :} Les bivalents 21 ne se séparent pas $\to$ une cellule fille reçoit les deux chromosomes 21 (chacun avec deux chromatides), l'autre n'en reçoit aucun. Après Méiose II : deux gamètes $n+1$ (deux chr 21) et deux gamètes $n-1$ (zéro chr 21).
    \item \textbf{Méiose II :} Séparation normale en MI, mais en MII les chromatides sœurs d'un chromosome 21 ne se séparent pas $\to$ un gamète $n+1$ (deux chromatides = un chromosome double), un gamète $n-1$, deux gamètes $n$ normaux.
\end{itemize}
\textbf{Conséquence :} Fécondation d'un gamète $n+1$ (21,21) par un gamète normal $n$ (21) $\to$ zygote $2n+1$ (trisomie 21). Fécondation d'un gamète $n-1$ $\to$ monosomie 21 (non viable, avortement précoce).
\end{experience}

\begin{savaistu}[Âge maternel et trisomie 21]
La fréquence de la trisomie 21 à la naissance augmente fortement avec l'âge de la mère au moment de la conception :
\begin{center}
\begin{tabular}{|c|c|c|}\hline
\textbf{Âge maternel} & \textbf{Risque approximatif à la naissance} & \textbf{Risque à 12 SA (amniocentèse)} \\ \hline
20 ans & 1 / 1\,500 & 1 / 1\,000 \\ \hline
30 ans & 1 / 900 & 1 / 600 \\ \hline
35 ans & 1 / 350 & 1 / 200 \\ \hline
40 ans & 1 / 100 & 1 / 65 \\ \hline
45 ans & 1 / 30 & 1 / 20 \\ \hline
\end{tabular}
\end{center}
\textbf{Explication :} Les ovocytes sont bloqués en prophase I (dictyotène) depuis la vie fœtale jusqu'à l'ovulation. Avec l'âge, la cohésion des chromatides sœurs (protéines cohesines) se dégrade, augmentant le risque de non-disjonction. Au Cameroun, le dépistage combiné (marqueurs sériques PAPP-A, $\beta$-hCG libre + clarté nucale à 11--13 SA) est recommandé pour les grossesses à risque.
\end{savaistu}

\courssub{Origine des mutations : spontanées et induites}

\begin{propriete}[Agents mutagènes]
\begin{itemize}
    \item \textbf{Mutations spontanées :} Erreurs de réplication de l'ADN (tautomérie des bases), dépurination, déamination spontanée (C$\to$U), instabilité des répétitions microsatellites. Taux de base $\approx 10^{-8}$ à $10^{-9}$ par nucléotide par génération.
    \item \textbf{Mutations induites :} Causées par des \textbf{agents mutagènes} physiques ou chimiques :
    \begin{itemize}
        \item \textbf{Rayonnements ionisants} (X, $\gamma$, UV) : cassures double brin, dimères de thymine (UV).
        \item \textbf{Chimiques :} Agents alkylants (moutarde azotée, EMS), analogues de bases (5-bromouracile), intercalants (bromure d'éthidium), métaux lourds, hydrocarbures aromatiques (tabac, fumées de cuisine au bois/charbon -- contexte camerounais pertinent).
    \end{itemize}
\end{itemize}
\end{propriete}

\begin{attention}[Mutagène $\neq$ Cancérigène $\neq$ Tératogène]
Un \textbf{mutagène} altère l'ADN (test d'Ames sur \emph{Salmonella} souche his-). Un \textbf{cancérigène} provoque le cancer (souvent mutagène, mais pas toujours : promoteurs tumoraux comme le TPA). Un \textbf{tératogène} perturbe le développement embryonnaire (ex. : alcool, thalidomide, virus rubéole, acide valproïque) sans nécessairement muter l'ADN. Ces notions sont distinctes mais peuvent se chevaucher.
\end{attention}

\courssub{Conséquences des mutations et rôle dans l'évolution}

\begin{aretenir}[Devenir des mutations]
\begin{center}
\begin{tabularx}{\linewidth}{|l|X|X|}\hline
\textbf{Type d'effet} & \textbf{Fréquence} & \textbf{Devenir évolutif} \\ \hline
Létale (embryonnaire) & Fréquente & Éliminée immédiatement (sélection purifiante) \\ \hline
Délétère (maladie) & Moyenne & Maintenue à basse fréquence (équilibre mutation-sélection) ou avantage hétérozygote (drépanocytose/paludisme) \\ \hline
Neutre (synonyme, non codant) & Majoritaire & Dérive génétique, horloge moléculaire \\ \hline
Favorable (avantage sélectif) & Rare & Sélection positive, fixation, adaptation \\ \hline
\end{tabularx}
\end{center}
\end{aretenir}

L'exemple de la drépanocytose illustre parfaitement l'\textbf{équilibre mutation-sélection} avec \textbf{avantage hétérozygote} : dans les zones d'endémie palustre (forêt équatoriale du Sud-Cameroun, bassin du lac Tchad), la fréquence de l'allèle $S$ se maintient autour de 10--15 \% car les AS résistent mieux au paludisme grave que les AA, tandis que les SS meurent jeunes. C'est la \textbf{polymorphisme équilibré} classique.

\courssub{Calcul du risque génétique : prévoir pour décider}

Le \textbf{risque génétique} est la probabilité qu'un couple ait un enfant atteint d'une affection génétique donnée. Il combine :
\begin{enumerate}
    \item La probabilité des génotypes des parents (déduite des pedigrees, tests, population).
    \item La probabilité de transmission des allèles pathogènes (lois de Mendel).
    \item La pénétrance (pourcentage de génotype muté qui s'exprime phénotypiquement) et l'expressivité.
\end{enumerate}

\begin{methode}[Calculer un risque génétique (approche par arbre de probabilités)]
\begin{enumerate}
    \item \textbf{Identifier le mode de transmission} (autosomal récessif/dominant, lié au X récessif/dominant) et la pénétrance.
    \item \textbf{Déterminer les probabilités des génotypes parentaux} :
    \begin{itemize}
        \item Si phénotype connu $\to$ génotype certain ou probabiliste (ex. : phénotype normal, parents d'un enfant atteint $\to$ hétérozygotes obligés pour récessif).
        \item Si antécédents familiaux (pedigree) $\to$ calcul bayésien (probabilité conditionnelle).
        \item Si population générale $\to$ fréquence allélique (Hardy-Weinberg : $q^2$ malades, $2pq$ porteurs).
    \end{itemize}
    \item \textbf{Construire l'arbre des probabilités} : croiser les gamètes parentaux (tableau de Punnett ou arbre).
    \item \textbf{Multiplier les probabilités} le long de chaque branche menant au génotype « malade ».
    \item \textbf{Sommer} les probabilités de toutes les branches « malade ».
    \item \textbf{Exprimer} en pourcentage ou fraction (1 sur $n$).
\end{enumerate}
\end{methode}

\begin{popfigure}
\begin{tikzpicture}[scale=0.8, every node/.style={font=\small},
    level 1/.style={sibling distance=4cm, level distance=2.5cm},
    level 2/.style={sibling distance=2cm, level distance=2.5cm},
    edge from parent/.style={draw, thick, ->},
    bag/.style={circle, draw, thick, minimum width=1.2cm, inner sep=2pt},
    leaf/.style={rectangle, draw, rounded corners, fill=popgreenL, thick, inner sep=3pt}
]
\node[bag] {Couple AS $\times$ AS}
    child {node[bag] {Gamète mère}
        child {node[leaf] {$A$ (0,5)}}
        child {node[leaf] {$S$ (0,5)}}
    }
    child {node[bag] {Gamète père}
        child {node[leaf] {$A$ (0,5)}}
        child {node[leaf] {$S$ (0,5)}}
    };
% Annotations des résultats
\node[anchor=west, font=\small] at (5.5, 3.5) {\textbf{Zygotes possibles :}};
\node[anchor=west, font=\small] at (5.5, 3.0) {$AA$ : $0,5 \times 0,5 = 0,25$ (normal)};
\node[anchor=west, font=\small] at (5.5, 2.5) {$AS$ : $0,5 \times 0,5 + 0,5 \times 0,5 = 0,50$ (porteur sain)};
\node[anchor=west, font=\small, text=poppink] at (5.5, 2.0) {$SS$ : $0,5 \times 0,5 = \mathbf{0,25}$ \textbf{(drépanocytaire)}};
\node[anchor=west, font=\small, text=popblue] at (5.5, 1.3) {\textbf{Risque de récurrence = 25 \% par grossesse}};
\end{tikzpicture}
\legende{Arbre de probabilités pour le croisement de deux hétérozygotes AS (drépanocytose, autosomique récessif). Chaque gamète a une probabilité de 0,5 de porter l'allèle $A$ ou $S$. La probabilité conjointe d'obtenir un zygote $SS$ est $0,5 \times 0,5 = 0,25$. Ce risque est \textbf{indépendant} pour chaque grossesse (événements indépendants).}
\end{popfigure}

\begin{exoresolu}[Risque génétique avec incertitude sur le génotype parental]
\textbf{Énoncé :} Un couple consanguin (cousins germains) consulte. L'homme a un oncle maternel atteint de drépanocytose (SS). La femme est de population générale (fréquence allèle $S = 0,12$). Quel est le risque qu'ils aient un enfant drépanocytaire (SS) ?
\textbf{Solution détaillée :}
\begin{enumerate}
    \item \textbf{Homme (cousin de l'oncle atteint) :}
    \begin{itemize}
        \item L'oncle est SS $\to$ les grands-parents maternels sont obligatoirement AS $\times$ AS.
        \item La mère de l'homme (sœur de l'oncle) a un phénotype normal. Parmi les descendants normaux d'un croisement AS $\times$ AS, les génotypes sont $AA$ (1/3) et $AS$ (2/3). \textbf{Probabilité que la mère soit porteuse $AS = 2/3$}.
        \item Si la mère est $AS$, elle transmet $S$ à son fils (l'homme) avec une probabilité 1/2.
        \item $\to$ Probabilité que l'homme soit porteur $AS = \frac{2}{3} \times \frac{1}{2} = \frac{1}{3}$.
    \end{itemize}
    \item \textbf{Femme (population générale) :} Fréquence porteurs $2pq \approx 2 \times 0,88 \times 0,12 = 0,2112 \approx 21,1 \%$.
    \item \textbf{Risque enfant SS :} Nécessite que les deux parents soient porteurs ET transmettent $S$.
    \[
    P(\text{enfant SS}) = P(\text{père } AS) \times P(\text{mère } AS) \times P(S \text{ de père}) \times P(S \text{ de mère})
    \]
    \[
    = \frac{1}{3} \times 0,2112 \times \frac{1}{2} \times \frac{1}{2} = \frac{0,2112}{12} = 0,0176 = \mathbf{1,76 \%}
    \]
    \item \textbf{Conclusion :} Risque $\approx 1/57$. Sans consanguinité (père population générale), risque serait $0,2112 \times 0,2112 \times 1/4 \approx 1,1 \%$. La consanguinité augmente le risque car elle augmente la probabilité de partager l'allèle ancestral.
\end{enumerate}
\end{exoresolu}

\begin{exemplebox}[Deux parents hétérozygotes AS]
Pour un couple où les deux partenaires sont \textbf{porteurs sains AS} (dépistés ou ayant eu un enfant atteint), le risque d'avoir un enfant drépanocytaire majeur (SS) à \textbf{chaque grossesse} est de :
\[
P(SS) = \frac{1}{2} \times \frac{1}{2} = \frac{1}{4} = \mathbf{25 \%}
\]
Ce risque est \textbf{constant} d'une grossesse à l'autre (indépendance des événements méiotiques). Avoir un enfant sain ne « protège » pas le suivant, et avoir un enfant atteint n'augmente pas le risque pour le suivant.
\end{exemplebox}

\courssub{Prévention des anomalies génétiques : du dépistage au conseil}

La prévention en génétique humaine repose sur trois piliers : \textbf{information}, \textbf{dépistage}, \textbf{accompagnement}.

\begin{propriete}[Stratégies de prévention au Cameroun]
\begin{itemize}
    \item \textbf{Conseil génétique (consultation spécialisée) :} Disponible dans les grands hôpitaux (CHU Yaoundé, Douala, Garoua, centres de référence drépanocytose). Objectif : informer le couple sur le risque, le mode de transmission, les options (diagnostic prénatal, adoption, don de gamètes, FIV avec DPI), sans directivité. Respect de l'autonomie et de la confidentialité.
    \item \textbf{Dépistage néonatal systématique de la drépanocytose :} Prélèvement sur buvard (carte Guthrie) à la naissance ou à la première consultation vaccinale (PEV) pour électrophorèse d'hémoglobine. Permet une prise en charge précoce (prophylaxie pénicilline, vaccination anti-pneumocoque, acide folique, éducation des parents) $\to$ réduction drastique de la mortalité infantile.
    \item \textbf{Dépistage prénuptial / préconceptionnel :} Électrophorèse d'hémoglobine proposée aux futurs mariés (sensibilisation dans les églises, mairies, médias). Éviter le mariage AS $\times$ AS (risque 25 \% SS) ou AS $\times$ SS (risque 50 \% SS) par choix éclairé. \emph{Attention :} ne pas stigmatiser les porteurs ; le conseil prime sur l'interdiction.
    \item \textbf{Incompatibilité Rhésus (Rh) :} Dépistage Rh dès la 1\textsuperscript{re} consultation prénatale (CPN). Femme Rh- (phénotype $rr$, fréquence faible $\approx$ 1--5 \% au Cameroun) avec conjoint Rh+ $\to$ risque d'allo-immunisation anti-D si fœtus Rh+. \textbf{Prophylaxie :} Injection d'immunoglobulines anti-D (RhoGAM) à 28 SA et dans les 72 h post-partum (ou après tout événement sensibilisant : IVG, amniocentèse, traumatisme). Gratuit dans le secteur public dans de nombreux districts.
    \item \textbf{Diagnostic prénatal (DPN) :}
    \begin{itemize}
        \item \textbf{Non invasif (DPNI) :} ADN fœtal libre dans le sang maternel (dès 10 SA). Dépistage trisomies 21, 18, 13, aneuploïdies sexuelles. Sensibilité $> 99 \%$ pour T21. Pas de risque fœtal. Coût encore élevé.
        \item \textbf{Invasif :} Amniocentèse (15--18 SA, liquide amniotique), biopsie chorialaire (11--13 SA, trophoblaste). Caryotypage (bande G, FISH, puces à ADN / CGH-array). Risque de fausse couche $\approx 0,5--1 \%$. Indiqué si DPNI positif, antécédents, âge maternel $\ge 38$ ans, anomalie échographique.
    \end{itemize}
    \item \textbf{Évitement des mariages consanguins à risque :} Sensibilisation sur l'augmentation du risque de maladies autosomiques récessives rares (ex. : albinisme, mucoviscidose, surdité génétique) chez les cousins (partagent 1/8 de leurs allèles identiques par descendance). Risque multiplié par 2 à 10 selon la pathologie.
    \item \textbf{Supplémentation en acide folique (vitamine B9) :} 0,4 mg/jour dès le désir de grossesse (ou 5 mg/jour si antécédent de DFTN) $\to$ prévention des \textbf{défauts de fermeture du tube neural} (anencéphalie, spina bifida), anomalies multifactorielles (gènes + environnement).
\end{itemize}
\end{propriete}

\begin{savaistu}[Drépanocytose et politique de santé au Cameroun]
Le Cameroun a adopté un \textbf{Plan Stratégique National de Lutte contre la Drépanocytose (2020--2024)}. Les axes majeurs : (1) Dépistage néonatal systématique dans les formations sanitaires de premier niveau ; (2) Création de centres de prise en charge intégrés (Yaoundé, Douala, Garoua, Bafoussam, Bertoua, Ngaoundéré, Maroua) ; (3) Subvention des traitements (hydroxycarbamide, vaccins, prophylaxie antipaludéenne) ; (4) Sensibilisation communautaire via les leaders d'opinion, les médias, les écoles. Le but : réduire la mortalité des enfants drépanocytaires de 50 \% d'ici 2025. La génétique n'est pas une fatalité : la connaissance permet l'action.
\end{savaistu}

\begin{attention}[Erreurs fréquentes en calcul de risque -- Vigilance Bac !]
\begin{itemize}
    \item \textbf{Confondre risque de récurrence et risque d'occurrence :} Le risque de récurrence (couple déjà parents d'un enfant atteint) est souvent plus élevé (ex. : 25 \% pour AR) que le risque d'occurrence dans la population générale (ex. : $q^2$).
    \item \textbf{Oublier la pénétrance incomplète :} Si une maladie dominante a une pénétrance de 80 \%, un parent porteur de la mutation a 80 \% de chance de l'exprimer. Un parent \textbf{phénotypiquement normal} peut quand même être porteur (probabilité $= \frac{0,2}{1,2} = 1/6$ si l'autre parent est mutant... calcul bayésien !).
    \item \textbf{Traiter les grossesses comme dépendantes :} « Ils ont déjà 3 enfants sains, le 4\textsuperscript{e} a plus de chances d'être malade » $\to$ \textbf{FAUX} (indépendance des méioses, loi des grands nombres).
    \item \textbf{Ne pas distinguer allèle mutant et allèle normal dans le tableau de Punnett :} Toujours écrire les allèles ($A$, $a$ ou $X^H$, $X^h$) et pas seulement les phénotypes.
    \item \textbf{Omettre l'unité ou le \% :} « Risque = 1/4 » $\to$ écrire « Risque = 1/4 = 25 \% = 1 sur 4 ».
\end{itemize}
\end{attention}

\begin{exercice}[Entraînement : Risque hémophilie A]
Un couple consulte. La femme a un frère hémophile A (maladie récessive liée à l'X). Son père est sain. La femme est phénotypiquement normale. L'homme est sain. Calculer le risque qu'ils aient un fils hémophile. (On supposera que la mutation est rare dans la population).
\end{exercice}


\begin{corrige}[Exercice n]
\begin{enumerate}
    \item \textbf{Mode de transmission :} Récessif lié à l'X ($X^h$ = allèle hémophilie, $X^H$ = normal).
    \item \textbf{Génotype du père de la femme :} Sain $\to$ $X^H Y$.
    \item \textbf{Génotype de la mère de la femme :} A un fils hémophile ($X^h Y$) $\to$ elle a transmis $X^h$ $\to$ elle est obligatoirement porteuse $X^H X^h$.
    \item \textbf{Génotype de la femme (consultante) :} Elle a reçu $X^H$ de son père. De sa mère ($X^H X^h$), elle a 1/2 chance de recevoir $X^H$ (génotype $X^H X^H$) et 1/2 chance de recevoir $X^h$ (génotype $X^H X^h$ porteuse).
    $\to P(\text{femme porteuse}) = 1/2$.
    \item \textbf{Homme :} Sain $\to$ $X^H Y$.
    \item \textbf{Risque fils hémophile :} Nécessite : (femme porteuse) ET (transmission $X^h$ à un fils) ET (sexe masculin).
    \[
    P = P(\text{femme } X^H X^h) \times P(\text{transmet } X^h) \times P(\text{fils}) = \frac{1}{2} \times \frac{1}{2} \times \frac{1}{2} = \frac{1}{8} = \mathbf{12,5 \%}
    \]
    \item \textbf{Risque fille hémophile :} 0 (père donne $X^H$).
\end{enumerate}
\end{corrige}

\begin{aretenir}[Synthèse : Mutations, risques, prévention]
\begin{itemize}
    \item La \textbf{mutation} (génique ou chromosomique) est la source de tout allèle nouveau. Elle peut être spontanée ou induite (mutagènes : UV, chimiques, radiations).
    \item \textbf{Mutation génique :} substitution, délétion, insertion. Ex. drépanocytose (substitution $\beta^6$ Glu$\to$Val).
    \item \textbf{Mutation chromosomique :} structurelle (délétion, duplication, inversion, translocation) ou numérique (aneuploïdie : trisomie 21 par non-disjonction méiotique maternelle ; polyploïdie).
    \item \textbf{Risque génétique :} Probabilité conditionnelle combinant génotypes parentaux (déduits du pedigree, tests, fréquence population) et lois de Mendel. Calcul par arbre de probabilités.
    \item \textbf{Prévention :} Conseil génétique (non directif), dépistage néonatal (drépanocytose), prénuptial (hémoglobine, Rh), prénatal (DPNI, amniocentèse), supplémentation folique, information sur la consanguinité.
    \item \textbf{Contexte camerounais :} Drépanocytose (prévalence élevée, plan national), incompatibilité Rh (prophylaxie anti-D), consanguinité (zones culturelles), défauts tube neural (acide folique).
\end{itemize}
\end{aretenir}

\newpage

\coursec{Activité d'intégration}

\coursec{Leçon 04 : La transmission des caractères au sein d'une espèce et quelques aspects de la génétique humaine}

\courssub{Activité d'intégration : Conseil génétique au Centre Hospitalier Essos (Yaoundé)}

\courssub{Objectifs de l'activité}
\begin{itemize}
    \item Mobiliser les connaissances sur la génétique formelle (monohybridisme, dihybridisme, brassage interchromosomique) et la génétique humaine (hérédité autosomique, groupes sanguins) pour analyser une situation réelle de conseil génétique.
    \item Construire et exploiter un pedigree complexe portant sur deux caractères autosomiques (albinisme et drépanocytose) et les systèmes ABO/Rhésus.
    \item Calculer des risques génétiques précis (probabilités conditionnelles et fréquences alléliques populationnelles) pour une descendance donnée.
    \item Rédiger un avis de conseil génétique structuré, argumenté scientifiquement et adapté au contexte socioculturel camerounais, intégrant les dimensions préventives et éthiques.
\end{itemize}

\begin{experience}[Situation clinique : Consultation de génétique prénuptiale]
M. \textbf{Jean Ngo Bell} (32 ans) et Mme \textbf{Marie Ngo Bell, née Essomba} (28 ans), couple résidant à Yaoundé, se présentent à la consultation de génétique médicale du Centre Hospitalier Essos pour un \emph{conseil prénuptial et prénatal}. Ils souhaitent avoir des enfants mais s'inquiètent de l'existence de maladies génétiques dans leurs familles respectives.

\textbf{Antécédents familiaux (recueil à l'anamnèse) :}
\begin{itemize}
    \item \textbf{Côté Ngo Bell (branche paternelle) :} Le grand-père paternel de Jean (I-1) était albinos. Le père de Jean (II-1, vivant, 60 ans) est phénotypiquement normal. Jean a une sœur aînée, \textbf{Anne} (III-2), albinos. Jean lui-même (III-1) est phénotypiquement normal (peau pigmentée). Aucun cas de drépanocytose n'est signalé dans cette branche.
    \item \textbf{Côté Essomba (branche maternelle) :} La mère de Marie (II-3, vivante, 55 ans) est drépanocytaire homozygote (SS). Le père de Marie (II-4, décédé) était phénotypiquement normal (AS ou AA). Marie a un frère cadet, \textbf{Paul} (III-4), drépanocytaire (SS). Marie elle-même (III-3) est phénotypiquement normale (pas de crises vaso-occlusives). Aucun cas d'albinisme n'est signalé dans cette branche.
    \item \textbf{Consanguinité :} Aucune consanguinité connue entre les deux familles.
\end{itemize}

\textbf{Examens biologiques du couple (résultats du laboratoire d'hématologie) :}
\begin{center}
\begin{tabularx}{\linewidth}{|l|X|X|}\hline
\rowcolor{popblueL} \textbf{Paramètre} & \textbf{M. Jean Ngo Bell} & \textbf{Mme Marie Ngo Bell} \\ \hline
Groupe sanguin ABO & A & B \\
Rhésus (D) & Positif (Rh+) & Négatif (Rh-) \\
Électrophorèse de l'hémoglobine (Hb) & AA (normale) & AS (hétérozygote drépanocytaire) \\ \hline
\end{tabularx}
\end{center}

\textbf{Question médicale posée :} Quel est le risque pour ce couple d'avoir un enfant atteint de drépanocytose (homozigote SS) ? D'avoir un enfant albinos ? D'avoir un enfant atteint des deux affections simultanément ? Quelles mesures de prévention et de prise en charge proposer ?
\end{experience}

\courssub{Consignes}
\begin{enumerate}
    \item \textbf{Analyse des modes de transmission :} À partir des informations du pedigree (à construire), déterminer le mode de transmission de l'albinisme et de la drépanocytose. Justifier en précisant les allèles en jeu ($A/a$ pour l'albinisme, $S^A/S^S$ pour l'hémoglobine) et les génotypes probables des individus clés (grands-parents, parents, fratrie, consultants).
    \item \textbf{Construction du pedigree :} Représenter le pedigree complet sur 4 générations (Grands-parents $\to$ Parents $\to$ Consultants/Fratrie $\to$ Descendance théorique) en utilisant les symboles standards (carré = homme, cercle = femme, symbole plein = atteint, symbole moitié plein = porteur hétérozygote connu, flèche = consultant). Indiquer les génotypes déduits sous chaque symbole.
    \item \textbf{Génotypes des consultants :} Déduire avec certitude les génotypes complets de Jean et Marie pour les trois systèmes étudiés : Albinisme ($A/a$), Drépanocytose ($S^A/S^S$), Groupes sanguins ABO ($I^A, I^B, i$) et Rhésus ($Rh^+/Rh^-$). Préciser le lien de liaison (indépendance) entre ces gènes.
    \item \textbf{Croisements et risques génétiques :} 
    \begin{enumerate}
        \item Écrire le croisement pour le gène de l'hémoglobine (drépanocytose) et calculer la probabilité d'avoir un enfant SS.
        \item Écrire le croisement pour le gène de l'albinisme et calculer la probabilité d'avoir un enfant albinos ($aa$).
        \item Calculer la probabilité d'avoir un enfant \textbf{à la fois} drépanocytaire (SS) \textbf{et} albinos ($aa$). Justifier le calcul (indépendance des gènes).
        \item Déterminer les phénotypes ABO et Rhésus possibles pour la descendance et leurs fréquences.
    \end{enumerate}
    \item \textbf{Avis de conseil génétique :} Rédiger un compte-rendu de consultation structuré (Anamnèse, Analyse génétique, Risques calculés, Recommandations/Prévention, Suivi) adressé au couple, en langage clair mais rigoureux.
\end{enumerate}

\courssub{Ressources à mobiliser}
\begin{propriete}[Rappels génétiques indispensables]
\textbf{1. Monohybridisme autosomique (Dominance complète) :} Allèles $A$ (dominant) et $a$ (récessif). Croisement $Aa \times Aa \rightarrow$ Phénotypes : $\frac{3}{4}\ A-$ (normal), $\frac{1}{4}\ aa$ (atteint). Un phénotype normal issu de ce croisement a $\frac{2}{3}$ de chances d'être $Aa$ et $\frac{1}{3}$ d'être $AA$.\\[4pt]
\textbf{2. Codominance (Groupe sanguin ABO) :} Allèles $I^A, I^B$ codominants, $i$ récessif. Génotypes : $[I^A I^A]$ ou $[I^A i]$ (Gpe A), $[I^B I^B]$ ou $[I^B i]$ (Gpe B), $[I^A I^B]$ (Gpe AB), $[ii]$ (Gpe O).\\[4pt]
\textbf{3. Système Rhésus (Facteur D) :} $Rh^+$ dominant sur $Rh^-$. Phénotype Rh- $\equiv$ génotype $[Rh^-//Rh^-]$.\\[4pt]
\textbf{4. Drépanocytose (Codominance au niveau moléculaire/phénotypique) :} $S^A S^A$ (AA, normal), $S^A S^S$ (AS, hétérozygote, trait drépanocytaire, phénotype normal en conditions ordinaires), $S^S S^S$ (SS, drépanocytose majeure). L'électrophorèse distingue les trois génotypes.\\[4pt]
\textbf{5. Indépendance des gènes (Brassage interchromosomique) :} Gènes sur chromosomes différents (ou lointains sur le même chromosome). $P(G_1 \cap G_2) = P(G_1) \times P(G_2)$. Localisations : \emph{HBB} (Chr 11), \emph{OCA2} (Chr 15), \emph{ABO} (Chr 9), \emph{RHD/RHCE} (Chr 1).
\end{propriete}

\begin{methode}[Calcul de risque conditionnel dans un pedigree]
\begin{enumerate}
    \item Identifier le mode de transmission (autosomique récessif pour les deux maladies ici).
    \item Déterminer le génotype des parents des consultants (obligatoirement hétérozygotes s'ils ont un enfant atteint et sont phénotypiquement normaux).
    \item Calculer la probabilité \emph{a priori} pour le consultant d'être hétérozygote (ex: $\frac{2}{3}$ pour un phénotype normal issu de $Aa \times Aa$).
    \item \textbf{Intégrer l'information biologique directe :} Si un test (électrophorèse, génotypage) confirme le statut (ex: Jean AA pour Hb, Marie AS), la probabilité devient 1 (certitude) pour ce gène. Le risque ne dépend plus du pedigree mais du croisement direct des génotypes connus.
    \item Pour un conjoint non informé par le pedigree (ex: Marie pour l'albinisme), utiliser la fréquence allélique populationnelle $q$ : $P(\text{porteur}) \approx 2q$ (si $q$ faible).
    \item Multiplier les probabilités indépendantes pour les caractères non liés (loi du produit).
\end{enumerate}
\end{methode}

\begin{attention}[Pièges fréquents à éviter]
\begin{itemize}
    \item Confondre \emph{phénotype normal} et \emph{génotype homozygote dominant}. Un individu sain issu de parents hétérozygotes a $\frac{2}{3}$ de chances d'être hétérozygote (pas $\frac{1}{2}$).
    \item Oublier que l'électrophorèse de l'hémoglobine distingue AA, AS et SS (codominance), ce qui lève l'ambiguïté sur le génotype des consultants pour la drépanocytose.
    \item Traiter les gènes ABO, Rhésus, Albinisme, Drépanocytose comme liés alors qu'ils sont sur des chromosomes différents $\rightarrow$ Indépendance stricte (brassage interchromosomique).
    \item Négliger le risque Rhésus (allo-immunisation maternelle anti-D) dans le conseil prénatal si mère Rh- et père Rh+.
    \item Utiliser le point décimal dans le code TikZ/pgfplots (ex: \texttt{domain=0:1.5}) et la virgule dans le texte (ex: \texttt{1,5}).
\end{itemize}
\end{attention}

\courssub{Démarche et corrigé commenté}
\begin{exoresolu}[Conseil génétique couple Ngo Bell : Résolution complète]
\textbf{Étape 1 : Construction du pedigree et déduction des génotypes}

\begin{popfigure}
\begin{tikzpicture}[
    scale=0.75, transform shape,
    male/.style={rectangle, draw=popdark, thick, minimum size=1.2cm},
    female/.style={circle, draw=popdark, thick, minimum size=1.2cm},
    affected/.style={fill=popred!60},
    carrier/.style={fill=poporange!40, pattern=north east lines, pattern color=poporange},
    normal/.style={fill=popgreen!20},
    consultant/.style={draw=popblue, ultra thick},
    marriage/.style={thick},
    desc/.style={thick}
]
% Generation I
\node[male, affected, label=below:I-1] (I1) at (0,0) {$aa$};
\node[female, carrier, label=below:I-2] (I2) at (3,0) {$Aa$};
\node[female, affected, label=below:I-3, xshift=8cm] (I3) at (8,0) {$S^S S^S$};
\node[male, carrier, label=below:I-4] (I4) at (11,0) {$S^A S^S$};

% Marriage lines Gen I
\draw[marriage] (I1) -- (I2);
\draw[marriage] (I3) -- (I4);

% Generation II
\node[male, carrier, label=below:II-1] (II1) at (1.5,-2) {$Aa$};
\node[female, carrier, label=below:II-2] (II2) at (1.5,-3.5) {$Aa$};
\node[female, affected, label=below:II-3] (II3) at (9.5,-2) {$S^S S^S$};
\node[male, carrier, label=below:II-4] (II4) at (9.5,-3.5) {$S^A S^S$};

% Descent lines Gen I -> II
\draw[desc] (1.5,0) -- (1.5,-2);
\draw[desc] (9.5,0) -- (9.5,-2);

% Marriage lines Gen II
\draw[marriage] (II1) -- (II2);
\draw[marriage] (II3) -- (II4);

% Generation III
\node[male, normal, consultant, label=below:{III-1 Jean}] (III1) at (0.5,-5.5) {$AA$};
\node[female, affected, label=below:{III-2 Anne}] (III2) at (2.5,-5.5) {$aa$};
\node[female, carrier, consultant, label=below:{III-3 Marie}] (III3) at (8.5,-5.5) {$S^A S^S$};
\node[male, affected, label=below:{III-4 Paul}] (III4) at (10.5,-5.5) {$S^S S^S$};

% Descent lines Gen II -> III
\draw[desc] (1.5,-3.5) -- (0.5,-5.5);
\draw[desc] (1.5,-3.5) -- (2.5,-5.5);
\draw[desc] (9.5,-3.5) -- (8.5,-5.5);
\draw[desc] (9.5,-3.5) -- (10.5,-5.5);

% Generation IV (Theoretical)
\node[draw=popmuted, dashed, label=below:IV-1?] (IV1) at (0.5,-7.5) {?};
\node[draw=popmuted, dashed, label=below:IV-2?] (IV2) at (2.5,-7.5) {?};
\draw[desc, dashed] (III1) -- (IV1);
\draw[desc, dashed] (III3) -- (IV1);
\draw[desc, dashed] (III1) -- (IV2);
\draw[desc, dashed] (III3) -- (IV2);

% Legend
\begin{scope}[shift={(13,-1)}]
\node[male, normal, minimum size=0.8cm] (l1) at (0,0) {}; \node[right] at (1,0) {Homme sain};
\node[female, affected, minimum size=0.8cm] (l2) at (0,-1.2) {}; \node[right] at (1,-1.2) {Femme atteinte};
\node[male, carrier, minimum size=0.8cm] (l3) at (0,-2.4) {}; \node[right] at (1,-2.4) {Homme porteur (hétérozygote)};
\node[male, consultant, minimum size=0.8cm] (l4) at (0,-3.6) {}; \node[right] at (1,-3.6) {Consultant (proband)};
\node[draw=popmuted, dashed, minimum size=0.8cm] (l5) at (0,-4.8) {}; \node[right] at (1,-4.8) {Descendance théorique};
\end{scope}
\end{tikzpicture}
\legende{Pedigree simplifié de la famille Ngo Bell / Essomba. Génération I (grands-parents), II (parents), III (consultants et fratrie), IV (descendance théorique). Symboles : carré = homme, cercle = femme. Rempli rouge = atteint (aa ou SS). Hachuré orange = porteur sain certain (Aa ou AS). Vert clair = phénotype normal (génotype AA ou Aa selon déduction). Épais bleu = consultants. Pointillés = génération future.}
\end{popfigure}

\medskip
\noindent \textbf{A. Albinisme (Gène \emph{OCA2}, Chr 15, autosomique récessif). Allèles : $A$ (normal), $a$ (albinos).}
\begin{itemize}
    \item \textbf{Génération I :} Grand-père paternel (I-1) albinos $\rightarrow$ $[a//a]$. Grand-mère paternelle (I-2) phénotype normal, a eu un enfant albinos (II-1) et une petite-fille albinos (III-2) $\rightarrow$ obligatoirement hétérozygote $[A//a]$.
    \item \textbf{Génération II :} Père de Jean (II-1) : phénotype normal, né de $[a//a] \times [A//a]$ $\rightarrow$ Génotype certain $[A//a]$. Mère de Jean (II-2) : phénotype normal, a eu une fille albinos (III-2) $\rightarrow$ obligatoirement $[A//a]$.
    \item \textbf{Génération III :} Anne (III-2) albinos $\rightarrow$ $[a//a]$. Jean (III-1) : phénotype normal, issu de $[A//a] \times [A//a]$. \textbf{Avant test biologique}, $P([A//A]) = \frac{1}{3}$, $P([A//a]) = \frac{2}{3}$. \textbf{L'électrophorèse Hb n'informe pas sur l'albinisme.} On conserve la probabilité conditionnelle : Jean est $[A//a]$ avec $p = \frac{2}{3}$.
\end{itemize}

\medskip
\noindent \textbf{B. Drépanocytose (Gène $\beta$-globine \emph{HBB}, Chr 11, autosomique à codominance). Allèles : $S^A$ (Hb A), $S^S$ (Hb S).}
\begin{itemize}
    \item \textbf{Génération I :} Grand-mère maternelle (I-3) drépanocytaire SS $\rightarrow$ $[S^S//S^S]$. Grand-père maternel (I-4) phénotype normal, a eu un enfant SS (III-4) $\rightarrow$ obligatoirement porteur $[S^A//S^S]$ (AS).
    \item \textbf{Génération II :} Mère de Marie (II-3) : $[S^S//S^S]$. Père de Marie (II-4) : $[S^A//S^S]$.
    \item \textbf{Génération III :} Croisement parental : $[S^A//S^S] \times [S^S//S^S] \rightarrow \frac{1}{2} [S^A//S^S]$ (AS) ; $\frac{1}{2} [S^S//S^S]$ (SS).
    \item Paul (III-4) est SS $\rightarrow$ $[S^S//S^S]$. Marie (III-3) phénotype normal. \textbf{Électrophorèse de Marie : AS} $\rightarrow$ Génotype certain $[S^A//S^S]$. L'examen biologique lève l'incertitude du pedigree ($1/2$ a priori $\rightarrow$ certitude 1).
    \item Jean (III-1) : \textbf{Électrophorèse AA} $\rightarrow$ Génotype certain $[S^A//S^A]$.
\end{itemize}

\medskip
\noindent \textbf{C. Groupes sanguins ABO (Chr 9) et Rhésus (Chr 1)}
\begin{itemize}
    \item Jean : Groupe A Rh+. Génotypes ABO possibles : $[I^A I^A]$ ou $[I^A i]$. Rhésus possibles : $[Rh^+//Rh^+]$ ou $[Rh^+//Rh^-]$.
    \item Marie : Groupe B Rh-. Génotype ABO possibles : $[I^B I^B]$ ou $[I^B i]$. Rhésus : $[Rh^-//Rh^-]$ (seul génotype possible pour phénotype Rh-).
\end{itemize}
\textbf{Indépendance :} Les quatre loci (\emph{HBB}, \emph{OCA2}, \emph{ABO}, \emph{RH}) sont sur des chromosomes distincts $\rightarrow$ ségrégation indépendante (loi de Mendel 2 / brassage interchromosomique).

\medskip
\noindent \textbf{Étape 2 : Croisements du couple consultatif (Jean $\times$ Marie) et Calcul des risques}

\medskip
\noindent \textbf{1. Risque Drépanocytose (Gène \emph{HBB})}
\begin{center}
\textbf{Croisement :} Jean $[S^A//S^A]$ $\times$ Marie $[S^A//S^S]$
\end{center}
\begin{center}
\begin{tabular}{|c|c|c|}\hline
\rowcolor{popblueL} \textbf{Gamètes $\downarrow$ / $\rightarrow$} & \textbf{Jean : $S^A$} &  \\ \hline
\textbf{Marie : $S^A$} & $[S^A//S^A]$ (AA) & $\frac{1}{2}$ \\ \hline
\textbf{Marie : $S^S$} & $[S^A//S^S]$ (AS) & $\frac{1}{2}$ \\ \hline
\end{tabular}
\end{center}
$\rightarrow$ \textbf{Risque enfant SS (Drépanocytaire majeur) = 0.} \\
$\rightarrow$ Risque enfant AS (Porteur sain, trait drépanocytaire) = $\frac{1}{2}$ (50\%). \\
\textit{Interprétation : Le père n'ayant pas l'allèle $S^S$, il ne peut pas le transmettre. L'enfant ne peut être au maximum que porteur hétérozygote (AS), cliniquement sain en conditions normales.}

\medskip
\noindent \textbf{2. Risque Albinisme (Gène \emph{OCA2})}
\begin{center}
\textbf{Croisement :} Jean (Probabilité $\frac{2}{3}$ d'être $[A//a]$) $\times$ Marie (Statut inconnu a priori pour albinisme).
\end{center}
\textit{Hypothèse populationnelle (Cameroun) :} Fréquence allélique $q$ de l'allèle $a$ $\approx 1/100$ (prévalence phénotypique $\approx q^2 = 1/10\,000$). Probabilité que Marie soit porteuse $[A//a] \approx 2q = 1/50 = 0,02$. Probabilité qu'elle soit $[A//A] \approx 0,98$.
\begin{itemize}
    \item \textbf{Cas 1 (seul cas à risque) :} Jean $[A//a]$ ($p=2/3$) \textbf{et} Marie $[A//a]$ ($p=1/50$). Croisement $[A//a] \times [A//a] \rightarrow \frac{1}{4} [a//a]$.
    \item \textbf{Autres cas :} Jean $[A//A]$ ($p=1/3$) ou Marie $[A//A]$ ($p=49/50$) $\rightarrow$ Risque 0.
\end{itemize}
\textbf{Risque global Albinisme} $= P(\text{Jean } Aa) \times P(\text{Marie } Aa) \times \frac{1}{4} = \frac{2}{3} \times \frac{1}{50} \times \frac{1}{4} = \frac{2}{600} = \frac{1}{300} \approx \mathbf{0,33\%}$ (\num{0,33}\%).
\textit{Note : Si un test moléculaire (\emph{OCA2}) montre que Jean n'est pas porteur, le risque tombe à 0. En conseil génétique, on annonce le risque résiduel calculé.}

\medskip
\noindent \textbf{3. Risque combiné (Drépanocytose SS ET Albinisme aa)}
Les gènes \emph{HBB} (Chr 11) et \emph{OCA2} (Chr 15) sont sur chromosomes différents $\rightarrow$ \textbf{Indépendance totale (Brassage interchromosomique)}.
\[
P(\text{SS et } aa) = P(\text{SS}) \times P(aa) = 0 \times \frac{1}{300} = \mathbf{0}.
\]
\textbf{Conclusion :} Il est impossible d'avoir un enfant drépanocytaire SS car le père est homozygote normal (AA) pour l'hémoglobine.

\medskip
\noindent \textbf{4. Groupes sanguins ABO et Rhésus de la descendance}
\textbf{ABO :} Jean (A : $I^A I^A$ ou $I^A i$) $\times$ Marie (B : $I^B I^B$ ou $I^B i$).
Quatre combinaisons parentales possibles selon les génotypes réels (inconnus sans typage parental ou étude de la descendance).
\begin{itemize}
    \item Si $I^A I^A \times I^B I^B$ $\rightarrow$ 100\% AB.
    \item Si $I^A I^A \times I^B i$ $\rightarrow$ 50\% A, 50\% AB.
    \item Si $I^A i \times I^B I^B$ $\rightarrow$ 50\% B, 50\% AB.
    \item Si $I^A i \times I^B i$ $\rightarrow$ 25\% A, 25\% B, 25\% AB, 25\% O.
\end{itemize}
\textbf{Phénotypes possibles au total : A, B, AB, O (tous possibles si les deux parents sont hétérozygotes).}

\textbf{Rhésus :} Jean (Rh+ : $Rh^+ Rh^+$ ou $Rh^+ Rh^-$) $\times$ Marie (Rh- : $Rh^- Rh^-$).
\begin{itemize}
    \item Si Jean $Rh^+ Rh^+$ (homozigote) : 100\% enfants Rh+.
    \item Si Jean $Rh^+ Rh^-$ (hétérozygote) : 50\% Rh+, 50\% Rh-.
\end{itemize}
\textbf{Risque hémolytique fœtal (Maladie hémolytique du nouveau-né par anti-D) :} Existe si fœtus Rh+ et mère Rh-. \textbf{Prophylaxie standard :} Injection d'Immunoglobulines Anti-D (RhoGAM) à 28 SA, puis dans les 72h post-partum si nouveau-né Rh+. Recherche d'anticorps irréguliers (RAI) mensuelle pendant la grossesse.

\medskip
\noindent \textbf{Étape 3 : Rédaction de l'avis de conseil génétique (Modèle)}

\medskip
\noindent \textbf{CENTRE HOSPITALIER ESSOS - SERVICE DE GÉNÉTIQUE MÉDICALE} \\
\textbf{COMPTE-RENDU DE CONSEIL GÉNÉTIQUE PRÉCONCEPTIONNEL}

\medskip
\noindent \textbf{Consultants :} M. Jean Ngo Bell (32 ans) \& Mme Marie Ngo Bell (28 ans). \\
\textbf{Date :} 15 Octobre 2025. \\
\textbf{Motif :} Projet parental, antécédents familiaux d'albinisme (côté paternel) et drépanocytose (côté maternel).

\medskip
\noindent \textbf{1. ANAMNÈSE ET ARBRE GÉNÉALOGIQUE} \\
Pedigree établi sur 4 générations (voir figure). Deux pathologies autosomiques distinctes identifiées :
\begin{itemize}
    \item \textbf{Albinisme oculocutané type 2 (OCA2) :} Transmission autosomique récessive. Cas index : Grand-père paternel (I-1) et Tante Anne (III-2). Père de Jean (II-1) et Mère de Jean (II-2) sont porteurs sains obligatoires ($Aa$). Jean (III-1) phénotype normal : probabilité $2/3$ d'être porteur.
    \item \textbf{Drépanocytose (Hb SS) :} Transmission autosomique récessive (codominance allélique). Cas index : Grand-mère maternelle (I-3) SS, Oncle Paul (III-4) SS. Mère de Marie (II-3) SS. Père de Marie (II-4) porteur obligé (AS). Marie (III-3) : Électrophorèse \textbf{AS confirmée} (porteuse saine). Jean (III-1) : Électrophorèse \textbf{AA confirmée} (non porteur).
\end{itemize}

\medskip
\noindent \textbf{2. ANALYSE GÉNÉTIQUE ET RISQUES CALCULÉS} \\
\textbf{Génotypes confirmés du couple :}
\begin{itemize}
    \item Jean : \textbf{Hb : $S^A S^A$ (AA)} ; \textbf{Albinisme : $A-$ (Porteur probant $p=2/3$)} ; \textbf{ABO : A} ; \textbf{Rh : +}.
    \item Marie : \textbf{Hb : $S^A S^S$ (AS)} ; \textbf{Albinisme : $A-$ (Risque porteur population $\approx 2\%$)} ; \textbf{ABO : B} ; \textbf{Rh : -}.
\end{itemize}

\textbf{Risques pour la descendance :}
\begin{center}
\begin{tabularx}{\linewidth}{|l|X|c|}\hline
\rowcolor{popblueL} \textbf{Pathologie / Caractère} & \textbf{Croisement effectif / Situation} & \textbf{Risque enfant atteint} \\ \hline
\textbf{Drépanocytose (SS)} & $S^A S^A \times S^A S^S$ & \textbf{0 \%} (Impossible) \\ \hline
\textbf{Trait drépanocytaire (AS)} & $S^A S^A \times S^A S^S$ & \textbf{50 \%} (Enfant porteur sain) \\ \hline
\textbf{Albinisme (aa)} & $Aa_{(2/3)} \times Aa_{(\approx 1/50)} \times 1/4$ & \textbf{$\approx 0,33 \%$ (1/300)} \\ \hline
\textbf{Double atteinte (SS + aa)} & Indépendance chromosomique (Chr 11 \& 15) & \textbf{0 \%} \\ \hline
\textbf{Incompatibilité Rhésus} & Mère Rh- / Père Rh+ (hétéro/homozygote) & \textbf{Risque réel si fœtus Rh+} (50 à 100\%) \\ \hline
\end{tabularx}
\end{center}

\medskip
\noindent \textbf{3. RECOMMANDATIONS ET PRÉVENTION}
\begin{enumerate}
    \item \textbf{Drépanocytose :} \textbf{Risque nul d'enfant malade (SS).} Rassurer le couple sur ce point majeur. L'enfant sera soit AA (50\%), soit AS (50\%). Expliquer que l'enfant AS est cliniquement sain (sauf conditions extrêmes : haute altitude $>2500$m, déshydratation sévère, anesthésie générale hypoxique). \textbf{Dépistage néonatal systématique} à la naissance (test de Guthrie / électrophorèse Hb) pour confirmer le statut et permettre une prise en charge précoce si AS (éducation thérapeutique) ou SS (impossible ici mais protocole standard).
    \item \textbf{Albinisme :} Risque faible mais non nul ($\approx 1/300$). Proposer un \textbf{test moléculaire (séquençage \emph{OCA2})} à Jean pour lever l'incertitude (porteur ou non). Si Jean non porteur $\rightarrow$ Risque 0 définitif. Si Jean porteur, proposer le test à Marie. En cas de grossesse et si les deux sont porteurs, \textbf{diagnostic prénatal (DPN)} possible sur ADN fœtal (liquide amniotique 16-18 SA ou trophoblaste 11-13 SA).
    \item \textbf{Allo-immunisation Rhésus :} Marie est Rh-. Risque d'hémolyse fœtale si fœtus Rh+. \textbf{Prophylaxie obligatoire :} Injection d'Immunoglobulines Anti-D (RhoGAM, 300 µg) à 28 SA, puis dans les 72h post-partum si nouveau-né Rh+. Recherche d'anticorps irréguliers (RAI) mensuelle pendant la grossesse. Surveillance échographique (vélocimétrie artère cérébrale moyenne) en cas d'allo-immunisation.
    \item \textbf{Suivi prénatal global :} Consultation prénatale précoce (avant 12 SA). Supplémentation en acide folique (0,4 mg/j) dès le désir de grossesse (prévention anomalies du tube neural). Échographies morphologiques aux 3 trimestres. Bilan infectieux standard (VIH, Syphilis, Hépatites B/C, Toxoplasmose, Rubéole).
    \item \textbf{Aspects psychosociaux et éthiques :} Aborder la stigmatisation potentielle de l'albinisme au Cameroun (protection solaire stricte indice 50+, vêtements couvrants, suivi dermatologique annuel pour dépistage cancers cutanés, suivi ophtalmologique pour photophobie/nystagmus/acuité visuelle). Insister sur la normalité de vie d'un enfant AS (porteur sain). Respect strict du secret médical et de l'autonomie décisionnelle du couple (droit de savoir / droit de ne pas savoir). Orientation vers les associations de patients (ex: ASMOCAM pour l'albinisme, associations de lutte contre la drépanocytose).
\end{enumerate}

\medskip
\noindent \textbf{Signature du Généticien Clinicien} \\
\textit{Dr. \_\_\_\_\_\_\_\_\_\_\_\_\_\_\_\_\_\_\_\_\_\_\_\_\_\_\_\_\_\_\_\_\_\_\_\_\_\_\_\_\_\_\_\_\_\_\_\_}
\end{exoresolu}

\courssub{Bilan de l'activité}
\begin{aretenir}[Compétences intégrées et points clés validés]
Cette activité de synthèse a permis de mobiliser l'ensemble des savoir-faire du Module I (Le Monde Vivant) dans une situation professionnelle authentique de conseil génétique :
\begin{enumerate}
    \item \textbf{Modélisation génétique rigoureuse :} Construction d'un pedigree multi-caractères sur 4 générations. Distinction claire entre déduction probabiliste par le pedigree (risque conditionnel $2/3$ pour l'albinisme de Jean) et certitude diagnostique apportée par le phénotype biochimique (électrophorèse Hb confirmant Jean AA et Marie AS).
    \item \textbf{Application du brassage interchromosomique :} Utilisation concrète de l'indépendance des gènes (\emph{HBB} Chr 11, \emph{OCA2} Chr 15, \emph{ABO} Chr 9, \emph{RH} Chr 1) pour calculer les risques combinés (produit des probabilités marginales). Vérification que le risque double (SS + aa) est nul car l'un des risques marginaux est nul.
    \item \textbf{Calcul de risque génétique complet :} Intégration de la fréquence allélique populationnelle (albinisme au Cameroun, $q \approx 1/100$) pour le conjoint non informé par le pedigree (Marie). Distinction fondamentale entre risque nul (certitude génétique : père AA pour HbS $\rightarrow$ pas d'allèle $S^S$ transmissible) et risque résiduel non nul (albinisme).
    \item \textbf{Communication médicale et prévention :} Traduction des probabilités ($1/2, 1/300, 0$) en langage accessible pour le couple. Distinction cruciale entre "porteur sain" (AS, vie normale) et "malade" (SS). Prise en charge proactive de l'incompatibilité Rhésus (prévention primaire par Anti-D, surveillance RAI).
    \item \textbf{Éthique et santé publique camerounaise :} Place du dépistage néonatal systématique (politique MINSANTE), indication du diagnostic prénatal (DPN) encadrée par la loi, lutte contre la stigmatisation des albinos (croyances, protection solaire), éducation thérapeutique pour la drépanocytose.
\end{enumerate}
\textbf{Erreur critique évitée :} Ne pas confondre le risque \emph{a priori} de Jean pour l'albinisme ($2/3$) avec son statut pour la drépanocytose (levé par l'électrophorèse : certitude d'être $S^A S^A$, risque 0 d'être porteur $S^S$). L'examen biologique direct prime sur la déduction indirecte par le pedigree.
\end{aretenir}

\begin{savaistu}[Contexte camerounais : Drépanocytose et Albinisme - Données épidémiologiques et prise en charge]
\begin{itemize}
    \item \textbf{Drépanocytose (Maladie génétique la plus fréquente au Cameroun) :} Fréquence allélique $S^S$ (HbS) : 10 à 20\% selon les régions (plus élevée au Sud/Forêt, plus basse au Nord/Sahel). Prévalence SS à la naissance $\approx 1,5\%$ (environ 6\,000 naissances SS/an). \textbf{Prise en charge standardisée (PNLCD - Programme National de Lutte Contre la Drépanocytose) :} Dépistage néonatal systématique (points de Guthrie sur papier filtre), prise en charge précoce avant 3 mois : Pénicilline V per os (prophylaxie anti-pneumococcique) jusqu'à 5 ans, Vaccination anti-pneumocoque (Prevenar 13) et anti-Haemophilus, Acide folique 5 mg/j à vie, Hydroxyurée (seule thérapie modificatrice de la maladie, accessible via pharmacies hospitalières), Éducation thérapeutique (hydratation 2-3L/j, évitement froid/excès d'effort, consultation fièvre $>38,5^\circ$C urgence). Suivi transcranien (Doppler) annuel dès 2 ans pour prévention AVC.
    \item \textbf{Albinisme (OCA2 - Albinisme oculocutané type 2) :} Prévalence $\approx 1/5\,000$ à $1/10\,000$ (plus fréquent qu'en Europe $1/20\,000$). Gène \emph{OCA2} (Chr 15q11.2-q13). Mutations fondatrices spécifiques en Afrique subsaharienne (ex: délétion 2,7 kb). \textbf{Prise en charge multidisciplinaire :} Protection solaire stricte (crème indice 50+ renouvelée 2h, vêtements couvrants, chapeau large bord, lunettes solaires UV400), Suivi dermatologique semestriel (dépistage carcinomes épidermoïdes précoces sur zones photo-exposées), Suivi ophtalmologique annuel (correction optique, filtres teintés, chirurgie nystagmus/strabisme si indiqué). \textbf{Enjeu sociétal :} Lutte contre les discriminations et croyances néfastes (sorcellerie, sacrifices rituels, exclusion scolaire) via les associations (ASMOCAM, APAC, Albinos sans Frontières) et l'éducation inclusive.
    \item \textbf{Conseil génétique au Cameroun :} En plein développement (Centres : Hôpital Central Yaoundé, Hôpital Général Douala, CHU Essos, Facultés de Médecine Yaoundé I/II, Buea). Nécessité de former des conseillers génétiques (Master spécialisé), de rendre les tests moléculaires (PCR, séquençage NGS) accessibles et remboursables, et d'intégrer la génétique clinique dans le cursus médical et paramédical.
\end{itemize}
\end{savaistu}

\newpage

\coursec{Méthodes et exercices résolus}

\courssub{Méthodes et exercices résolus}

\begin{methode}[Analyser un croisement monohybride (dominance ou codominance)]
    \textbf{Objectif :} Prédire les phénotypes et génotypes de la descendance (génération F1, F2 ou test-cross) à partir des allèles parentaux.
    \begin{enumerate}
        \item \textbf{Identifier le gène et ses allèles :} Noter les allèles (ex. $A$ dominant, $a$ récessif ; ou $I^A, I^B, i$ pour ABO). Préciser le type de relation allélique (dominance complète, codominance, allèles multiples).
        \item \textbf{Déterminer les gamètes parentaux :} Un individu diploïde produit des gamètes haploïdes par méiose. Un homozygote ne produit qu'un type de gamète ; un hétérozygote en produit deux (séparation des allèles).
        \item \textbf{Construire le carré de Punnett :} Croiser les gamètes mâles (colonnes) et femelles (lignes). Chaque case = un génotype possible du zygote.
        \item \textbf{Compter les génotypes et phénotypes :} Regrouper les cases identiques. Exprimer en pourcentages ou rapports (ex. 3/4, 1/4).
        \item \textbf{Vérifier la cohérence :} Somme des fréquences = 1 (ou 100\%). Respecter la notation demandée (fractions, pourcentages, rapports).
    \end{enumerate}
    \textbf{Formules clés :} Fréquence d'un génotype = (nb cases correspondantes) / (nb total cases). Test-cross : phénotype récessif $\times$ phénotype dominant $\to$ si 1/2 récessifs dans descendance $\Rightarrow$ parent dominant hétérozygote.
\end{methode}

\begin{exoresolu}[Monohybridisme et groupes sanguins ABO (Codominance / Allèles multiples)]
    \textbf{Énoncé :} Au Cameroun, les groupes sanguins ABO suivent la codominance des allèles $I^A$ et $I^B$ (tous deux dominants sur $i$). Un homme de groupe A (génotype $I^A i$) épouse une femme de groupe B (génotype $I^B i$).
    \begin{enumerate}
        \item Déterminer les phénotypes et leurs fréquences attendus parmi leurs enfants.
        \item L'un de leurs fils, de groupe AB, épouse une femme de groupe O. Quel est le risque d'avoir un enfant de groupe O ?
    \end{enumerate}
    \tcblower
    \textbf{Solution détaillée :}
    \begin{enumerate}
        \item \textbf{Parents :} Père $I^A i$ (Groupe A) $\times$ Mère $I^B i$ (Groupe B).
        \textbf{Gamètes :} Père $\to$ $I^A$ ou $i$ (fréquence 1/2 chacun). Mère $\to$ $I^B$ ou $i$ (fréquence 1/2 chacun).
        \textbf{Carré de Punnett :}
        \begin{center}
        \begin{tabular}{|c|c|c|}\hline
        & $I^B$ (1/2) & $i$ (1/2) \\ \hline
        $I^A$ (1/2) & $I^A I^B$ (Groupe AB) & $I^A i$ (Groupe A) \\ \hline
        $i$ (1/2) & $I^B i$ (Groupe B) & $ii$ (Groupe O) \\ \hline
        \end{tabular}
        \end{center}
        \textbf{Fréquences phénotypiques (F1) :}
        \begin{itemize}
            \item Groupe AB ($I^A I^B$) : 1/4 = 25\%
            \item Groupe A ($I^A i$) : 1/4 = 25\%
            \item Groupe B ($I^B i$) : 1/4 = 25\%
            \item Groupe O ($ii$) : 1/4 = 25\%
        \end{itemize}
        \item \textbf{Croisement :} Fils Groupe AB ($I^A I^B$) $\times$ Femme Groupe O ($ii$).
        \textbf{Gamètes :} Fils $\to$ $I^A$ ou $I^B$ (pas de $i$). Femme $\to$ $i$ uniquement.
        \textbf{Descendance possible :} $I^A i$ (Groupe A) à 50\% et $I^B i$ (Groupe B) à 50\%.
        \textbf{Réponse :} Le risque d'avoir un enfant de groupe O ($ii$) est \textbf{0\% (impossible)} car le père ne possède pas l'allèle $i$ à transmettre.
    \end{enumerate}
    \textbf{Conclusion :} La codominance $I^A/I^B$ et la récessivité de $i$ permettent 4 phénotypes en F1. Un individu AB ne peut transmettre que $I^A$ ou $I^B$, excluant le groupe O chez sa descendance avec un partenaire O.
\end{exoresolu}

\begin{methode}[Analyser un dihybridisme (brassage interchromosomique) et réaliser un test-cross]
    \textbf{Objectif :} Étudier la transmission simultanée de deux gènes non liés (sur chromosomes différents ou lointains sur le même chromosome) et tester le génotype d'un individu à phénotype dominant.
    \begin{enumerate}
        \item \textbf{Vérifier l'indépendance :} Les deux gènes sont sur des paires de chromosomes différentes (loi de Mendel 2) $\Rightarrow$ brassage interchromosomique aléatoire.
        \item \textbf{Écrire les génotypes parentaux (P) :} Noter les allèles pour chaque gène (ex. $A/a$ et $B/b$). Préciser la phase (cis/trans) si hétérozygote double.
        \item \textbf{Déterminer les gamètes (Méiose) :} 4 types de gamètes en fréquences égales (1/4 chacun) pour un double hétérozygote : $AB, Ab, aB, ab$. Pour un homozygote : 1 type. Pour un simple hétérozygote : 2 types.
        \item \textbf{Croiser les gamètes (Carré 4$\times$4 = 16 cases) :} Obtenir les 16 génotypes de F2. Ratio phénotypique classique (dominance complète) : 9/16 double dominant : 3/16 dom A rec B : 3/16 rec A dom B : 1/16 double récessif.
        \item \textbf{Test-cross (croisement test) :} Croiser l'individu à phénotype dominant (génotype inconnu : $A\_ B\_$) avec un double récessif ($aa bb$). Le double récessif ne donne que des gamètes $ab$. Le phénotype de la descendance reflète \emph{directement} le génotype des gamètes du parent testé.
        \item \textbf{Interpréter :} 4 phénotypes en proportions égales (1/4 chacun) $\Rightarrow$ parent testé double hétérozygote ($Aa Bb$). 2 phénotypes (1/2 chacun) $\Rightarrow$ parent hétérozygote pour un seul gène. 1 seul phénotype $\Rightarrow$ parent double homozygote dominant.
    \end{enumerate}
\end{methode}

\begin{exoresolu}[Dihybridisme et test-cross : Couleur et forme de graines (Pois)]
    \textbf{Énoncé :} Chez le pois, la couleur des graines (Jaune $J$ dominant sur Vert $j$) et la forme (Lisse $L$ dominant sur Ridée $l$) sont déterminées par deux gènes non liés. On croise une lignée pure à graines jaunes lisses avec une lignée pure à graines vertes ridées. On obtient la F1. On pratique ensuite un test-cross sur un individu F1.
    \begin{enumerate}
        \item Préciser les génotypes des parents (P) et de la F1.
        \item Déterminer les phénotypes et leurs fréquences dans la F2 (autofécondation F1).
        \item Déterminer les phénotypes et leurs fréquences attendus dans la descendance du test-cross (F1 $\times$ $jj ll$).
    \end{enumerate}
    \tcblower
    \textbf{Solution détaillée :}
    \begin{enumerate}
        \item \textbf{P :} Lignée pure Jaune Lisse = $JJ LL$. Lignée pure Verte Ridée = $jj ll$.
        \textbf{Gamètes P :} $JL$ et $jl$ uniquement.
        \textbf{F1 :} Tous $Jj Ll$ (Double hétérozygotes). Phénotype : Jaunes Lisses.
        \item \textbf{F1 $\times$ F1 (Autofécondation) :}
        \textbf{Gamètes F1 (4 types, 1/4 chacun) :} $JL, Jl, jL, jl$.
        \textbf{Carré de Punnett (16 cases) :} Ratio phénotypique classique 9:3:3:1.
        \begin{itemize}
            \item Jaunes Lisses ($J\_ L\_$) : 9/16 = 56,25\%
            \item Jaunes Ridée ($J\_ ll$) : 3/16 = 18,75\%
            \item Vertes Lisses ($jj L\_$) : 3/16 = 18,75\%
            \item Vertes Ridée ($jj ll$) : 1/16 = 6,25\%
        \end{itemize}
        \item \textbf{Test-cross :} F1 ($Jj Ll$) $\times$ Témoin récessif ($jj ll$).
        \textbf{Gamètes F1 :} $JL, Jl, jL, jl$ (1/4 chacun).
        \textbf{Gamètes Témoin :} $jl$ (100\%).
        \textbf{Fécondation :} Chaque gamète F1 fusionne avec $jl$.
        \textbf{Descendance (phénotypes = gamètes F1) :}
        \begin{itemize}
            \item $JL$ + $jl$ $\to$ $Jj Ll$ : Jaune Lisse (1/4)
            \item $Jl$ + $jl$ $\to$ $Jj ll$ : Jaune Ridée (1/4)
            \item $jL$ + $jl$ $\to$ $jj Ll$ : Verte Lisse (1/4)
            \item $jl$ + $jl$ $\to$ $jj ll$ : Verte Ridée (1/4)
        \end{itemize}
        \textbf{Résultat :} 4 phénotypes en proportions égales (25\% chacun). Cela confirme que l'individu F1 était bien double hétérozygote et que les gènes sont en brassage interchromosomique (indépendants).
    \end{enumerate}
    \textbf{Conclusion :} Le test-cross est l'outil décisif pour révéler la constitution gamétique d'un individu à phénotype dominant.
\end{exoresolu}

\begin{methode}[Établir une carte factorielle (Linkage partiel / Crossing-over)]
    \textbf{Objectif :} Calculer la distance génétique entre deux gènes liés sur le même chromosome à partir des fréquences de recombinants (crossing-over).
    \begin{enumerate}
        \item \textbf{Identifier le linkage :} Les fréquences phénotypiques de la descendance (souvent issue d'un test-cross) s'écartent du 1:1:1:1. Deux classes majoritaires (parentaux / non-recombinants) et deux classes minoritaires (recombinants).
        \item \textbf{Déterminer la phase allélique (Cis ou Trans) du parent hétérozygote :}
            \begin{itemize}
                \item \textbf{Cis (couplage) :} Allèles dominants sur un chromosome, récessifs sur l'autre : $AB // ab$. Parentaux = $AB, ab$. Recombinants = $Ab, aB$.
                \item \textbf{Trans (répulsion) :} Allèles dominants/récessifs mélangés : $Ab // aB$. Parentaux = $Ab, aB$. Recombinants = $AB, ab$.
            \end{itemize}
        \item \textbf{Calculer la fréquence de crossing-over ($F_{CO}$) :}
        \[ F_{CO} = \frac{\text{Nombre d'individus recombinants}}{\text{Total des individus}} \times 100 \]
        (Exprimée en \%, équivalente à la distance en \textbf{centiMorgans (cM)} ou \textbf{unités de carte (UC)}).
        \item \textbf{Construire la carte :} Placer les gènes sur une ligne. Distance = $F_{CO}$ cM. Si 3 gènes : utiliser les doubles recombinants (rares) pour ordonner le gène du milieu. Distance $A-C$ = Distance $A-B$ + Distance $B-C$ (si pas d'interférence forte).
        \item \textbf{Vérifier :} $F_{CO} \le 50\%$. Si $F_{CO} \approx 50\%$ $\Rightarrow$ gènes indépendants (brassage interchromosomique ou très éloignés).
    \end{enumerate}
    \textbf{Formule (3 gènes) :} Distance (cM) = $\%$ Recombinants = $\frac{2 \times \text{Nb doubles CO} + \text{Nb simples CO}}{\text{Total}} \times 100$.
\end{methode}

\begin{exoresolu}[Carte factorielle à 3 gènes : Drosophila (Yeux, Ailes, Corps)]
    \textbf{Énoncé :} Chez la drosophile, trois gènes liés : $se$ (yeux sépia, récessif), $e$ (corps ébène, récessif), $cv$ (ailes courtes, récessif). Une femelle triple hétérozygote (phénotype sauvage) issue d'un croisement entre une mouche sauvage pure et une mouche triple mutante est test-croisée avec un mâle triple mutant. Résultats sur 1000 descendants :
    \begin{center}
    \begin{tabular}{|l|c|}\hline Phénotype & Effectif \\ \hline
    Sauvage ($se^+ e^+ cv^+$) & 402 \\ \hline
    Sepia, Ébène, Courte ($se\ e\ cv$) & 398 \\ \hline
    Sepia, Courte ($se\ cv$) & 88 \\ \hline
    Ébène ($e$) & 92 \\ \hline
    Sepia, Ébène ($se\ e$) & 8 \\ \hline
    Courte ($cv$) & 7 \\ \hline
    Sauvage, Courte ($se^+ e^+ cv$) & 3 \\ \hline
    Sepia ($se$) & 2 \\ \hline
    \end{tabular}
    \end{center}
    \textbf{Établir la carte factorielle (ordre des gènes et distances).}
    \tcblower
    \textbf{Solution détaillée :}
    \begin{enumerate}
        \item \textbf{Identifier les classes parentales (non-recombinants) :} Les plus fréquentes.
        $se^+ e^+ cv^+$ (402) et $se\ e\ cv$ (398). Total parentaux = 800.
        \textbf{Phase de la mère (Cis) :} $se^+ e^+ cv^+ // se\ e\ cv$.
        \item \textbf{Identifier les doubles recombinants (rares) :} Les deux classes les plus rares : $se^+ e^+ cv$ (3) et $se\ e\ cv^+$ (phénotype "Sepia" seul, 2).
        Dans les doubles CO, l'allèle du gène du milieu est inversé par rapport aux parentaux.
        Parentaux : $+ + +$ et $se\ e\ cv$.
        Doubles CO observés : $+ + cv$ (3) et $se\ e\ +$ (2).
        Comparaison : $+ + +$ $\to$ $+ + cv$ : $cv$ a changé. $se\ e\ cv$ $\to$ $se\ e\ +$ : $cv$ a changé.
        \textbf{Le gène du milieu est $cv$ (ailes courtes).}
        \textbf{Ordre :} $se$ -- $cv$ -- $e$ (ou $e$ -- $cv$ -- $se$).
        \item \textbf{Calculer les distances :}
        \textbf{Intervalle $se - cv$ :} Recombinants = Simples CO ($se-cv$) + Doubles CO.
        Phénotypes recombinants $se-cv$ : $se\ cv$ (88) et $e$ (92) [car $e$ = $se^+ e\ cv^+$]. + Doubles CO (5).
        Total $RF_{se-cv} = \frac{88 + 92 + 5}{1000} = 18,5\% = 18,5 \text{ cM}$.
        \textbf{Intervalle $cv - e$ :} Recombinants = Simples CO ($cv-e$) + Doubles CO.
        Phénotypes recombinants $cv-e$ : $se\ e$ (8) et $cv$ (7) [car $cv$ = $se^+ e^+ cv$]. + Doubles CO (5).
        Total $RF_{cv-e} = \frac{8 + 7 + 5}{1000} = 2,0\% = 2,0 \text{ cM}$.
        \textbf{Distance $se - e$ :} $18,5 + 2,0 = 20,5 \text{ cM}$ (vérification : recombinants directs $se-e$ = 189 $\to$ 18,9\% $\approx$ 20,5\% moins doubles CO comptés 2 fois).
    \end{enumerate}
    \textbf{Carte factorielle :}
    \begin{popfigure}
    \begin{tikzpicture}[scale=0.8]
        \draw[thick, popdark] (0,0) -- (20.5,0);
        \foreach \x/\label in {0/se, 18.5/cv, 20.5/e} {
            \draw (\x,0.2) -- (\x,-0.2) node[below] {\label\ (\x\ cM)};
        }
    \end{tikzpicture}
    \legende{Carte factorielle des gènes $se$, $cv$, $e$ chez la drosophile. Distances en centiMorgans (cM).}
    \end{popfigure}
    \textbf{Conclusion :} L'ordre est $se$ -- $cv$ -- $e$. Distances : $se-cv = 18,5$ cM, $cv-e = 2,0$ cM.
\end{exoresolu}

\begin{methode}[Construire et exploiter un pedigree (Hérédité autosomique)]
    \textbf{Objectif :} Représenter l'hérédité d'un caractère sur plusieurs générations et déduire le mode de transmission (dominant/récessif, autosomique) et les génotypes probables.
    \begin{enumerate}
        \item \textbf{Symboles standards :} $\square$ = Homme, $\bigcirc$ = Femme. Symboles pleins (noirs) = atteints. Symboles vides = sains. Trait horizontal = union. Trait vertical = descendance. Générations numérotées I, II, III... Individus numérotés 1, 2, 3... dans chaque génération.
        \item \textbf{Analyser le mode de transmission :}
            \begin{itemize}
                \item \textbf{Récessif autosomique :} L'affection saute des générations. Parents sains peuvent avoir des enfants atteints (hétérozygotes porteurs). Atteints $\times$ Atteints $\to$ 100\% atteints. Consanguinité fréquente.
                \item \textbf{Dominant autosomique :} L'affection ne saute pas de génération (sauf non-pénétrance). Un parent atteint $\times$ sain $\to$ 50\% atteints. Hommes et femmes touchés équitablement.
            \end{itemize}
        \item \textbf{Déduire les génotypes :} Partir des individus atteints (génotype connu si récessif : $aa$ ; si dominant : $Aa$ ou $AA$). Remonter vers les parents et descendre vers les enfants. Noter les allèles sur le pedigree.
        \item \textbf{Calculer le risque génétique (Prévision) :} Pour un couple consultant (ex. II-3 et II-4), déterminer la probabilité qu'ils aient un enfant atteint.
            \begin{itemize}
                \item Probabilité que le parent soit porteur (si phénotype sain mais parent d'atteint ou frère/soeur d'atteint).
                \item Probabilité de transmission de l'allèle muté (1/2 si hétérozygote).
                \item Produit des probabilités (indépendance des événements méiotiques).
            \end{itemize}
    \end{enumerate}
\end{methode}

\begin{exoresolu}[Pedigree et risque génétique : Drépanocytose (Cameroun)]
    \textbf{Énoncé :} La drépanocytose (anémie falciforme) est une maladie autosomique récessive (allèle $S$ muté, allèle $A$ normal). Dans une famille de Yaoundé :
    \begin{itemize}
        \item Génération I : Parents sains (I-1, I-2).
        \item Génération II : 3 enfants. II-1 (homme) sain. II-2 (femme) atteinte (drépanocytaire majeure $SS$). II-3 (homme) sain, épouse une femme saine (II-4) sans antécédent familial connu.
        \item Génération III : II-1 a 2 enfants sains. II-3 et II-4 attendent un enfant.
    \end{itemize}
    \begin{enumerate}
        \item Représenter le pedigree (schéma TikZ).
        \item Déterminer les génotypes de I-1, I-2, II-1, II-3.
        \item Calculer le risque pour le couple II-3 $\times$ II-4 d'avoir un enfant drépanocytaire ($SS$).
    \end{enumerate}
    \tcblower
    \textbf{Solution détaillée :}
    \begin{enumerate}
        \item \textbf{Pedigree :}
        \begin{popfigure}
        \begin{tikzpicture}[node distance=1.2cm and 0.8cm, every node/.style={draw, thick, minimum size=0.7cm}]
            % Gen I
            \node[rectangle] (I1) {I-1};
            \node[circle, right=of I1] (I2) {I-2};
            \draw (I1) -- (I2);
            % Gen II
            \node[rectangle, below left=of I1] (II1) {II-1};
            \node[circle, fill=popdark, below=of I1] (II2) {II-2};
            \node[rectangle, below right=of I1] (II3) {II-3};
            \node[circle, right=of II3] (II4) {II-4};
            \draw (I1) |- (II1); \draw (I1) |- (II2); \draw (I1) |- (II3);
            \draw (II3) -- (II4);
            % Gen III
            \node[rectangle, below left=of II1] (III1) {III-1};
            \node[rectangle, below right=of II1] (III2) {III-2};
            \draw (II1) |- (III1); \draw (II1) |- (III2);
            \node[rectangle, dashed, below=of II3] (III3) {III-3 ?};
            \draw (II3) |- (III3);
        \end{tikzpicture}
        \legende{Pédigree drépanocytose. Symboles pleins (noirs) = atteints ($SS$). Cercles = femmes, Carrés = hommes. Trait pointillé = grossesse en cours.}
        \end{popfigure}
        \item \textbf{Génotypes :}
        \begin{itemize}
            \item II-2 atteinte $\Rightarrow$ $SS$.
            \item I-1 et I-2 parents sains d'un enfant $SS$ $\Rightarrow$ obligatoirement porteurs sains (hétérozygotes) $\Rightarrow$ $AS$.
            \item II-1 et II-3 : Phénotype sain, parents $AS \times AS$. Probabilités conditionnelles : $P(AA) = 1/3$, $P(AS) = 2/3$. (On exclut $SS$ car sains).
        \end{itemize}
        \item \textbf{Risque pour II-3 $\times$ II-4 :}
        \begin{itemize}
            \item II-3 : $P(AS) = 2/3$, $P(AA) = 1/3$.
            \item II-4 : Femme saine, sans antécédent. Fréquence allélique $S$ au Cameroun $\approx 0,10$ à $0,20$ selon régions. Prendre $q = 0,15$. $P(AS) \approx 2pq = 2 \times 0,85 \times 0,15 = 0,255$. \textbf{Donnée Bac souvent simplifiée :} $P(\text{II-4 porteuse}) = 1/10$.
            \item Enfant $SS$ si : II-3 transmet $S$ ET II-4 transmet $S$.
            \item $P(\text{II-3 transmet } S) = P(AS) \times 1/2 = (2/3) \times (1/2) = 1/3$.
            \item $P(\text{II-4 transmet } S) = P(AS) \times 1/2 = (1/10) \times (1/2) = 1/20$.
            \item \textbf{Risque total} = $(1/3) \times (1/20) = 1/60 \approx 1,67\%$.
        \end{itemize}
    \end{enumerate}
    \textbf{Conclusion :} Le risque est faible (1,67\%) mais non nul. Le conseil génétique recommande le diagnostic prénatal (DPN) ou le dépistage néonatal systématique (réalisé au Cameroun via le Programme National de Lutte contre la Drépanocytose).
\end{exoresolu}

\begin{methode}[Analyser un pedigree d'hérédité gonosomique (liée au sexe)]
    \textbf{Objectif :} Différencier hérédité récessive liée à l'X (hémophilie, daltonisme) et dominante liée à l'X (rare), et calculer les risques pour les descendants.
    \begin{enumerate}
        \item \textbf{Repérer les indices clés :}
            \begin{itemize}
                \item \textbf{Récessive X :} Prédominance masculine (hommes $X^h Y$ atteints, femmes $X^h X^h$ rares). Transmission "en saut de génération" par les femmes porteuses ($X^H X^h$). \textbf{Pas de transmission père $\to$ fils} (père donne Y au fils). Père atteint $\to$ toutes les filles porteuses (reçoivent son $X^h$).
                \item \textbf{Dominante X :} Femmes plus touchées (2X). Père atteint $\to$ toutes les filles atteintes. Mère atteinte $\to$ 50\% enfants atteints (sexe indifférent).
            \end{itemize}
        \item \textbf{Notation :} Allèles sur le chromosome X : $X^A$ (normal), $X^a$ (muté). Chromosome Y : pas d'allèle pour ce gène. Génotypes : Homme $X^A Y$ (sain), $X^a Y$ (atteint). Femme $X^A X^A$ (saine), $X^A X^a$ (porteuse saine si récessif), $X^a X^a$ (atteinte).
        \item \textbf{Calculer les risques :} Croisement classique. Ex : Mère porteuse ($X^A X^a$) $\times$ Père sain ($X^A Y$).
            \begin{itemize}
                \item Filles : 1/2 $X^A X^A$ (saines), 1/2 $X^A X^a$ (porteuses).
                \item Garçons : 1/2 $X^A Y$ (sains), 1/2 $X^a Y$ (atteints).
            \end{itemize}
        \item \textbf{Attention :} Ne pas confondre "lié au sexe" (gène sur X ou Y) et "limité au sexe" (gène autosomique mais exprimé dans un seul sexe, ex. barbe). Au programme : lié au sexe (X).
    \end{enumerate}
\end{methode}

\begin{exoresolu}[Hérédité gonosomique récessive : Hémophilie A]
    \textbf{Énoncé :} L'hémophilie A (déficit en facteur VIII) est récessive liée à l'X. Un homme hémophile (I-1) épouse une femme saine non porteuse (I-2). Ils ont une fille (II-1) et un fils (II-2). La fille II-1 épouse un homme sain (II-3). Ils ont un fils (III-1) atteint et une fille (III-2) saine.
    \begin{enumerate}
        \item Représenter le pedigree avec les génotypes.
        \item Calculer la probabilité que la fille III-2 soit porteuse.
        \item Si III-2 épouse un homme sain, quel est le risque d'avoir un fils hémophile ?
    \end{enumerate}
    \tcblower
    \textbf{Solution détaillée :}
    \begin{enumerate}
        \item \textbf{Génotypes :}
        \begin{itemize}
            \item I-1 (Homme atteint) : $X^h Y$.
            \item I-2 (Femme saine non porteuse) : $X^H X^H$.
            \item II-1 (Fille) : Reçoit $X^h$ du père, $X^H$ de la mère $\Rightarrow$ \textbf{$X^H X^h$ (Porteuse saine)}.
            \item II-2 (Fils) : Reçoit $Y$ du père, $X^H$ de la mère $\Rightarrow$ \textbf{$X^H Y$ (Sain)}.
            \item II-3 (Mari de II-1, sain) : \textbf{$X^H Y$}.
            \item III-1 (Fils atteint) : $X^h Y$. Il a reçu $Y$ du père (II-3) et $X^h$ de la mère (II-1). Confirme que II-1 est porteuse.
            \item III-2 (Fille saine) : Reçoit $X^H$ du père (II-3). Reçoit $X^H$ ou $X^h$ de la mère (II-1).
        \end{itemize}
        \item \textbf{Probabilité III-2 porteuse :} Elle a reçu le $X$ maternel. Mère II-1 = $X^H X^h$. Probabilité de transmettre $X^h$ = 1/2.
        \textbf{Réponse :} $P(\text{III-2 porteuse}) = \mathbf{1/2}$ (Génotype $X^H X^h$).
        \item \textbf{Risque pour descendance III-2 $\times$ Homme sain ($X^H Y$) :}
        Si III-2 est porteuse (prob 1/2) : Risque fils atteint = 1/2.
        Si III-2 non porteuse (prob 1/2) : Risque = 0.
        \textbf{Risque global d'avoir un fils hémophile} = $P(\text{porteuse}) \times P(\text{transmet } X^h) \times P(\text{fils})$.
        $= (1/2) \times (1/2) \times (1/2) = \mathbf{1/8 = 12,5\%}$.
        (Risque par grossesse d'un garçon atteint = 1/8. Risque par grossesse d'un enfant atteint = 1/8).
    \end{enumerate}
    \textbf{Conclusion :} L'absence de transmission père-fils et la présence de porteuses féminines asymptomatiques confirment le mode récessif lié à l'X.
\end{exoresolu}

\begin{methode}[Identifier et classer les mutations (géniques vs chromosomiques)]
    \textbf{Objectif :} Distinguer les types de mutations à partir de leur description (échelle, mécanisme, conséquence) et prévoir leur impact phénotypique.
    \begin{enumerate}
        \item \textbf{Mutation génique (à l'échelle du gène / nucléotide) :} Modification de la séquence d'ADN d'un gène.
            \begin{itemize}
                \item \textbf{Substitution :} Un nucléotide remplacé par un autre.
                    \begin{itemize}
                        \item Synonyme (silencieuse) : Codon modifié mais même acide aminé (redondance code génétique).
                        \item Non-synonyme (missense) : Changement d'acide aminé (ex. Drépanocytose : $GAG \to GTG$, Glu $\to$ Val).
                        \item Non-sens (nonsense) : Création d'un codon STOP $\to$ protéine tronquée, souvent non fonctionnelle.
                    \end{itemize}
                \item \textbf{Insertion / Délétion (Indel) :} Ajout ou perte de 1, 2... nucléotides.
                    \item Si multiple de 3 : Ajout/perte d'acides aminés (in-frame), protéine souvent fonctionnelle.
                    \item Si non multiple de 3 : \textbf{Mutation à décalage de cadre (frameshift)} $\to$ modification totale de la séquence en aval, codon STOP précoce fréquent $\to$ protéine non fonctionnelle.
            \end{itemize}
        \item \textbf{Mutation chromosomique (à l'échelle du chromosome / génome) :} Modification de la structure ou du nombre de chromosomes. Détectable au caryotype.
            \begin{itemize}
                \item \textbf{Structurales :} Délétion (perte fragment), Duplication (gain), Inversion (renversement), Translocation (échange entre chromosomes non homologues).
                \item \textbf{Numériques (Anomalies de nombre) :} Aneuploïdies (Trisomie 21, 13, 18 ; Syndrome de Turner XO ; Klinefelter XXY). Polyploïdie (Triploïdie 3n - létale chez l'homme).
            \end{itemize}
        \item \textbf{Origine :} Spontanée (erreurs réplication, rayons UV, chimiques) ou Induite (mutagènes). Transmission : Seules les mutations des \textbf{cellules germinales} (gamètes) sont transmissibles à la descendance. Mutations somatiques $\to$ cancers, mosaïcisme, non transmises.
    \end{enumerate}
\end{methode}

\begin{exoresolu}[Classification de mutations et conséquences]
    \textbf{Énoncé :} Classer les mutations suivantes et préciser leur conséquence probable sur la protéine / l'individu.
    \begin{enumerate}
        \item Remplacement d'une base A par une base T dans le codon 6 du gène $\beta$-globine (HBB) : $GAG \to GTG$.
        \item Perte de 2 paires de bases dans l'exon 10 du gène $CFTR$ (fibrose kystique).
        \item Perte d'un fragment du bras court du chromosome 5 (délétion 5p).
        \item Non-disjonction des chromosomes 21 lors de la méiose I chez une femme de 40 ans.
        \item Translocation équilibrée $t(14;21)$ chez un parent sain.
    \end{enumerate}
    \tcblower
    \textbf{Solution détaillée :}
    \begin{enumerate}
        \item \textbf{Mutation génique : Substitution (Transversion A$\to$T).} Conséquence : Missense (non-synonyme). Codon 6 : Glu (Acide glutamique) $\to$ Val (Valine). \textbf{Maladie : Drépanocytose (Anémie falciforme).} Protéine $\beta$-globine mutée (HbS) polymérise en faible $O_2$.
        \item \textbf{Mutation génique : Délétion de 2 pb (Indel).} 2 n'est pas multiple de 3. \textbf{Conséquence : Mutation à décalage de cadre (Frameshift).} Modification de la lecture codante en aval, apparition rapide d'un codon STOP. Protéine CFTR tronquée, non fonctionnelle. \textbf{Maladie : Fibrose kystique (Mucoviscidose).} (Note : $\Delta F508$ est une délétion de 3 pb, in-frame, mais instable).
        \item \textbf{Mutation chromosomique : Structurelle (Délétion partielle).} Perte de matériel génétique (haploinsuffisance pour les gènes de la région). \textbf{Syndrome : Cri du chat (5p-).} Retard mental, cri strident, microcéphalie.
        \item \textbf{Mutation chromosomique : Numérique (Aneuploïdie).} Non-disjonction $\to$ Gamète avec 2 chr 21 (n+1) + Gamète normal (n) $\to$ Zygote trisomique (2n+1). \textbf{Syndrome : Trisomie 21 (Syndrome de Down).} Risque maternel lié à l'âge (cohésion chromatidienne défaillante).
        \item \textbf{Mutation chromosomique : Structurelle (Translocation réciproque équilibrée).} Échange de segments entre chr 14 et 21. Parent sain (matériel génétique complet, juste réorganisé). \textbf{Risque reproductif :} Gamètes déséquilibrés possibles. Risque majoré de trisomie 21 par translocation (chr 21 collé au 14) chez la descendance. Conseil génétique indispensable.
    \end{enumerate}
    \textbf{Conclusion :} L'échelle de la mutation (nucléotide vs chromosome) dicte l'outil de diagnostic (séquençage Sanger/NGS vs Caryotype/FISH/ACPA) et la prévisibilité du phénotype.
\end{exoresolu}

\begin{methode}[Calculer un risque génétique et proposer une prévention (Conseil génétique)]
    \textbf{Objectif :} Synthétiser les savoir-faire pour établir un risque de récidivité et proposer les mesures de prévention adaptées au contexte camerounais.
    \begin{enumerate}
        \item \textbf{Diagnostiquer le mode de transmission :} Pedigree + Biochimie/Moléculaire $\to$ Autosomique Dom/Réc, Gonosomique, Mitochondrial, Multifactoriel.
        \item \textbf{Calculer le risque de récidivité (R) :} Probabilité qu'un couple ait un enfant atteint.
            \begin{itemize}
                \item \textbf{Autosomique Récessif (porteurs sains) :} $R = 1/4$ (si les deux porteurs confirmés). Si statut incertain : $R = P(\text{Mère porteuse}) \times P(\text{Père porteur}) \times 1/4$.
                \item \textbf{Autosomique Dominant (un parent atteint) :} $R = 1/2$. (Attention : néomutation, non-pénétrance, expressivité variable).
                \item \textbf{Gonosomique Récessif (mère porteuse) :} $R(\text{fils atteint}) = 1/4$ (1/2 fils $\times$ 1/2 transmet allèle). $R(\text{fille atteinte}) \approx 0$ (sauf père atteint). $R(\text{enfant atteint}) = 1/4$.
                \item \textbf{Trisomie 21 (âge maternel) :} Risque croissant avec l'âge (ex. 1/1000 à 20 ans, 1/350 à 35 ans, 1/100 à 40 ans).
            \end{itemize}
        \item \textbf{Prévention primaire (avant conception) :}
            \begin{itemize}
                \item Conseil prénuptial / prémarital (obligatoire au Cameroun pour drépanocytose : test électrophorèse Hb).
                \item Dépistage porteurs (famille à risque, populations cibles).
                \item Supplémentation en Acide Folique (0,4 mg/j) $\to$ prévention anomalies tube neural (spina bifida).
                \item Évitement mutagènes (alcool, tabac, rayonnements, certains médicaments tératogènes) pendant la période périconceptionnelle.
            \end{itemize}
        \item \textbf{Prévention secondaire (pendant la grossesse - DPN) :}
            \begin{itemize}
                \item \textbf{Non invasif :} DPNI (ADN fœtal circulant maternel) $\to$ Trisomies 21, 18, 13, sexe. Dépistage combiné (marqueurs sériques + clarté nucale échographie 11-13 SA).
                \item \textbf{Invasif (diagnostic certain) :} Trophobiopsie (11-14 SA), Amniocentèse (15-18 SA), Cordocentèse. Risque fausse couche $\approx$ 0,5-1\%. Caryotype / ACPA / Séquençage.
            \end{itemize}
        \item \textbf{Prévention tertiaire (néonatale) :} Dépistage néonatal systématique (Guthrie) : Phénylcétonurie, Hypothyroïdie congénitale, Drépanocytose (au Cameroun), Fibrose kystique. Prise en charge précoce = pronostic vital/fonctionnel amélioré.
    \end{enumerate}
\end{methode}

\begin{exoresolu}[Conseil génétique : Couple à risque Drépanocytose + Âge maternel]
    \textbf{Énoncé :} Un couple camerounais consulte avant grossesse. L'homme est AS (porteur sain drépanocytose). La femme est AA (non porteuse). La femme a 38 ans. Ils demandent le risque global d'avoir un enfant atteint d'une anomalie génétique grave (Drépanocytose majeure OU Trisomie 21).
    \begin{enumerate}
        \item Calculer le risque de drépanocytose majeure ($SS$).
        \item Estimer le risque de Trisomie 21 lié à l'âge maternel (donnée : 38 ans $\approx$ 1/200).
        \item Calculer le risque global (indépendance des événements).
        \item Proposer la stratégie de prévention adaptée.
    \end{enumerate}
    \tcblower
    \textbf{Solution détaillée :}
    \begin{enumerate}
        \item \textbf{Risque Drépanocytose :} Père $AS$, Mère $AA$.
        Gamètes Père : $A$ (1/2), $S$ (1/2). Gamètes Mère : $A$ (1).
        Descendance : $AA$ (1/2, sain non porteur), $AS$ (1/2, porteur sain).
        \textbf{Risque $SS$ = 0\%.} (Impossible car mère ne donne que $A$).
        \item \textbf{Risque Trisomie 21 :} Âge maternel 38 ans. Risque approximatif à la naissance $\approx$ \textbf{1/200 = 0,5\%}. (Risque au 2ème trimestre plus élevé $\approx$ 1/100, mais avortements spontanés fréquents).
        \item \textbf{Risque global (Anomalie grave) :} Événements indépendants (gènes différents, mécanismes différents).
        $P(\text{Anomalie}) = P(SS) + P(T21) - P(SS \cap T21)$.
        $P(SS) = 0$. $P(T21) = 0,005$.
        \textbf{Risque global = 0,5\% (soit 1 chance sur 200).}
        \item \textbf{Stratégie de prévention :}
        \begin{itemize}
            \item \textbf{Information :} Risque drépanocytose nul pour ce couple précis. Risque Trisomie 21 lié à l'âge, non modifiable.
            \item \textbf{DPNI (Dépistage Prénatal Non Invasif) :} Fortement recommandé dès 10-11 SA (ADN fœtal circulant). Sensibilité Trisomie 21 > 99\%, Spécificité > 99,9\%. Évite ponction invasive initiale.
            \item \textbf{Si DPNI positif :} Confirmation par Amniocentèse (caryotype rapide FISH + ACPA/seq) vers 16 SA.
            \item \textbf{Supplémentation :} Acide Folique 0,4 mg/j (déjà commencé idéalement 3 mois avant conception) pour prévention anomalies tube neural.
            \item \textbf{Suivi échographique :} Échographie morphologique 22 SA (recherche malformations associées).
            \item \textbf{Dépistage néonatal :} Test de Guthrie à J3 (drépanocytose, phénylcétonurie, hypothyroïdie) systématique à la maternité.
        \end{itemize}
    \end{enumerate}
    \textbf{Conclusion :} Le conseil génétique doit être chiffré, non directif, et contextualisé (accès aux techniques DPNI/Amniocentèse à Yaoundé/Douala, prise en charge drépanocytose en centres spécialisés type Centre Pasteur ou Hôpital Général).
\end{exoresolu}

\begin{attention}[Les 5 erreurs qui coûtent des points au Bac]
    \begin{enumerate}
        \item \textbf{Confondre "Brassage interchromosomique" et "Brassage intrachromosomique" :} Ne pas écrire "brassage interchromosomique" pour des gènes liés sur le même chromosome (linkage). Le brassage inter = loi de Mendel 2 (indépendance, 4 gamètes égaux). Le brassage intra = linkage (gènes liés), fréquences déviées, crossing-over possible. \textbf{Piège :} Ne pas préciser "gènes non liés" ou "sur chromosomes différents" pour justifier le 9:3:3:1.
        \item \textbf{Oublier la condition "Test-cross = croisé avec double récessif" :} Dans un test-cross, le parent récessif ne produit qu'UN type de gamète. Le phénotype de la descendance = génotype du gamète de l'autre parent. Erreur fréquente : faire un carré 4x4 pour un test-cross (inutile, 4 cases suffisent) ou oublier que le parent testé peut être homozygote dominant (1 seul phénotype descendant).
        \item \textbf{Erreurs sur le Pédigree Gonosomique :} (1) Faire transmettre l'allèle muté du père à son fils (père donne Y au fils !). (2) Oublier que la fille d'un père atteint (récessif X) est \textbf{obligatoirement porteuse} (reçoit son unique X muté). (3) Calculer le risque "par enfant" au lieu de "par grossesse" ou oublier le sexe de l'enfant (ex: risque fils hémophile = 1/4 par grossesse, pas 1/2).
        \item \textbf{Confondre Mutation Génique et Chromosomique :} Dire "Trisomie 21 est une mutation génique" (FAUX, c'est numérique chromosomique). Dire "Drépanocytose est une mutation chromosomique" (FAUX, c'est substitution ponctuelle, mutation génique). Savoir citer l'outil de détection : Séquençage (génique) vs Caryotype/FISH/ACPA (chromosomique).
        \item \textbf{Calcul de risque : Oublier la probabilité conditionnelle (statut porteur incertain).} Ex: Frère/soeur d'un atteint (récessif) = 2/3 de chance d'être porteur (pas 1/2 car on sait qu'il n'est PAS atteint). Ne pas multiplier les probabilités indépendantes (Père transmet $\times$ Mère transmet $\times$ Sexe enfant). Ne pas mettre d'unité (% ou fraction) ni de phrase de conclusion ("Le risque est de...").
    \end{enumerate}
\end{attention}

\newpage

\coursec{Exercices}

\courssub{Je vérifie mes connaissances}

\begin{exercice}[QCM : Allèles et génotypes]
Parmi les affirmations suivantes, une seule est correcte. Laquelle ?
\begin{enumerate}
    \item Un gène ne possède jamais plus de deux allèles dans une population.
    \item Un individu diploïde homozygote pour un gène autosomique possède deux allèles identiques portés par des chromosomes homologues.
    \item La dominance d'un allèle sur un autre modifie la fréquence de cet allèle dans la population.
    \item Le brassage interchromosomique concerne uniquement les gènes situés sur le même chromosome.
\end{enumerate}
\end{exercice}

\begin{exercice}[Vrai ou Faux justifié : Mécanismes du brassage]
Dire si les affirmations suivantes sont vraies ou fausses en justifiant brièvement votre réponse.
\begin{enumerate}
    \item Lors de la méiose, la ségrégation des allèles d'un même gène (loi de la pureté des gamètes) a lieu au cours de l'anaphase I.
    \item Le crossing-over est la cause unique du brassage interchromosomique.
    \item Un taux de recombinaison de 50 \% entre deux gènes prouve obligatoirement qu'ils sont situés sur des chromosomes différents.
    \item Chez l'homme, le déterminisme du groupe sanguin ABO est un exemple de codominance et d'allèles multiples.
\end{enumerate}
\end{exercice}

\begin{exercice}[Définitions et notations]
\begin{enumerate}
    \item Définir les termes suivants : \cle{allèle}, \cle{locus}, \cle{génotype}, \cle{phénotype}, \cle{test-cross}.
    \item Chez la Drosophile, la couleur des yeux est un caractère lié au chromosome X. L'allèle $X^R$ (yeux rouges) est dominant sur l'allèle $X^r$ (yeux blancs). Écrire le génotype d'une femelle aux yeux rouges homozygote et celui d'un mâle aux yeux blancs. Préciser la composition en allèles des gamètes produits par chacun.
\end{enumerate}
\end{exercice}

\begin{exercice}[Calculs directs : Probabilités et taux de recombinaison]
\begin{enumerate}
    \item Un individu de génotype $A/a ; B/b$ (gènes non liés, sur chromosomes différents) produit des gamètes. Calculer la probabilité d'obtenir un gamète de génotype $A ; b$.
    \item Chez le maïs, deux gènes liés $G$ (graines lisses, dominant) et $C$ (graines colorées, dominant) sont séparés par un taux de recombinaison de \SI{18}{\percent}. Un individu doublement hétérozygote en \cle{phase de couplage} ($G\;C // g\;c$) est croisé avec un double récessif ($g\;c // g\;c$). Calculer les fréquences phénotypiques attendues dans la descendance.
\end{enumerate}
\end{exercice}

\courssub{Je m'entraîne}

\begin{exercice}[Monohybridisme et codominance : Groupes sanguins ABO]
Un homme de groupe sanguin A épouse une femme de groupe sanguin B. Ils ont un enfant de groupe O.
\begin{enumerate}
    \item Déterminer le génotype précis de chacun des deux parents.
    \item Lister les phénotypes (groupes sanguins) possibles pour leurs enfants à venir et donner la probabilité de chaque phénotype.
    \item Cet enfant (groupe O) peut-il recevoir une transfusion sanguine de son père ? De sa mère ? Justifier en précisant les agglutinines présentes dans le sérum du receveur et les agglutinogènes sur les hématies du donneur.
\end{enumerate}
\end{exercice}

\begin{exercice}[Dihybridisme et test-cross : Indépendance des caractères]
Chez la souris, la couleur du pelage est déterminée par un gène autosomique ($N$ = noir dominant sur $n$ = blanc) et la longueur des poils par un autre gène autosomique ($L$ = poils longs dominant sur $l$ = poils courts). Les deux gènes sont non liés (sur chromosomes différents).
On croise une souris noire à poils longs (issue d'une lignée pure) avec une souris blanche à poils courts. Tous les individus F1 sont noirs à poils longs.
On réalise ensuite un test-cross : on croise une femelle F1 avec un mâle blanc à poils courts.
\begin{enumerate}
    \item Écrire les génotypes des parents P, des F1 et du mâle testeur.
    \item Déterminer les phénotypes et leurs fréquences théoriques attendus dans la descendance du test-cross.
    \item Si l'on observe 240 descendants au total, calculer les effectifs théoriques pour chaque classe phénotypique.
\end{enumerate}
\end{exercice}

\begin{exercice}[Linkage partiel : Cartographie génétique chez la Drosophile]
On étudie deux gènes chez \emph{Drosophila melanogaster} : $e$ (corps ébène, récessif) et $v$ (ailes vestigiales, récessif). Les allèles sauvages sont $e^+$ (corps gris) et $v^+$ (ailes normales).
Une femelle doublement hétérozygote en phase de couplage ($e^+\;v^+ // e\;v$) est croisée avec un mâle double mutant ($e\;v // e\;v$).
On obtient la descendance suivante (effectifs) :
\begin{center}
\begin{tabular}{|l|c|}\hline
\rowcolor{popblueL} Phénotype & Effectif \\ \hline
Corps gris, ailes normales & 412 \\
Corps ébène, ailes vestigiales & 398 \\
Corps gris, ailes vestigiales & 88 \\
Corps ébène, ailes normales & 92 \\ \hline
Total & 990 \\ \hline
\end{tabular}
\end{center}
\begin{enumerate}
    \item Identifier les classes parentales et les classes recombinées. En déduire si les gènes sont liés ou non.
    \item Calculer le taux de recombinaison (fréquence de crossing-over) entre les gènes $e$ et $v$. Exprimer le résultat en \cle{centiMorgans (cM)}.
    \item Représenter la carte factorielle partielle correspondante.
\end{enumerate}
\end{exercice}

\begin{popfigure}
\begin{tikzpicture}[
    node distance=1.5cm and 1cm,
    male/.style={draw, rectangle, minimum width=1cm, minimum height=1cm, thick},
    female/.style={draw, circle, minimum width=1cm, minimum height=1cm, thick},
    affected/.style={fill=popred!80},
    carrier/.style={fill=poporange!50, pattern=north east lines, pattern color=poporange},
    normal/.style={fill=popgreen!20},
    marriage/.style={thick},
    descent/.style={thick}
]
% Generation I (Grand-parents)
\node[female, normal] (GM1) at (0,3) {};
\node[male, normal] (GF1) at (2,3) {};
\node[female, affected] (GM2) at (5,3) {};
\node[male, normal] (GF2) at (7,3) {};

% Marriage lines Gen I
\draw[marriage] (GM1) -- (GF1);
\draw[marriage] (GM2) -- (GF2);

% Generation II (Parents)
\node[female, carrier] (M) at (1,1.5) {};
\node[male, carrier] (F) at (6,1.5) {};

% Descent lines Gen I -> Gen II
\draw[descent] (1,2.5) -- (1,1.5); % from GF1/GM1 center approx
\draw[descent] (6,2.5) -- (6,1.5); % from GF2/GM2 center approx

% Marriage Gen II
\draw[marriage] (M) -- (F);

% Generation III (Enfants)
\node[male, normal] (C1) at (3.5,0) {};
\node[male, normal] (C2) at (5,0) {};
\node[female, affected] (C3) at (6.5,0) {};

% Descent lines Gen II -> Gen III
\draw[descent] (3.5,0.75) -- (3.5,0);
\draw[descent] (5,0.75) -- (5,0);
\draw[descent] (6.5,0.75) -- (6.5,0);

% Labels
\node[above=2pt] at (GM1.north) {I-1};
\node[above=2pt] at (GF1.north) {I-2};
\node[above=2pt] at (GM2.north) {I-3};
\node[above=2pt] at (GF2.north) {I-4};
\node[left=2pt] at (M.west) {II-1 (Mère)};
\node[right=2pt] at (F.east) {II-2 (Père)};
\node[below=2pt] at (C1.south) {III-1};
\node[below=2pt] at (C2.south) {III-2};
\node[below=2pt] at (C3.south) {III-3 (SS)};

% Legend
\node[anchor=north west, align=left, font=\small] at (-2,-1.5) {
    \textbf{Légende :}\\
    \tikz\draw[male, normal] (0,0) rectangle (0.3,0.3); \hspace{2pt} Homme sain (AA ou AS)\\
    \tikz\draw[female, normal] (0,0) circle (0.15); \hspace{2pt} Femme saine (AA ou AS)\\
    \tikz\draw[female, affected] (0,0) circle (0.15); \hspace{2pt} Femme atteinte (SS)\\
    \tikz\draw[female, carrier] (0,0) circle (0.15); \hspace{2pt} Porteuse saine (AS - déduit)\\
};
\end{tikzpicture}
\legende{Pedigree de la famille camerounaise (Exercice Drépanocytose). Les symboles hachurés indiquent les porteurs sains (AS) déduits du raisonnement génétique. La fille III-3 est atteinte (SS, symbole plein).}
\end{popfigure}

\begin{exercice}[Pedigree : Drépanocytose au Cameroun]
Dans une famille camerounaise, les deux parents sont phénotypiquement normaux (hémoglobine A). Ils ont déjà trois enfants : deux garçons normaux et une fille atteinte de drépanocytose majeure (hémoglobine S homozygote, SS). Le père a une sœur atteinte (décédée jeune) ; la mère a un frère sain.
\begin{enumerate}
    \item Rappeler le mode de transmission de la drépanocytose (autosomique récessif). Écrire les allèles en jeu : $A$ (normal) et $S$ (drépanocytaire).
    \item Déterminer le génotype des deux parents. Justifier.
    \item Le pedigree de cette famille sur trois générations est représenté ci-dessus. Vérifier qu'il correspond à l'énoncé et compléter la numérotation des individus (générations I, II, III).
    \item Calculer le risque (probabilité) pour qu'un quatrième enfant soit atteint (SS). Calculer le risque pour qu'il soit porteur sain (AS).
    \item Proposer une mesure de prévention adaptée au contexte camerounais pour ce couple s'ils envisagent une nouvelle grossesse (dépistage prénatal, conseil génétique, prise en charge néonatale).
\end{enumerate}
\end{exercice}

\courssub{Je me prépare au Bac}

\begin{exercice}[Type Bac : Analyse de croisements chez la Drosophile — Linkage et carte factorielle]
\textbf{Contexte :} Un généticien étudie trois gènes chez \emph{Drosophila melanogaster} :
\begin{itemize}
    \item $y$ : corps jaune (récessif) / $y^+$ : corps gris (sauvage)
    \item $w$ : yeux blancs (récessif) / $w^+$ : yeux rouges (sauvage)
    \item $m$ : ailes miniatures (récessif) / $m^+$ : ailes normales (sauvage)
\end{itemize}
Ces trois gènes sont situés sur le chromosome X.

\textbf{Expérience 1 :} On croise une femelle triple hétérozygote ($y^+\;w^+\;m^+ // y\;w\;m$) avec un mâle triple mutant ($y\;w\;m // Y$). On observe 1500 descendants mâles (hémizygotes pour le chromosome X maternel) répartis ainsi :
\begin{center}
\begin{tabular}{|l|c|}\hline
\rowcolor{popblueL} Phénotype (mâles) & Effectif \\ \hline
Gris, rouges, normales ($y^+\;w^+\;m^+$) & 580 \\
Jaune, blancs, miniatures ($y\;w\;m$) & 592 \\
Gris, blancs, normales ($y^+\;w\;m^+$) & 42 \\
Jaune, rouges, miniatures ($y\;w^+\;m$) & 38 \\
Gris, rouges, miniatures ($y^+\;w^+\;m$) & 124 \\
Jaune, blancs, normales ($y\;w\;m^+$) & 118 \\
Gris, blancs, miniatures ($y^+\;w\;m$) & 3 \\
Jaune, rouges, normales ($y\;w^+\;m^+$) & 3 \\ \hline
Total & 1500 \\ \hline
\end{tabular}
\end{center}

\textbf{Expérience 2 :} On croise une femelle double hétérozygote pour $y$ et $w$ ($y^+\;w^+ // y\;w$) avec un mâle double mutant ($y\;w // Y$). On observe un taux de recombinaison de \SI{1,5}{\percent} entre $y$ et $w$.

\textbf{Questions :}
\begin{enumerate}
    \item À partir des résultats de l'Expérience 1, identifier les classes parentales et les classes recombinées simples et doubles. En déduire l'ordre des trois gènes sur le chromosome X.
    \item Calculer les taux de recombinaison (en \%) entre $y$ et $w$, entre $w$ et $m$, et entre $y$ et $m$. Exprimer les distances en cM.
    \item Vérifier la cohérence de vos résultats avec la formule de la fonction de mapping (absence d'interférence supposée) : $RF_{y-m} \approx RF_{y-w} + RF_{w-m} - 2 \times RF_{y-w} \times RF_{w-m}$.
    \item Établir la carte factorielle des trois gènes en indiquant les distances en cM.
    \item Expliquer pourquoi l'Expérience 2 ne permet pas de déterminer l'ordre des gènes par rapport à $m$.
\end{enumerate}
\end{exercice}

\begin{popfigure}
\begin{tikzpicture}[scale=0.8, >=Stealth]
% Gene order: y -- w -- m
\draw[thick, popblue] (0,0) -- (10,0);
% y at 0
\node[below=3pt] at (0,0) {$y$};
\node[above=5pt] at (0,0) {\textbf{0,0 cM}};
% w at 5.73 cM -> scaled ~2.86 cm
\draw[thick, popblue] (2.86, -0.1) -- (2.86, 0.1);
\node[below=3pt] at (2.86,0) {$w$};
\node[above=5pt] at (2.86,0) {\textbf{5,7 cM}};
% m at 22.26 cM -> scaled ~11.13 cm (too wide). Rescale.
% Let's rescale: 1 cM = 0.3 cm. Total ~6.7 cm.
\draw[thick, popblue] (0,0) -- (6.7,0);
\node[below=3pt] at (0,0) {$y$};
\node[above=5pt] at (0,0) {\textbf{0,0 cM}};
\draw[thick, popblue] (1.72, -0.1) -- (1.72, 0.1); % 5.73 * 0.3
\node[below=3pt] at (1.72,0) {$w$};
\node[above=5pt] at (1.72,0) {\textbf{5,7 cM}};
\draw[thick, popblue] (6.68, -0.1) -- (6.68, 0.1); % 22.26 * 0.3
\node[below=3pt] at (6.68,0) {$m$};
\node[above=5pt] at (6.68,0) {\textbf{22,3 cM}};
% Distance labels
\node[above=12pt] at (0.86,0) {5,7 cM};
\node[above=12pt] at (4.2,0) {16,5 cM};
\node[above=22pt] at (3.34,0) {22,3 cM (y-m)};
\end{tikzpicture}
\legende{Carte factorielle des gènes $y$, $w$, $m$ sur le chromosome X de la Drosophile (Expérience 1). L'ordre est $y - w - m$. Distances calculées à partir des fréquences de recombinaison observées.}
\end{popfigure}

\begin{exercice}[Type Bac : Hérédité gonosomique — Hémophilie A et conseil génétique]
\textbf{Contexte :} L'hémophilie A est une maladie récessive liée au chromosome X (gène $F8$, allèle muté $h$, allèle sain $H$). Au Cameroun, la prévalence est estimée à 1/5000 naissances mâles.
On considère la famille suivante :
\begin{itemize}
    \item Génération I : Grand-père (I-1) hémophile ; Grand-mère (I-2) phénotypiquement normale, sans antécédent familial connu.
    \item Génération II : Ils ont une fille (II-1) phénotypiquement normale et un fils (II-2) sain.
    \item Génération III : La fille (II-1) épouse un homme sain (II-3). Ils ont deux enfants : un garçon (III-1) hémophile et une fille (III-2) phénotypiquement normale.
\end{itemize}
\textbf{Questions :}
\begin{enumerate}
    \item Construire le pedigree complet de cette famille sur trois générations. Indiquer les génotypes probables de chaque individu (utiliser la notation $X^H$, $X^h$, $Y$).
    \item Justifier le génotype de la grand-mère (I-2) et de la mère (II-1).
    \item Calculer la probabilité que la fille (III-2) soit porteuse de l'allèle $h$.
    \item La fille (III-2), maintenant adulte, envisage une grossesse avec un homme sain. Calculer le risque (probabilité) qu'elle ait :
    \begin{enumerate}
        \item un garçon hémophile ;
        \item une fille porteuse ;
        \item un enfant (garçon ou fille) atteint ou porteur.
    \end{enumerate}
    \item En tant que conseiller génétique dans un hôpital de district au Cameroun, rédiger un court message (3-4 lignes) à destination de ce couple pour expliquer le risque et les options de diagnostic prénatal disponibles (ex. : PCR sur ADN fœtal, prélèvement de villosités choriales).
\end{enumerate}
\end{exercice}

\begin{popfigure}
\begin{tikzpicture}[
    node distance=1.5cm and 1cm,
    male/.style={draw, rectangle, minimum width=1cm, minimum height=1cm, thick},
    female/.style={draw, circle, minimum width=1cm, minimum height=1cm, thick},
    affected/.style={fill=popred!80},
    carrier/.style={fill=poporange!50, pattern=north east lines, pattern color=poporange},
    normal/.style={fill=popgreen!20},
    marriage/.style={thick},
    descent/.style={thick}
]
% Gen I
\node[male, affected] (I1) at (0,3) {I-1};
\node[female, normal] (I2) at (2,3) {I-2};
\draw[marriage] (I1) -- (I2);

% Gen II
\node[female, carrier] (II1) at (0,1.5) {II-1};
\node[male, normal] (II2) at (2,1.5) {II-2};
\draw[descent] (1,2.5) -- (1,1.5); % center
\draw[descent] (1,2.5) -- (1,1.5); % dummy for alignment
% Actually draw from middle of marriage line
\draw[descent] (1,3) -- (0,1.5);
\draw[descent] (1,3) -- (2,1.5);

% Gen II Marriage
\node[male, normal] (II3) at (5,1.5) {II-3};
\draw[marriage] (II1) -- (II3);

% Gen III
\node[male, affected] (III1) at (4,0) {III-1};
\node[female, carrier] (III2) at (6,0) {III-2};
\draw[descent] (5,0.75) -- (4,0);
\draw[descent] (5,0.75) -- (6,0);

% Genotypes labels
\node[left=5pt, font=\small] at (I1.west) {$X^h Y$};
\node[right=5pt, font=\small] at (I2.east) {$X^H X^H$};
\node[left=5pt, font=\small] at (II1.west) {$X^H X^h$};
\node[right=5pt, font=\small] at (II2.east) {$X^H Y$};
\node[right=5pt, font=\small] at (II3.east) {$X^H Y$};
\node[left=5pt, font=\small] at (III1.west) {$X^h Y$};
\node[right=5pt, font=\small] at (III2.east) {$X^H X^h$ (1/2) ou $X^H X^H$ (1/2)};

% Legend
\node[anchor=north west, align=left, font=\small] at (-3,-1) {
    \textbf{Légende :}\\
    \tikz\draw[male, affected] (0,0) rectangle (0.3,0.3); \hspace{2pt} Homme atteint ($X^h Y$)\\
    \tikz\draw[male, normal] (0,0) rectangle (0.3,0.3); \hspace{2pt} Homme sain ($X^H Y$)\\
    \tikz\draw[female, carrier] (0,0) circle (0.15); \hspace{2pt} Femme porteuse ($X^H X^h$)\\
    \tikz\draw[female, normal] (0,0) circle (0.15); \hspace{2pt} Femme saine non porteuse ($X^H X^H$)\\
};
\end{tikzpicture}
\legende{Pedigree de la famille avec hémophilie A (liée à l'X). La grand-mère I-2 est homozygote normale ($X^H X^H$) car sans antécédent ; la mère II-1 est porteuse obligée ($X^H X^h$) car fille de père atteint. La fille III-2 a 1/2 chance d'être porteuse.}
\end{popfigure}

\begin{exercice}[Type Bac : Mutations, drépanocytose et prévention en zone d'endémie palustre]
\textbf{Document 1 : Séquence d'ADN codant pour la chaîne $\beta$-globine (début de l'exon 1, brin codant).}
\begin{itemize}
    \item Allèle normal ($A$) : 5' - \textbf{CTG ACT CCT GAG GAG AAG TCT} - 3'
    \item Allèle drépanocytaire ($S$) : 5' - \textbf{CTG ACT CCT \underline{GTG} GAG AAG TCT} - 3'
\end{itemize}
\textbf{Document 2 : Caryotype d'un individu atteint de trisomie 21 (47,XX,+21).}
\textbf{Document 3 : Résultats d'une électrophorèse d'hémoglobine sur gel d'acétate de cellulose (pH 8,6).}
Migration vers l'anode (+) : Hb A (rapide) > Hb F > Hb S (lente) > Hb C (très lente).
Un échantillon montre deux bandes : une à la position Hb A, une à la position Hb S.

\textbf{Questions :}
\begin{enumerate}
    \item \textbf{Analyse moléculaire :} En utilisant le code génétique (fourni ci-dessous), traduire les deux séquences d'ADN (brin codant) en séquence d'acides aminés. Identifier le type de mutation (substitution, insertion, délétion) et sa conséquence sur la protéine (synonyme, non-sens, faux-sens). Préciser le changement d'acide aminé (position 6).
    \item \textbf{Analyse cytogénétique :} Classer la mutation illustrée par le Document 2 (mutation chromosomique : type et conséquence sur le nombre de chromosomes). Donner la formule chromosomique d'un gamète produit par cet individu (méiose) qui aboutirait à un enfant trisomique après fécondation par un gamète normal.
    \item \textbf{Interprétation phénotypique :} À partir du Document 3, déterminer le phénotype (génotype à l'échelle de l'hémoglobine) de l'individu testé. Est-il malade ? Est-il porteur ? Justifier.
    \item \textbf{Épidémiologie et prévention :} Au Cameroun, la fréquence de l'allèle $S$ atteint 10 à 25 \% selon les régions (forte prévalence dans le Nord et l'Extrême-Nord, zone de savane soudano-sahélienne).
    \begin{enumerate}
        \item Expliquer le maintien de cette fréquence élevée par la sélection naturelle (avantage sélectif des hétérozygotes $AS$ face au paludisme à \emph{Plasmodium falciparum}).
        \item Proposer une stratégie nationale de prévention de la drépanocytose majeure (SS) adaptée au système de santé camerounais (dépistage néonatal, conseil génétique prénuptial, prise en charge précoce : pénicilline V, acide folique, vaccination anti-pneumocoque, hydroxyurée). Citer le rôle du \emph{Centre de Prise en Charge de la Drépanocytose} (ex. : Hôpital Laquintinie Douala, Hôpital Central Yaoundé).
    \end{enumerate}
\end{enumerate}

\textbf{Code génétique (extraits) :}
\begin{center}
\begin{tabular}{|c|c|c|c|}\hline
\rowcolor{popblueL} Codon & AA & Codon & AA \\ \hline
GAG / GAA & Glu (E) & GTG / GTA / GTC / GTT & Val (V) \\ \hline
CTG / CTC / CTT / CTA & Leu (L) & ACT / ACC / ACA / ACG & Thr (T) \\ \hline
CCT / CCC / CCA / CCG & Pro (P) & AAG / AAA & Lys (K) \\ \hline
TCT / TCC / TCA / TCG & Ser (S) & & \\ \hline
\end{tabular}
\end{center}
\end{exercice}

\begin{aretenir}[Résumé des formules clés pour le Bac]
\begin{itemize}
    \item \textbf{Taux de recombinaison (RF) :} $RF = \frac{\text{Nombre de recombinants}}{\text{Total des descendants}} \times 100$ (en \%) ou en cM (1\% = 1 cM).
    \item \textbf{Test-cross (dihybridisme lié) :} Phénotypes parentaux = majoritaires ; Phénotypes recombinés = minoritaires.
    \item \textbf{Fonction de mapping (Haldane, sans interférence) :} $RF_{AC} = RF_{AB} + RF_{BC} - 2 \times RF_{AB} \times RF_{BC}$.
    \item \textbf{Risque génétique (autosomique récessif, parents hétérozygotes) :} Enfant atteint = 1/4 ; Porteur sain = 1/2 ; Sain non porteur = 1/4.
    \item \textbf{Risque génétique (lié à l'X récessif, mère porteuse) :} Fils atteint = 1/2 ; Fille porteuse = 1/2.
\end{itemize}
\end{aretenir}

\begin{attention}[Pièges fréquents à éviter]
\begin{itemize}
    \item \textbf{Notation Drosophile :} Ne pas écrire $R/r$ pour la couleur des yeux (gène sur X). Utiliser $X^R/X^r$ (femelle) et $X^r/Y$ (mâle).
    \item \textbf{Taux de recombinaison 50\% :} Ne pas conclure « gènes sur chromosomes différents ». Conclure : « gènes non liés (indépendants) : soit sur chromosomes différents, soit éloignés sur le même chromosome ».
    \item \textbf{Pedigree :} Ne pas oublier de numéroter les individus par génération (I-1, I-2, II-1...). Un individu atteint (homozygote récessif) a obligatoirement deux parents porteurs (hétérozygotes) pour un caractère autosomique récessif.
    \item \textbf{Trisomie 21 :} La non-disjonction peut avoir lieu en anaphase I (séparation homologues) ou anaphase II (séparation chromatides). Le gamète « diploïde » pour le chr 21 a 24 chromosomes (23 + 1 chr 21 en double).
    \item \textbf{Électrophorèse Hb :} Hb A migre plus vite (plus chargé négativement) que Hb S. Bande unique Hb S = SS (malade). Bandes Hb A + Hb S = AS (porteur sain).
\end{itemize}
\end{attention}

\newpage

\coursec{Corrigés des exercices}

\begin{corrige}[Exercice 1 — QCM : Allèles et génotypes]
\textbf{Réponse correcte : b.}

\begin{itemize}
    \item \textbf{a. Fausse.} Un gène peut posséder plus de deux allèles dans une population : c'est le cas de l'\cle{allélisme multiple}. Exemple : le gène du système ABO possède trois allèles ($I^A$, $I^B$, $i$). En revanche, un \emph{individu} diploïde ne porte que deux allèles de ce gène (un sur chaque chromosome homologue).
    \item \textbf{b. Vraie.} Chez un individu diploïde, les deux allèles d'un gène autosomique occupent le même \cle{locus} sur deux chromosomes homologues. S'ils sont identiques, l'individu est \cle{homozygote} ; s'ils sont différents, il est \cle{hétérozygote}.
    \item \textbf{c. Fausse.} La dominance est une relation \emph{phénotypique} entre allèles chez l'hétérozygote ; elle ne confère aucun avantage de transmission : un allèle dominant se transmet avec la même probabilité (1/2) qu'un allèle récessif. La fréquence allélique dans une population dépend de la sélection naturelle, de la dérive génétique, des mutations et des migrations, pas de la dominance.
    \item \textbf{d. Fausse.} C'est exactement l'inverse : le brassage \emph{inter}chromosomique concerne des gènes situés sur des chromosomes \emph{différents} (disjonction indépendante des paires d'homologues en anaphase I). Le brassage \emph{intra}chromosomique (crossing-over) concerne les gènes situés sur le \emph{même} chromosome.
\end{itemize}

\begin{attention}[Point de vigilance]
Distinguer « nombre d'allèles dans la population » (potentiellement illimité, $\geq 2$) et « nombre d'allèles chez un individu diploïde » (exactement 2 par gène autosomique, sauf les gènes portés par le chromosome X chez l'homme, qui est hémizygote).
\end{attention}
\end{corrige}

\begin{corrige}[Exercice 2 — Vrai ou Faux : Mécanismes du brassage]
\begin{enumerate}
    \item \textbf{Vrai.} En anaphase I de méiose, les chromosomes homologues (chacun formé de deux chromatides sœurs) se séparent et migrent vers des pôles opposés. Les deux allèles d'un même gène, portés par ces homologues, sont ainsi séparés dans des cellules différentes : chaque gamète ne reçoit qu'\emph{un} allèle par gène. C'est la base cytologique de la première loi de Mendel (pureté des gamètes).
    \item \textbf{Faux.} Le crossing-over est le mécanisme du brassage \emph{intra}chromosomique (échange réciproque de segments entre chromatides non-sœurs de chromosomes homologues, en prophase I, au niveau des chiasmas). Le brassage \emph{inter}chromosomique résulte de la \emph{disjonction indépendante} des paires de chromosomes homologues en anaphase I (orientation aléatoire des bivalents en métaphase I).
    \item \textbf{Faux.} Un taux de recombinaison de \SI{50}{\percent} signifie que les deux gènes se comportent comme s'ils étaient indépendants. Deux interprétations possibles : (i) ils sont sur des chromosomes différents ; (ii) ils sont sur le même chromosome mais très éloignés l'un de l'autre (les crossing-over sont alors si fréquents entre eux que les gamètes recombinés atteignent 50 \%). Un taux de 50 \% ne permet donc pas de trancher entre ces deux situations.
    \item \textbf{Vrai.} Le gène ABO présente trois allèles : $I^A$, $I^B$ (codominants entre eux : l'hétérozygote $I^A/I^B$ est de groupe AB, les deux agglutinogènes A et B étant exprimés à la surface des hématies) et $i$ (récessif devant les deux autres). C'est le cas classique combinant allélisme multiple et codominance.
\end{enumerate}
\end{corrige}

\begin{corrige}[Exercice 3 — Définitions et notations]
\textbf{1. Définitions.}
\begin{itemize}
    \item \cle{Allèle} : version (forme) d'un gène, occupant un locus donné ; les allèles d'un même gène diffèrent par leur séquence nucléotidique et peuvent déterminer des versions différentes du caractère.
    \item \cle{Locus} : emplacement précis et constant d'un gène sur un chromosome.
    \item \cle{Génotype} : ensemble des allèles portés par un individu pour le ou les gènes considérés (ex. $A/a$).
    \item \cle{Phénotype} : manifestation observable du génotype, résultant de l'expression des gènes dans un environnement donné (ex. [graines lisses]).
    \item \cle{Test-cross} (croisement-test) : croisement d'un individu de phénotype dominant (génotype inconnu) avec un individu homozygote récessif (testeur) ; l'analyse des proportions de la descendance permet de déterminer le génotype de l'individu testé et, pour deux gènes, de révéler un éventuel linkage et de mesurer le taux de recombinaison.
\end{itemize}

\textbf{2. Génotypes et gamètes.}
\begin{itemize}
    \item Femelle aux yeux rouges homozygote : $X^R / X^R$. Gamètes : 100 \% d'ovules portant $X^R$.
    \item Mâle aux yeux blancs : $X^r / Y$ (hémizygote : un seul allèle, porté par le X ; le chromosome Y ne porte pas ce gène). Gamètes : 50 \% de spermatozoïdes $X^r$ et 50 \% de spermatozoïdes $Y$.
\end{itemize}

\begin{attention}[Point de vigilance]
Chez la Drosophile comme chez l'Homme, le mâle est hémizygote pour les gènes portés par X : il exprime son unique allèle, même s'il est « récessif ». On écrit $X^r/Y$ et jamais $r/r$.
\end{attention}
\end{corrige}

\begin{corrige}[Exercice 4 — Calculs directs : Probabilités et taux de recombinaison]
\textbf{1. Probabilité du gamète $A ; b$.}

Les gènes étant indépendants : $P(A) = \dfrac{1}{2}$ (ségrégation du couple $A/a$) et $P(b) = \dfrac{1}{2}$ (ségrégation du couple $B/b$). Les deux événements sont indépendants :
\[ P(A ; b) = \frac{1}{2} \times \frac{1}{2} = \frac{1}{4} = 0,25 \; \text{soit} \; \SI{25}{\percent}. \]
L'individu $A/a ; B/b$ produit en fait 4 types de gamètes équiprobables : $A;B$, $A;b$, $a;B$, $a;b$ (1/4 chacun), par brassage interchromosomique.

\textbf{2. Test-cross avec gènes liés (phase de couplage).}

L'individu $G\;C // g\;c$ produit :
\begin{itemize}
    \item gamètes \textbf{parentaux} : $G\;C$ et $g\;c$, au total $100 - 18 = \SI{82}{\percent}$, soit \SI{41}{\percent} chacun ;
    \item gamètes \textbf{recombinés} (par crossing-over) : $G\;c$ et $g\;C$, au total \SI{18}{\percent}, soit \SI{9}{\percent} chacun.
\end{itemize}
Le testeur $g\;c // g\;c$ ne produit que des gamètes $g\;c$ : le phénotype des descendants révèle directement le gamète reçu du parent hétérozygote.

\begin{center}
\begin{tabular}{|l|c|c|}\hline
\rowcolor{popblueL} Phénotype de la descendance & Génotype & Fréquence \\ \hline
Graines lisses, colorées & $G\;C // g\;c$ & \SI{41}{\percent} \\
Graines rugueuses, incolores & $g\;c // g\;c$ & \SI{41}{\percent} \\
Graines lisses, incolores & $G\;c // g\;c$ & \SI{9}{\percent} \\
Graines rugueuses, colorées & $g\;C // g\;c$ & \SI{9}{\percent} \\ \hline
\end{tabular}
\end{center}

Vérification : $41 + 41 + 9 + 9 = \SI{100}{\percent}$. Les deux classes parentales sont majoritaires et équifréquentes, les deux classes recombinées minoritaires et équifréquentes : c'est la signature d'un \cle{linkage partiel}.
\end{corrige}

\begin{corrige}[Exercice 5 — Monohybridisme et codominance : Groupes sanguins ABO]
\textbf{1. Génotypes des parents.}

L'enfant est de groupe O, donc de génotype $i/i$ : il a reçu un allèle $i$ de \emph{chaque} parent. Le père (groupe A) possède au moins un allèle $I^A$ ; la mère (groupe B) au moins un allèle $I^B$. Donc :
\[ \text{Père : } I^A/i \qquad \text{Mère : } I^B/i. \]

\textbf{2. Phénotypes possibles et probabilités.}

Gamètes du père : $I^A$ (1/2), $i$ (1/2). Gamètes de la mère : $I^B$ (1/2), $i$ (1/2). Échiquier de croisement :

\begin{center}
\begin{tabular}{|c|c|c|}\hline
\rowcolor{popblueL} & $I^B$ (1/2) & $i$ (1/2) \\ \hline
$I^A$ (1/2) & $I^A/I^B$ : groupe AB (1/4) & $I^A/i$ : groupe A (1/4) \\ \hline
$i$ (1/2) & $I^B/i$ : groupe B (1/4) & $i/i$ : groupe O (1/4) \\ \hline
\end{tabular}
\end{center}

Les quatre groupes sont possibles, chacun avec la probabilité $1/4 = \SI{25}{\percent}$.

\textbf{3. Compatibilité transfusionnelle.}

L'enfant O possède des hématies sans agglutinogène (ni A ni B) mais son sérum contient les \textbf{deux agglutinines} anti-A \emph{et} anti-B.
\begin{itemize}
    \item \textbf{Sang du père (A)} : hématies portant l'agglutinogène A $\Rightarrow$ elles seraient agglutinées par les anti-A du receveur. \textbf{Transfusion interdite.}
    \item \textbf{Sang de la mère (B)} : hématies portant l'agglutinogène B $\Rightarrow$ agglutination par les anti-B du receveur. \textbf{Transfusion interdite.}
\end{itemize}
L'enfant O ne peut recevoir que du sang de groupe O (donneur universel pour les hématies). C'est la règle appliquée dans les banques de sang des hôpitaux camerounais : on vérifie toujours la compatibilité ABO \emph{et} Rhésus avant toute transfusion.
\end{corrige}

\begin{corrige}[Exercice 6 — Dihybridisme et test-cross : Indépendance des caractères]
\textbf{1. Génotypes.}
\begin{itemize}
    \item Parent P (lignée pure, noire à poils longs) : $N/N ; L/L$ — gamètes $N;L$.
    \item Parent P (blanche à poils courts) : $n/n ; l/l$ — gamètes $n;l$.
    \item F1 : $N/n ; L/l$ (double hétérozygote, phénotype [noir, poils longs], ce qui confirme la dominance de $N$ sur $n$ et de $L$ sur $l$).
    \item Mâle testeur (blanc, poils courts) : $n/n ; l/l$ — gamètes $n;l$ uniquement.
\end{itemize}

\textbf{2. Résultats théoriques du test-cross.}

Les gènes étant indépendants, la femelle F1 produit 4 types de gamètes équiprobables (brassage interchromosomique) : $N;L$, $N;l$, $n;L$, $n;l$ (1/4 chacun). Le testeur ne fournissant que $n;l$, les phénotypes de la descendance reflètent directement les gamètes de la F1 :

\begin{center}
\begin{tabular}{|l|l|c|}\hline
\rowcolor{popblueL} Génotype & Phénotype & Fréquence \\ \hline
$N/n ; L/l$ & noir, poils longs & 1/4 \\
$N/n ; l/l$ & noir, poils courts & 1/4 \\
$n/n ; L/l$ & blanc, poils longs & 1/4 \\
$n/n ; l/l$ & blanc, poils courts & 1/4 \\ \hline
\end{tabular}
\end{center}

\textbf{3. Effectifs théoriques pour 240 descendants.}
\[ 240 \times \frac{1}{4} = 60 \text{ individus par classe}. \]
Soit 60 noirs à poils longs, 60 noirs à poils courts, 60 blancs à poils longs, 60 blancs à poils courts.

\begin{attention}[Point de vigilance]
Un test-cross donnant 4 classes équiprobables (1/4 ; 1/4 ; 1/4 ; 1/4) est la signature du \textbf{brassage interchromosomique} (gènes indépendants). Si les gènes étaient liés, on obtiendrait 2 classes majoritaires (parentales) et 2 classes minoritaires (recombinées).
\end{attention}
\end{corrige}

\begin{corrige}[Exercice 7 — Linkage partiel : Cartographie chez la Drosophile]
\textbf{1. Classes parentales et recombinées.}

La femelle est en phase de couplage $e^+\;v^+ // e\;v$ : ses gamètes parentaux sont $e^+\;v^+$ et $e\;v$. Le mâle testeur $e\;v // e\;v$ étant double récessif, le phénotype des descendants révèle le gamète maternel.
\begin{itemize}
    \item \textbf{Classes parentales} (majoritaires) : [gris, normales] = 412 et [ébène, vestigiales] = 398, soit $412 + 398 = 810$ individus.
    \item \textbf{Classes recombinées} (minoritaires) : [gris, vestigiales] = 88 et [ébène, normales] = 92, soit $88 + 92 = 180$ individus.
\end{itemize}
Total : $810 + 180 = 990$ descendants. Les classes ne sont pas équiprobables (810 parentaux contre 180 recombinés) : les gènes sont \textbf{liés} (sur le même chromosome), avec un linkage \emph{partiel} car des recombinants existent (crossing-over en prophase I).

\textbf{2. Taux de recombinaison.}
\[ RF = \frac{\text{recombinés}}{\text{total}} \times 100 = \frac{88 + 92}{990} \times 100 = \frac{180}{990} \times 100 \approx \SI{18,2}{\percent}. \]
Les gènes $e$ et $v$ sont distants de \textbf{18,2 cM} (centiMorgans) : 1 \% de recombinaison correspond par convention à 1 cM.

\textbf{3. Carte factorielle partielle.}

\begin{popfigure}
\begin{tikzpicture}[scale=1, >=Stealth]
\draw[thick, popblue] (0,0) -- (5.46,0);
\draw[thick, popblue] (0,-0.12) -- (0,0.12);
\draw[thick, popblue] (5.46,-0.12) -- (5.46,0.12);
\node[below=4pt] at (0,0) {$e$};
\node[below=4pt] at (5.46,0) {$v$};
\node[above=6pt] at (2.73,0) {\textbf{18,2 cM}};
\end{tikzpicture}
\legende{Carte factorielle des gènes $e$ (corps ébène) et $v$ (ailes vestigiales) chez \emph{Drosophila melanogaster} : distance génétique de 18,2 cM, déduite du taux de recombinaison de \SI{18,2}{\percent}.}
\end{popfigure}

\textbf{Barème indicatif (sur 4 pts) :} identification parentaux/recombinés avec justification par les effectifs (1 pt) ; conclusion « gènes liés » argumentée (0,5 pt) ; calcul du RF avec formule et application numérique (1,5 pt) ; carte factorielle légendée avec distance en cM (1 pt).

\begin{attention}[Point de vigilance]
Le taux de recombinaison se calcule \emph{uniquement} à partir des classes recombinées. Ne jamais inclure les parentaux au numérateur. Vérifier que le total utilisé au dénominateur (990) est bien la somme de toutes les classes.
\end{attention}
\end{corrige}

\begin{corrige}[Exercice 8 — Pedigree : Drépanocytose au Cameroun]
\textbf{1. Mode de transmission.} La drépanocytose est une maladie à transmission \cle{autosomique récessive} : pour le caractère « maladie », l'allèle $S$ (drépanocytaire) est récessif devant l'allèle $A$ (normal). Seuls les homozygotes $S/S$ sont atteints de drépanocytose majeure. (Remarque : au niveau moléculaire, il y a codominance : les hétérozygotes $A/S$ synthétisent à la fois Hb A et Hb S, décelables par électrophorèse.)

\textbf{2. Génotype des parents.} Les deux parents sont sains mais ont une fille atteinte ($S/S$). Cette fille a reçu un allèle $S$ de chaque parent. Chaque parent, phénotypiquement normal, possède donc au moins un allèle $A$. Conclusion :
\[ \text{Père : } A/S \qquad \text{Mère : } A/S \quad \text{(porteurs sains, hétérozygotes)}. \]
Cohérent avec les antécédents : la sœur du père était atteinte ($S/S$), ce qui impose que les grands-parents paternels étaient tous deux porteurs $A/S$ ; le père a alors pu recevoir l'allèle $S$ de l'un d'eux.

\textbf{3. Vérification du pedigree.} Génération I : les deux couples de grands-parents (I-1/I-2 côté maternel, I-3/I-4 côté paternel) sont sains ; du côté paternel, ils sont obligatoirement porteurs $A/S$ puisqu'une de leurs filles (la tante paternelle, génération II) est atteinte $S/S$ (symbole plein). Génération II : II-1 (mère) et II-2 (père), tous deux porteurs sains $A/S$ (symboles à moitié remplis ou hachurés). Génération III : III-1 et III-2 (garçons sains, $A/A$ ou $A/S$), III-3 (fille atteinte $S/S$, symbole plein). Le pedigree correspond bien à l'énoncé : des parents sains ayant un enfant atteint est la signature classique d'une transmission autosomique récessive.

\textbf{4. Risques pour un quatrième enfant.} Croisement $A/S \times A/S$ :

\begin{center}
\begin{tabular}{|c|c|c|}\hline
\rowcolor{popblueL} & $A$ (1/2) & $S$ (1/2) \\ \hline
$A$ (1/2) & $A/A$ sain (1/4) & $A/S$ porteur sain (1/4) \\ \hline
$S$ (1/2) & $A/S$ porteur sain (1/4) & $S/S$ atteint (1/4) \\ \hline
\end{tabular}
\end{center}

\begin{itemize}
    \item Probabilité d'un enfant atteint ($S/S$) : $\dfrac{1}{4} = \SI{25}{\percent}$.
    \item Probabilité d'un enfant porteur sain ($A/S$) : $\dfrac{1}{4} + \dfrac{1}{4} = \dfrac{1}{2} = \SI{50}{\percent}$.
\end{itemize}
Chaque grossesse est un événement indépendant : le fait d'avoir déjà une fille atteinte ne modifie pas ces probabilités.

\textbf{5. Prévention (contexte camerounais).}
\begin{itemize}
    \item \textbf{Conseil génétique} : information du couple sur le risque de 1/4 à chaque grossesse, dans un centre de référence (ex. Hôpital Central de Yaoundé, Hôpital Laquintinie de Douala).
    \item \textbf{Dépistage prénatal} : analyse de l'ADN fœtal après prélèvement de villosités choriales, lorsque cette technique est accessible.
    \item \textbf{Dépistage néonatal systématique} (électrophorèse de l'hémoglobine ou test de solubilité à la naissance) : un diagnostic précoce permet la prise en charge immédiate (pénicilline V en prévention des infections, acide folique, vaccination anti-pneumocoque, hydratation, hydroxyurée), qui améliore considérablement la survie des enfants SS.
\end{itemize}
\end{corrige}

\begin{corrige}[Exercice 9 — Type Bac : Linkage et carte factorielle (3 gènes sur X)]
\textbf{1. Classes parentales, recombinées simples et doubles ; ordre des gènes.}

Les mâles F1 sont hémizygotes : leur phénotype révèle directement le chromosome X reçu de la mère triple hétérozygote.
\begin{itemize}
    \item \textbf{Classes parentales} (les plus fréquentes) : $y^+\;w^+\;m^+$ (580) et $y\;w\;m$ (592), soit 1172 individus.
    \item \textbf{Doubles crossing-over (DCO)} (les plus rares) : $y^+\;w\;m$ (3) et $y\;w^+\;m^+$ (3), soit 6 individus.
    \item \textbf{Simples crossing-over (SCO)} : $y^+\;w\;m^+$ (42) + $y\;w^+\;m$ (38) = 80 ; et $y^+\;w^+\;m$ (124) + $y\;w\;m^+$ (118) = 242.
\end{itemize}
Vérification du total : $1172 + 6 + 80 + 242 = 1500$ individus.

\textbf{Détermination de l'ordre :} on compare une classe DCO à la classe parentale la plus proche. Parental $y^+\;w^+\;m^+$ $\rightarrow$ DCO $y^+\;w\;m^+$ : seul le marqueur $w$ a « changé de chromosome ». Or le gène qui permute seul lors d'un double crossing-over est le gène \emph{médian}. Donc \textbf{$w$ est au milieu} ; l'ordre est :
\[ y \;—\; w \;—\; m \qquad (\text{ou } m \;—\; w \;—\; y). \]

\textbf{2. Taux de recombinaison.} Taux de DCO $= \dfrac{6}{1500} \times 100 = \SI{0,4}{\percent}$.
\begin{align*}
RF_{y-w} &= \frac{80 + 6}{1500} \times 100 = \frac{86}{1500} \times 100 \approx \SI{5,73}{\percent} \;\Rightarrow\; 5,73 \text{ cM} \\
RF_{w-m} &= \frac{242 + 6}{1500} \times 100 = \frac{248}{1500} \times 100 \approx \SI{16,53}{\percent} \;\Rightarrow\; 16,53 \text{ cM} \\
RF_{y-m} &= \frac{80 + 242}{1500} \times 100 = \frac{322}{1500} \times 100 \approx \SI{21,47}{\percent} \;\Rightarrow\; 21,47 \text{ cM}
\end{align*}

\begin{attention}[Point de vigilance majeur]
Pour la distance entre les gènes \emph{extrêmes} ($y$ et $m$), les DCO rétablissent la disposition parentale des marqueurs extrêmes : ils ne sont \textbf{pas} comptés dans $RF_{y-m}$. En revanche, ils sont comptés dans chaque intervalle partiel ($y$-$w$ et $w$-$m$). C'est pourquoi la somme des intervalles ($5,73 + 16,53 = 22,26$ cM) dépasse $RF_{y-m}$ de exactement deux fois le taux de DCO : $21,47 + 2 \times 0,4 = 22,27 \approx 22,26$ cM.
\end{attention}

\textbf{3. Vérification par la fonction de mapping.}
\[ RF_{y-m}^{\text{attendu}} = RF_{y-w} + RF_{w-m} - 2 \times RF_{y-w} \times RF_{w-m} \]
\[ = 5,73 + 16,53 - 2 \times \frac{5,73 \times 16,53}{100} = 22,26 - 1,89 \approx \SI{20,37}{\percent}. \]
Valeur observée : \SI{21,47}{\percent}. L'écart est faible ($\approx$ 1 point) : les résultats sont \textbf{cohérents} avec l'hypothèse d'absence d'interférence. Vérification directe plus simple : $RF_{y-m} = RF_{y-w} + RF_{w-m} - 2 \times \text{DCO} = 22,26 - 0,8 = \SI{21,46}{\percent} \approx \SI{21,47}{\percent}$, exactement la valeur observée (à l'arrondi près).

\textbf{4. Carte factorielle.} Ordre $y - w - m$, avec $y$-$w$ = 5,73 cM et $w$-$m$ = 16,53 cM (distance totale $y$-$m$ = 22,26 cM en unités de carte) : voir la figure de l'énoncé.

\textbf{5. Limite de l'Expérience 2.} L'Expérience 2 ne porte que sur deux gènes ($y$ et $w$) : elle donne la distance $y$-$w$ (\SI{1,5}{\percent} — valeur différente de celle de l'Expérience 1 car obtenue dans un autre croisement ; les taux de recombinaison varient selon les souches et les conditions expérimentales) mais ne fournit \emph{aucune information} sur la position de $m$. Pour ordonner trois gènes, il faut un croisement impliquant \emph{simultanément} les trois marqueurs (test-cross d'un triple hétérozygote) : seules les classes DCO, en révélant le gène médian, permettent d'établir l'ordre.

\textbf{Barème indicatif (sur 10 pts) :} identification des classes (parentales/SCO/DCO) avec justification par les effectifs (2 pts) ; démonstration de l'ordre par comparaison DCO/parentaux (2 pts) ; trois calculs de RF avec formules (3 pts) ; vérification de cohérence (1 pt) ; carte factorielle (1 pt) ; explication de la limite de l'Expérience 2 (1 pt).
\end{corrige}

\begin{corrige}[Exercice 10 — Type Bac : Hémophilie A et conseil génétique]
\textbf{1. Pedigree et génotypes.} Voir la figure de l'énoncé. Génotypes :
\begin{itemize}
    \item I-1 (grand-père hémophile) : $X^h Y$ ; I-2 (grand-mère saine, sans antécédent) : $X^H X^H$.
    \item II-1 (fille) : $X^H X^h$ (porteuse obligée) ; II-2 (fils sain) : $X^H Y$ ; II-3 (époux sain) : $X^H Y$.
    \item III-1 (garçon hémophile) : $X^h Y$ ; III-2 (fille saine) : $X^H X^H$ ou $X^H X^h$.
\end{itemize}

\textbf{2. Justifications.}
\begin{itemize}
    \item \textbf{I-2} est $X^H X^H$ : elle est saine et sans antécédent familial ; l'allèle $h$ étant rare dans la population, on la considère homozygote normale. (Cohérent : son fils II-2 est sain, il a bien reçu $X^H$.)
    \item \textbf{II-1} est \emph{obligatoirement} $X^H X^h$ : son père I-1 ($X^h Y$) transmet son unique chromosome X, porteur de $h$, à \emph{toutes} ses filles. Étant saine, elle possède aussi un $X^H$ (reçu de sa mère). C'est une \cle{porteuse obligée}. Confirmation : elle a un fils hémophile (III-1), qui a reçu son $X^h$.
\end{itemize}

\textbf{3. Probabilité que III-2 soit porteuse.} Croisement II-1 ($X^H X^h$) $\times$ II-3 ($X^H Y$). Une fille reçoit obligatoirement $X^H$ de son père et, de sa mère, $X^H$ ou $X^h$ avec probabilité 1/2 chacun. Toutes les filles de ce croisement sont saines ; sachant que III-2 est une fille :
\[ P(\text{III-2 porteuse}) = \frac{1}{2}. \]

\textbf{4. Risques pour une grossesse de III-2 avec un homme sain ($X^H Y$).}

Raisonnement en deux étapes : $P(\text{III-2 est } X^H X^h) = 1/2$ ; si elle est porteuse, le croisement $X^H X^h \times X^H Y$ donne pour chaque enfant : fille $X^H X^H$ (1/4), fille $X^H X^h$ (1/4), garçon $X^H Y$ (1/4), garçon $X^h Y$ (1/4).
\begin{enumerate}
    \item \textbf{Garçon hémophile} : $\dfrac{1}{2} \times \dfrac{1}{4} = \dfrac{1}{8} = \SI{12,5}{\percent}$.
    \item \textbf{Fille porteuse} : $\dfrac{1}{2} \times \dfrac{1}{4} = \dfrac{1}{8} = \SI{12,5}{\percent}$.
    \item \textbf{Enfant atteint ou porteur} : seuls les garçons $X^h Y$ (atteints, 1/8) et les filles $X^H X^h$ (porteuses, 1/8) comptent : $\dfrac{1}{8} + \dfrac{1}{8} = \dfrac{1}{4} = \SI{25}{\percent}$.
\end{enumerate}

\textbf{5. Message de conseil génétique (exemple).}
« Madame, votre famille est atteinte d'hémophilie A, maladie liée au chromosome X. Vous avez une chance sur deux d'être porteuse ; dans ce cas, chaque garçon à naître aurait une chance sur deux d'être hémophile, et chaque fille une chance sur deux d'être porteuse saine. Un test génétique (recherche de la mutation sur un échantillon de votre sang) peut déterminer si vous êtes porteuse. Si oui, un diagnostic prénatal (prélèvement de villosités choriales vers la 12\textsuperscript{e} semaine, puis analyse de l'ADN fœtal) permettra de savoir si l'enfant est atteint et de préparer sa prise en charge. Nous vous invitons à en discuter ensemble en consultation de génétique. »

\textbf{Barème indicatif (sur 10 pts) :} pedigree correct avec conventions (carré/rond, rempli, numérotation) (2 pts) ; génotypes justifiés, en particulier la porteuse obligée II-1 (2 pts) ; probabilité 1/2 pour III-2 (1 pt) ; calculs des trois risques avec arbre ou tableau (3 pts) ; message clair, scientifiquement exact et adapté (2 pts).

\begin{attention}[Point de vigilance]
Ne pas oublier de conditionner par la probabilité que III-2 soit porteuse (1/2) : un garçon hémophile n'est possible que si la mère transmet $X^h$, ce qui exige qu'elle le possède. Répondre « 1/4 » à la question 4a est l'erreur classique.
\end{attention}
\end{corrige}

\begin{corrige}[Exercice 11 — Type Bac : Mutations, drépanocytose et prévention]
\textbf{1. Analyse moléculaire.}

Le brin donné est le brin \emph{codant} : l'ARNm a la même séquence (avec U au lieu de T). Traduction codon par codon :
\begin{itemize}
    \item Allèle $A$ : CTG–ACT–CCT–\textbf{GAG}–GAG–AAG–TCT $\rightarrow$ Leu–Thr–Pro–\textbf{Glu}–Glu–Lys–Ser
    \item Allèle $S$ : CTG–ACT–CCT–\textbf{GTG}–GAG–AAG–TCT $\rightarrow$ Leu–Thr–Pro–\textbf{Val}–Glu–Lys–Ser
\end{itemize}
Il s'agit d'une \cle{substitution} (remplacement d'une paire de bases A–T par T–A : transversion) au premier nucléotide du 6\textsuperscript{e} codon. Conséquence : mutation \cle{faux-sens} (missense) — le codon reste sens mais code un autre acide aminé : \textbf{Glutamate (Glu) $\rightarrow$ Valine (Val) en position 6} de la chaîne $\beta$-globine (notation E6V). Le remplacement d'un acide aminé hydrophile chargé (Glu) par un acide aminé hydrophobe (Val) provoque la polymérisation de l'hémoglobine S en conditions de faible oxygénation, d'où la falciformation des hématies.

\textbf{2. Analyse cytogénétique.}

Le Document 2 montre une \cle{mutation chromosomique numérique} (aneuploïdie) : présence d'un chromosome 21 surnuméraire ($2n = 47$ au lieu de 46). C'est une \textbf{trisomie 21}, résultant d'une non-disjonction du chromosome 21 en méiose (anaphase I ou II) chez un parent. Le gamète anormal à l'origine d'un enfant trisomique contient \textbf{deux chromosomes 21}, soit 24 chromosomes au total :
\[ \text{Formule du gamète : } 24, X, +21 \quad \text{ou} \quad 24, Y, +21. \]
Après fécondation par un gamète normal (23, X ou 23, Y) : zygote à 47 chromosomes (47, XX, +21 ou 47, XY, +21).

\textbf{3. Interprétation de l'électrophorèse.}

Deux bandes : Hb A \emph{et} Hb S $\Rightarrow$ l'individu synthétise les deux hémoglobines : son génotype est \textbf{$A/S$ (hétérozygote)}. C'est un \textbf{porteur sain} (trait drépanocytaire) : il n'est \emph{pas} malade (l'allèle $A$ assure une production suffisante de Hb A fonctionnelle), mais il peut transmettre l'allèle $S$ à sa descendance. (Rappel : une bande unique en position S correspondrait à un homozygote $S/S$, malade ; une bande unique en position A à un homozygote $A/A$.)

\textbf{4. Épidémiologie et prévention.}

\begin{enumerate}
    \item \textbf{Maintien de l'allèle $S$ — avantage sélectif de l'hétérozygote.} En zone d'endémie palustre (savanes du Nord et de l'Extrême-Nord camerounais), \emph{Plasmodium falciparum} se développe mal dans les hématies des hétérozygotes $A/S$ (falciformation partielle en cas d'infection, élimination plus rapide des hématies parasitées). Les $A/S$ survivent mieux au paludisme que les $A/A$, tandis que les $S/S$ meurent souvent jeunes de la drépanocytose. C'est un \textbf{polymorphisme équilibré} (avantage de l'hétérozygote) : la sélection élimine l'allèle $S$ chez les homozygotes mais le maintient chez les hétérozygotes, d'où sa fréquence élevée (10 à \SI{25}{\percent}) en zone d'endémie.
    \item \textbf{Stratégie nationale de prévention.}
    \begin{itemize}
        \item \textbf{Dépistage néonatal systématique} (électrophorèse ou test de solubilité) dans les maternités, pour diagnostiquer tôt les enfants SS.
        \item \textbf{Conseil génétique et dépistage prénuptial/préconceptionnel} : information des couples dont les deux membres sont $A/S$ (risque de 1/4 d'enfant SS à chaque grossesse).
        \item \textbf{Prise en charge précoce des enfants SS} : pénicilline V quotidienne (prévention des infections pneumococciques), acide folique, vaccination anti-pneumocoque et anti-méningocoque, hydratation, hydroxyurée (réduit les crises vaso-occlusives), transfusions si nécessaire.
        \item \textbf{Rôle des centres de référence} : le Centre de Prise en Charge de la Drépanocytose de l'Hôpital Laquintinie de Douala et le service d'hématologie de l'Hôpital Central de Yaoundé assurent diagnostic, suivi, éducation thérapeutique des familles et formation du personnel de santé.
    \end{itemize}
\end{enumerate}

\textbf{Barème indicatif (sur 10 pts) :} traduction correcte des deux séquences et identification Glu$\rightarrow$Val en position 6 (2,5 pts) ; mutation chromosomique numérique identifiée + formule du gamète 24, X, +21 (2 pts) ; génotype $A/S$ déduit de l'électrophorèse avec justification porteur sain (1,5 pts) ; explication du polymorphisme équilibré (2 pts) ; stratégie de prévention cohérente et contextualisée (2 pts).

\begin{attention}[Points de vigilance]
\begin{itemize}
    \item Le brin donné est le brin \emph{codant} : on traduit directement en remplaçant T par U. Ne pas transcrire en brin complémentaire.
    \item Une substitution n'entraîne \emph{pas} de décalage du cadre de lecture (contrairement à une insertion ou une délétion) : seul un codon est modifié.
    \item Pour l'avantage sélectif, citer explicitement le parasite \emph{Plasmodium falciparum} et le terme « hétérozygote $A/S$ » : l'avantage ne concerne ni les $A/A$ ni les $S/S$.
\end{itemize}
\end{attention}
\end{corrige}

\newpage

\coursec{Fiche bilan}

\courssub{Carte mentale de la leçon}
\begin{popfigure}
\begin{center}\tcbox[colback=popink,colframe=popink,coltext=white,arc=9pt,boxsep=4pt,left=6pt,right=6pt]{\bfseries\small Transmission des caractères}\end{center}
\vspace{-2pt}
\begin{tcbraster}[raster columns=3,raster equal height=rows,raster column skip=6pt,raster row skip=6pt,enhanced,blankest]
\begin{tcolorbox}[enhanced,colback=poporange,colframe=poporange,coltext=white,boxrule=0.9pt,arc=3pt,left=4pt,right=4pt,top=3pt,bottom=3pt,boxsep=1pt,halign=left]\bfseries\scriptsize 1. Allèles \& Monohybridisme (Dominance, Codominance, Test-cross)\end{tcolorbox}
\begin{tcolorbox}[enhanced,colback=popgreen,colframe=popgreen,coltext=white,boxrule=0.9pt,arc=3pt,left=4pt,right=4pt,top=3pt,bottom=3pt,boxsep=1pt,halign=left]\bfseries\scriptsize 2. Dihybridisme \& Brassage inter (Indépendance, 9:3:3:1)\end{tcolorbox}
\begin{tcolorbox}[enhanced,colback=poppurple,colframe=poppurple,coltext=white,boxrule=0.9pt,arc=3pt,left=4pt,right=4pt,top=3pt,bottom=3pt,boxsep=1pt,halign=left]\bfseries\scriptsize 3. Linkage \& Crossing-over (Gènes liés, Carte factorielle)\end{tcolorbox}
\begin{tcolorbox}[enhanced,colback=poppink,colframe=poppink,coltext=white,boxrule=0.9pt,arc=3pt,left=4pt,right=4pt,top=3pt,bottom=3pt,boxsep=1pt,halign=left]\bfseries\scriptsize 4. Pédigrées \& Autosomique (Albinisme, Drépanocytose, Groupes sanguins)\end{tcolorbox}
\begin{tcolorbox}[enhanced,colback=popblue,colframe=popblue,coltext=white,boxrule=0.9pt,arc=3pt,left=4pt,right=4pt,top=3pt,bottom=3pt,boxsep=1pt,halign=left]\bfseries\scriptsize 5. Hérédité Gonosomique (Daltonisme, Hémophilie, Transmission croisée)\end{tcolorbox}
\begin{tcolorbox}[enhanced,colback=popgold,colframe=popgold,coltext=white,boxrule=0.9pt,arc=3pt,left=4pt,right=4pt,top=3pt,bottom=3pt,boxsep=1pt,halign=left]\bfseries\scriptsize 6. Mutations \& Risques (Géniques, Chromosomiques, Prévention, Conseil)\end{tcolorbox}
\end{tcbraster}

\legende{Carte mentale récapitulative de la leçon 04 : les 6 piliers de la génétique formelle et humaine au programme de Terminale C/TI.}
\end{popfigure}

\courssub{Bilan des savoirs et savoir-faire}
\begin{aretenir}[Bilan des savoirs et savoir-faire]

\textbf{1. Définitions clés (à connaître par cœur)}
\begin{itemize}
    \item \cle{Gène} : portion d'ADN codant pour une protéine ou un ARN fonctionnel, déterminant un caractère. Localisé sur un chromosome à un \cle{locus} précis.
    \item \cle{Allèle} : forme alternative d'un gène occupant le même locus sur les chromosomes homologues. Notation : $A$, $a$, $I^{A}$, $I^{B}$, $i$, $X^{h}$, $X^{H}$.
    \item \cle{Génotype} : constitution allélique d'un individu pour un ou plusieurs gènes (ex. $[A//a]$, $[I^{A}//i]$, $[X^{h}//Y]$). \cle{Phénotype} : expression observable du caractère (morphologique, biochimique, moléculaire).
    \item \cle{Homozygote} : deux allèles identiques sur les chromosomes homologues ($[A//A]$ ou $[a//a]$). \cle{Hétérozygote} : deux allèles différents ($[A//a]$).
    \item \cle{Dominance} : allèle s'exprimant phénotypiquement chez l'hétérozygote (ex. $A$ sur $a$). \cle{Récessivité} : allèle masqué chez l'hétérozygote. \cle{Codominance} : les deux allèles s'expriment distinctement chez l'hétérozygote (ex. groupes sanguins $I^{A}$ et $I^{B}$).
    \item \cle{Brassage interchromosomique} : répartition aléatoire et indépendante des paires de chromosomes homologues lors de l'anaphase I de la méiose (loi de Mendel 2) $\rightarrow$ concerne des gènes sur chromosomes \emph{différents} (ou très éloignés sur le même).
    \item \cle{Brassage intrachromosomique (Crossing-over)} : échange réciproque de segments entre chromatides non-sœurs de chromosomes homologues lors de la prophase I (stade pachytène) $\rightarrow$ concerne des gènes \cle{liés} sur le \emph{même} chromosome.
    \item \cle{Linkage total} : gènes très proches, pas de crossing-over observable (fréquence recombinants $\approx 0\%$). \cle{Linkage partiel} : crossing-over possible, fréquence recombinants $> 0\%$ et $< 50\%$.
    \item \cle{Pédigrée (arbre généalogique)} : représentation codifiée de la transmission d'un caractère dans une famille (carré $\sigma$, cercle $\varphi$, symbole plein = atteint, double trait horizontal = consanguinité).
    \item \cle{Hérédité autosomique} : gène porté par un autosome (chromosome non sexuel). Transmission identique chez $\sigma$ et $\varphi$.
    \item \cle{Hérédité gonosomique (liée au sexe)} : gène porté par le chromosome X (hémizygie chez l'homme $XY$). Règle absolue : \textbf{pas de transmission père $\rightarrow$ fils}.
    \item \cle{Mutation génique} : modification de la séquence nucléotidique d'un gène (substitution, insertion, délétion d'un ou quelques nucléotides).
    \item \cle{Mutation chromosomique} : modification de la structure (délétion, duplication, inversion, translocation) ou du nombre de chromosomes (anéuploïdie : trisomie 21 ; polyploïdie).
\end{itemize}

\textbf{2. Lois, formules et règles numériques (Tableau de révision)}
\begin{center}
\begin{tabularx}{\linewidth}{|l|X|X|}\hline
\rowcolor{popblueL}
\textbf{Situation / Croisement} & \textbf{Résultat phénotypique F2 (ou descendance test-cross)} & \textbf{Formule / Règle clé} \\ \hline
Monohybridisme $A/a \times A/a$ (Dominance) & $3\ [A\text{-phéno}] : 1\ [a\text{-phéno}]$ & Loi de ségrégation : $p(A)=p(a)=1/2$ \\ \hline
Monohybridisme $A/a \times A/a$ (Codominance) & $1\ [A/A] : 2\ [A/a] : 1\ [a/a]$ & Phénotype = Génotype (pas de dominance) \\ \hline
Test-cross Monohybride $A/a \times a/a$ & $1\ [A\text{-phéno}] : 1\ [a\text{-phéno}]$ & Vérifie l'hétérozygotie de l'individu dominant \\ \hline
Dihybridisme $A/a; B/b \times A/a; B/b$ (Indépendants) & $9:3:3:1$ (phénotypique) & Loi de l'indépendance : $p(AB)=p(Ab)=p(aB)=p(ab)=1/4$ \\ \hline
Test-cross Dihybride $A/a; B/b \times a/a; b/b$ (Indépendants) & $1:1:1:1$ & 4 gamètes équiprobables \\ \hline
Linkage partiel (Phase Cis $AB/ab$) & Majorité parentaux ($AB$, $ab$) + Minorité recombinants ($Ab$, $aB$) & Fréquence recombinaison $d(\%) = \frac{N_{rec}}{N_{tot}} \times 100$ \\ \hline
Carte factorielle & Distance en \cle{Morganides (cM)} ou \cle{Morgans (M)} & $1\%$ recombinaison $= 1$ cM $= 0,01$ M. Additivité si $d < 10-15\%$ \\ \hline
Risque génétique (Auto récessif, parents sains $A/a$) & Enfant atteint $[a/a]$ & $P = 1/4$ (par grossesse, non conditionnée) \\ \hline
Risque génétique (Gonoso récessif X, mère porteuse $X^{H}X^{h}$) & Fils atteint $X^{h}Y$ & $P = 1/2$ (parmi les fils) ; $P = 1/4$ (parmi tous les enfants) \\ \hline
\end{tabularx}
\end{center}

\end{aretenir}

\begin{methode}[Croisement monohybride / dihybride (brassage interchromosomique)]
\begin{itemize}
    \item Identifier les allèles et la relation de dominance/codominance (lire l'énoncé).
    \item Écrire le génotype des parents (génération P) avec la notation $[Allèle_1//Allèle_2]$ ; séparer les paires de gènes par un point-virgule si dihybridisme.
    \item Déterminer les gamètes (méiose : ségrégation allélique + brassage interchromosomique si gènes sur chromosomes différents). Vérifier l'équiprobabilité ($1/2$ par allèle, $1/4$ par combinaison dihybride).
    \item Croiser les gamètes (tableau de Punnett ou produit cartésien).
    \item Lire les génotypes et phénotypes de la descendance (F1 ou F2). Exprimer les ratios en nombres entiers simplifiés.
\end{itemize}
\end{methode}

\begin{methode}[Test-cross (croisement test)]
\begin{itemize}
    \item Croiser l'individu à phénotype dominant (génotype inconnu $[A\text{-}]$ ou $[A\text{-};B\text{-}]$) avec un individu \cle{homozigote récessif} ($[a//a]$ ou $[a//a;b//b]$).
    \item Les phénotypes de la descendance révèlent \emph{directement} les gamètes produits par l'individu testé (l'apport du récessif est constant).
    \item Ratio $1:1$ (monohybride) ou $1:1:1:1$ (dihybride indépendant) $\rightarrow$ individu testé hétérozygote. Ratio phénotype dominant seul $\rightarrow$ individu testé homozygote dominant.
\end{itemize}
\end{methode}

\begin{methode}[Construction d'une carte factorielle (3 gènes liés)]
\begin{itemize}
    \item Analyser une descendance de test-cross trihybride (ou données fournies). Compter les effectifs des 8 phénotypes : 2 parentaux (majoritaires), 4 recombinants simples, 2 recombinants doubles (très rares).
    \item Calculer les fréquences de recombinaison 2 à 2 : $d_{AB}$, $d_{BC}$, $d_{AC}$ (en \%, arrondi au dixième).
    \item Ordonner les gènes : la plus grande distance = gènes extrêmes. Le gène du milieu est celui dont la somme des distances aux deux autres $\approx$ distance max.
    \item Vérifier la cohérence : $d_{\text{extrêmes}} \approx d_1 + d_2$ (si pas d'interférence forte / distances $< 10-15\%$).
    \item Dessiner la carte : \textendash[ $d_1$ cM ]\textendash Gène Milieu \textendash[ $d_2$ cM ]\textendash.
\end{itemize}
\end{methode}

\begin{methode}[Analyse d'un pedigree (Autosomique vs Gonosomique)]
\begin{itemize}
    \item \textbf{Autosomique dominant} : caractère présent à chaque génération (si pénétrance complète) ; transmission $\sigma \rightarrow \sigma$, $\varphi \rightarrow \varphi$, $\sigma \rightarrow \varphi$, $\varphi \rightarrow \sigma$ ; atteints = hétérozygotes le plus souvent ; parents atteints $\times$ sains $\rightarrow$ $1/2$ atteints.
    \item \textbf{Autosomique récessif} : saute des générations ; consanguinité favorise l'apparition ; parents sains (hétérozygotes) peuvent avoir enfants atteints ($A/a \times A/a \rightarrow 1/4$) ; $\sigma$ et $\varphi$ atteints en proportion égale.
    \item \textbf{Gonosomique récessif (lié au X)} : $\sigma$ beaucoup plus atteints que $\varphi$ (hémizygie) ; \textbf{pas de transmission père $\rightarrow$ fils} ; père atteint $\rightarrow$ toutes les filles porteuses (obligées) ; mère porteuse $\rightarrow$ $1/2$ fils atteints, $1/2$ filles porteuses.
    \item \textbf{Gonosomique dominant (lié au X)} : père atteint $\rightarrow$ toutes les filles atteintes, aucun fils atteint ; mère atteinte (hétérozygote) $\rightarrow$ $1/2$ enfants atteints (les deux sexes).
    \item \textbf{Y-lié (holandrique)} : seulement les $\sigma$ ; transmission père $\rightarrow$ \emph{tous} les fils.
\end{itemize}
\end{methode}

\begin{methode}[Calcul de risque génétique (Prévision génétique)]
\begin{itemize}
    \item Déduire les génotypes probables des consultants (probabilités conditionnées par leur phénotype, leur fratrie, le pedigree).
    \item Déterminer les gamètes possibles pour chaque parent et leurs probabilités.
    \item Croiser les gamètes (tableau de Punnett pondéré) pour obtenir les probabilités des génotypes de la descendance.
    \item Calculer la probabilité du génotype « à risque » (ex. $[a//a]$ pour maladie récessive, $X^{h}Y$ pour hémophilie).
    \item \textbf{Attention} : distinguer « risque par enfant (par grossesse) » et « risque pour la descendance globale » ; ne pas confondre probabilité \emph{a priori} et probabilité \emph{conditionnelle} (ex. : « sachant qu'ils ont déjà un enfant sain »).
\end{itemize}
\end{methode}

\begin{aretenir}[Checklist avant le Bac]
    $\square$ Je sais définir : gène, allèle, locus, génotype, phénotype, homozygote, hétérozygote, dominance, codominance, récessivité.\par
    $\square$ Je sais écrire et interpréter les croisements monohybrides (F1, F2, test-cross) avec dominance et codominance (groupes sanguins ABO : allèles $I^{A}$, $I^{B}$, $i$).\par
    $\square$ Je maîtrise le dihybridisme : loi de l'indépendance (gènes sur chromosomes différents), ratios $9:3:3:1$ (F2) et $1:1:1:1$ (test-cross), détermination des gamètes parentaux.\par
    $\square$ Je distingue brassage interchromosomique (gènes sur chromosomes différents, indépendance) et intrachromosomique (gènes liés sur même chromosome, crossing-over).\par
    $\square$ Je sais calculer une fréquence de recombinaison $d(\%) = \frac{N_{rec}}{N_{tot}} \times 100$ et construire une carte factorielle à 3 gènes (ordonnancement, distances en cM, additivité pour $d < 15\%$).\par
    $\square$ Je sais lire, annoter et interpréter un pedigree : reconnaître autosomique dominant/récessif et gonosomique récessif/dominant (règles clés : pas père$\rightarrow$fils, hémizygie homme, consanguinité).\par
    $\square$ Je connais les exemples types du programme : albinisme (auto récessif), drépanocytose (auto codominant : $[Hb^A//Hb^S]$ phénotype intermédiaire), achondroplasie (auto dominant), groupes ABO/Rhésus (systèmes multi-alléliques), daltonisme/hémophilie (gonoso récessif X).\par
    $\square$ Je sais calculer un risque génétique (probabilité conditionnelle) pour un couple consultant (ex: risque drépanocytose si parents $[Hb^A//Hb^S]$, risque hémophilie si mère porteuse).\par
    $\square$ Je différencie mutation génique (échelle nucléotidique : substitution, frameshift) et mutation chromosomique (structure : délétion, inversion, translocation ; nombre : anéuploïdie trisomie 21, polyploïdie).\par
    $\square$ Je cite les méthodes de prévention : conseil génétique, diagnostic prénatal (amniocentèse $\approx$ 16-18 SA, trophobiopsie $\approx$ 11-13 SA), DPI (diagnostic préimplantatoire sur embryon FIV), dépistage néonatal systématique (drépanocytose au Cameroun : test de Guthrie).\par
    $\square$ Je rédige rigoureusement : allèles en italique ($A$, $a$, $X^{h}$), génotypes entre crochets $[A//a]$, phénotypes en romain, unités (cM, \%), flèches $\rightarrow$ pour gamètes/fécondation, ratios en nombres entiers.
\end{aretenir}

\findecours

\end{document}
